% Modéliser un phénomène : exemples de modèles utilisés en physique et en chimie — Physique 1ère TI — Cours signé OnBuch+
% Programme officiel MINESEC. Compiler avec : tectonic N02-modeles-physique-chimie.tex
\newcommand{\DOCMATIERE}{Physique}
\newcommand{\DOCNIVEAU}{1ère TI}
\newcommand{\DOCMODULE}{Module 1 : Mesures et incertitudes}
\newcommand{\DOCLECON}{N02}
\newcommand{\DOCTITRE}{Modéliser un phénomène : exemples de modèles utilisés en physique et en chimie}
% OnBuch+ — préambule « cours signés » ENRICHI (compilé avec tectonic / XeLaTeX)
% Système visuel « Pop » d'OnBuch+ : fond crème (jamais blanc), bordures encre
% épaisses, ombres « dures » décalées, titres Archivo Black en capitales.
% Version MATHÉMATIQUES : graphes (pgfplots/tikz), boîtes pédagogiques
% (définition, propriété, méthode, attention, le savais-tu, exercice résolu,
% exercice, TP, bilan), page de garde et signature OnBuch+ ancrée en bas.
%
% Macros attendues dans main.tex AVANT \input{preamble} :
%   \DOCMATIERE  \DOCNIVEAU  \DOCMODULE  \DOCLECON (numéro)  \DOCTITRE
\documentclass[11pt]{article}

\usepackage{fontspec}
\usepackage[french]{babel}
\usepackage[a4paper,margin=2cm,top=3.4cm,bottom=2.8cm,headheight=1.4cm,footskip=1.3cm]{geometry}
\usepackage{amsmath,amssymb,amsfonts}
\usepackage{mathtools} % \xmapsto, \coloneqq... souvent utilisés par les modèles
\usepackage{cancel} % simplifications barrées (bilans de forces)
% \abs{z} (module, valeur absolue) et \norm{u} (norme), utilisés par les modèles
\providecommand{\abs}{}\let\abs\relax\DeclarePairedDelimiter{\abs}{\lvert}{\rvert}
\providecommand{\norm}{}\let\norm\relax\DeclarePairedDelimiter{\norm}{\lVert}{\rVert}
\usepackage[table]{xcolor} % [table] : fournit aussi \rowcolors (alternance de lignes)
\usepackage{pagecolor}
\usepackage{tikz}
\usetikzlibrary{arrows.meta,positioning,calc,shapes.geometric,shapes.misc,shapes.arrows,decorations.pathmorphing,decorations.pathreplacing,decorations.markings,patterns,shadows,mindmap,trees,fit,backgrounds,matrix,intersections,angles,quotes}
\usepackage{pgfplots}
\usepackage[version=4]{mhchem} % équations nucléaires \ce{^{235}_{92}U}
\usepackage[european,siunitx]{circuitikz} % schémas électriques
\pgfplotsset{compat=1.18}
\usepgfplotslibrary{fillbetween,groupplots} % aires entre courbes (calcul intégral)
\usepackage{enumitem}
\usepackage{fancyhdr}
% [most] charge toutes les bibliothèques tcolorbox (listings, minted, raster,
% documentation...) et gonfle inutilement la mémoire TeX sur un document long
% (~300 boîtes) : on ne charge que ce qui est réellement utilisé.
\usepackage[skins,breakable]{tcolorbox}
\usepackage{graphicx}
\usepackage{colortbl}
\usepackage{booktabs}
\usepackage{tabularx}
\usepackage{diagbox} % case d'en-tête diagonale des tableaux à double entrée
% Colonnes extensibles centrée / gauche / droite (C, L, R), souvent utilisées par les modèles
\newcolumntype{C}{>{\centering\arraybackslash}X}
\newcolumntype{L}{>{\raggedright\arraybackslash}X}
\newcolumntype{R}{>{\raggedleft\arraybackslash}X}
\usepackage{array}
\usepackage{multicol}
\usepackage{siunitx}
% parse-numbers=false : accepte aussi \SI{2,5 \times 10^{-3}}{...}
\sisetup{locale=FR,output-decimal-marker={,},group-minimum-digits=4,parse-numbers=false}
\DeclareSIUnit\ohm{\ensuremath{\symup{\Omega}}} % U+2126 absent de la police
\DeclareSIUnit\division{div} % réglages d'oscilloscope (ms/div, V/div)
\DeclareSIUnit\tr{tr} % tours (vitesse de rotation, ex. \SI{1500}{\tr\per\minute})
\DeclareSIUnit\ch{ch} % cheval-vapeur (puissance moteur, unité usuelle hors SI)
\DeclareSIUnit\franccfa{FCFA} % franc CFA (coûts, contexte camerounais)
\DeclareSIUnit\dioptre{\ensuremath{\delta}} % vergence d'une lentille (dioptrie)
\usepackage{qrcode}
% \degree : commande fréquemment utilisée par les modèles pour un angle ou une
% température en dehors de \SI{}{} (non fournie par les paquets chargés).
\providecommand{\degree}{^{\circ}}
% \qed (fin de démonstration), utilisé par les modèles sans amsthm
\providecommand{\qed}{\hfill$\square$}
% \barre : double barre des valeurs interdites dans un tableau de variations (tkz-tab)
\providecommand{\barre}{\ensuremath{\|}}
% environnement proof (amsthm non chargé) utilisé par les modèles
\ifdefined\proof\else\newenvironment{proof}[1][Démonstration]{\par\noindent\textit{#1.}\ }{\qed\par}\fi
\usepackage{lastpage}
\usepackage{hyperref}
\hypersetup{hidelinks}

% ============================================================
% Palette « Pop » OnBuch+ — mêmes hex que la classe P (plus_ui.dart)
% ============================================================
\definecolor{poporange}{HTML}{F59321}
\definecolor{popdark}{HTML}{E07A0C}
\definecolor{popcream}{HTML}{FFF6E8}
\definecolor{popink}{HTML}{1C1714}
\definecolor{poppurple}{HTML}{7A5AE0}
\definecolor{popgreen}{HTML}{2E9E6A}
\definecolor{poppink}{HTML}{E0457B}
\definecolor{popblue}{HTML}{2F7DE1}
\definecolor{popgold}{HTML}{C98A12}
\definecolor{popmuted}{HTML}{8A7F76}
\definecolor{popline}{HTML}{EADFD2}
\definecolor{popgray}{HTML}{8A7F76}
\colorlet{popgrey}{popgray}
% Alias tolérants (le modèle nomme parfois les couleurs différemment)
\colorlet{popwhite}{white}
\colorlet{popblack}{popink}
\colorlet{popred}{poppink}
\colorlet{popyellow}{popgold}
% Teintes claires pour fonds de tableaux / zones
\colorlet{popblueL}{popblue!10!white}
\colorlet{poporangeL}{poporange!14!white}
\colorlet{popgreenL}{popgreen!12!white}
\colorlet{poppurpleL}{poppurple!12!white}
\colorlet{poppinkL}{poppink!12!white}
\colorlet{popgoldL}{popgold!14!white}

\pagecolor{popcream}

% ============================================================
% Typographie
% ============================================================
\defaultfontfeatures{Ligatures=TeX}
\setmainfont{PlusJakartaSans-Medium.ttf}[Path=../../../fonts/,BoldFont=PlusJakartaSans-Bold.ttf]
% Maths en sans-serif (Fira Math) avec chiffres et lettres droites de Plus
% Jakarta, pour que les formules se fondent dans le texte.
\usepackage[math-style=ISO,bold-style=ISO]{unicode-math}
\setmathfont{FiraMath-Regular.otf}[Path=../../../fonts/]
\setmathfont{PlusJakartaSans-Medium.ttf}[Path=../../../fonts/,range={up/num,up/{Latin,latin},\mathup}]
\usepackage{icomma} % virgule décimale sans espace en mode math
% Caractères mathématiques Unicode absents des polices de texte (Plus Jakarta,
% Archivo Black) : rendus par leur équivalent mathématique, en texte comme en
% formule — sinon ils s'affichent en carré vide (ex. « Arithmétique dans ℤ »).
\usepackage{newunicodechar}
\newunicodechar{ℝ}{\ensuremath{\mathbb{R}}}
\newunicodechar{ℤ}{\ensuremath{\mathbb{Z}}}
\newunicodechar{ℂ}{\ensuremath{\mathbb{C}}}
\newunicodechar{ℕ}{\ensuremath{\mathbb{N}}}
\newunicodechar{ℚ}{\ensuremath{\mathbb{Q}}}
\newunicodechar{𝔻}{\ensuremath{\mathbb{D}}}
\newunicodechar{α}{\ensuremath{\alpha}}
\newunicodechar{β}{\ensuremath{\beta}}
\newunicodechar{γ}{\ensuremath{\gamma}}
\newunicodechar{δ}{\ensuremath{\delta}}
\newunicodechar{ε}{\ensuremath{\varepsilon}}
\newunicodechar{θ}{\ensuremath{\theta}}
\newunicodechar{λ}{\ensuremath{\lambda}}
\newunicodechar{μ}{\ensuremath{\mu}}
\newunicodechar{σ}{\ensuremath{\sigma}}
\newunicodechar{τ}{\ensuremath{\tau}}
\newunicodechar{φ}{\ensuremath{\varphi}}
\newunicodechar{ψ}{\ensuremath{\psi}}
\newunicodechar{ω}{\ensuremath{\omega}}
\newunicodechar{Σ}{\ensuremath{\Sigma}}
% majuscules produites par \MakeUppercase dans les titres (α-aminés -> Α)
\newunicodechar{Α}{\ensuremath{\alpha}}
\newunicodechar{Β}{\ensuremath{\beta}}
\newunicodechar{Γ}{\ensuremath{\Gamma}}
\newunicodechar{Δ}{\ensuremath{\Delta}}
\newunicodechar{Ε}{\ensuremath{\varepsilon}}
\newunicodechar{Θ}{\ensuremath{\Theta}}
\newunicodechar{Λ}{\ensuremath{\Lambda}}
\newunicodechar{Μ}{\ensuremath{\mu}}
\newunicodechar{Τ}{\ensuremath{\tau}}
\newunicodechar{Φ}{\ensuremath{\Phi}}
\newunicodechar{Ψ}{\ensuremath{\Psi}}
\newunicodechar{Ω}{\ensuremath{\Omega}}
\newunicodechar{↦}{\ensuremath{\mapsto}}
\newunicodechar{∈}{\ensuremath{\in}}
\newunicodechar{∉}{\ensuremath{\notin}}
\newunicodechar{∧}{\ensuremath{\wedge}}
\newunicodechar{⇔}{\ensuremath{\Leftrightarrow}}
\newunicodechar{⟺}{\ensuremath{\Longleftrightarrow}}
\newunicodechar{⇒}{\ensuremath{\Rightarrow}}
\newunicodechar{⟹}{\ensuremath{\Longrightarrow}}
\newunicodechar{∘}{\ensuremath{\circ}}
\newunicodechar{∪}{\ensuremath{\cup}}
\newunicodechar{∩}{\ensuremath{\cap}}
\newunicodechar{≡}{\ensuremath{\equiv}}
\newunicodechar{⊂}{\ensuremath{\subset}}
\newunicodechar{⊥}{\ensuremath{\perp}}
\newunicodechar{∃}{\ensuremath{\exists}}
\newunicodechar{∀}{\ensuremath{\forall}}
\newunicodechar{∛}{\ensuremath{\sqrt[3]{\,}}}
\newunicodechar{∜}{\ensuremath{\sqrt[4]{\,}}}
\newunicodechar{⃗}{\ensuremath{{}^{\to}}}
\newunicodechar{ⁿ}{\ifmmode^{n}\else\textsuperscript{\ensuremath{n}}\fi}
\newunicodechar{ˣ}{\ifmmode^{x}\else\textsuperscript{\ensuremath{x}}\fi}
\newunicodechar{ᵇ}{\ifmmode^{b}\else\textsuperscript{\ensuremath{b}}\fi}
\newunicodechar{ʳ}{\ifmmode^{r}\else\textsuperscript{\ensuremath{r}}\fi}
\newunicodechar{ᵘ}{\ifmmode^{u}\else\textsuperscript{\ensuremath{u}}\fi}
\newunicodechar{ʸ}{\ifmmode^{y}\else\textsuperscript{\ensuremath{y}}\fi}
\newunicodechar{ᵃ}{\ifmmode^{a}\else\textsuperscript{\ensuremath{a}}\fi}
\newunicodechar{ᵉ}{\ifmmode^{e}\else\textsuperscript{\ensuremath{e}}\fi}
\newunicodechar{ᵏ}{\ifmmode^{k}\else\textsuperscript{\ensuremath{k}}\fi}
\newunicodechar{ᵖ}{\ifmmode^{p}\else\textsuperscript{\ensuremath{p}}\fi}
\newunicodechar{ᴺ}{\ifmmode^{N}\else\textsuperscript{\ensuremath{N}}\fi}
\newunicodechar{ᶜ}{\ifmmode^{c}\else\textsuperscript{\ensuremath{c}}\fi}
\newunicodechar{ʰ}{\ifmmode^{h}\else\textsuperscript{\ensuremath{h}}\fi}
\newunicodechar{ᵠ}{\ifmmode^{q}\else\textsuperscript{\ensuremath{q}}\fi}
\newunicodechar{ₙ}{\ifmmode_{n}\else\textsubscript{\ensuremath{n}}\fi}
\newunicodechar{ₐ}{\ifmmode_{a}\else\textsubscript{\ensuremath{a}}\fi}
\newunicodechar{ᵢ}{\ifmmode_{i}\else\textsubscript{\ensuremath{i}}\fi}
\newunicodechar{ᵣ}{\ifmmode_{r}\else\textsubscript{\ensuremath{r}}\fi}
\newunicodechar{ₖ}{\ifmmode_{k}\else\textsubscript{\ensuremath{k}}\fi}
\newunicodechar{ᵦ}{\ifmmode_{\beta}\else\textsubscript{\ensuremath{\beta}}\fi}
\newunicodechar{ₓ}{\ifmmode_{x}\else\textsubscript{\ensuremath{x}}\fi}
\newfontfamily\popdisplay{ArchivoBlack-Regular.ttf}[Path=../../../fonts/]
\newfontfamily\poplogo{SpaceGrotesk-Bold.ttf}[Path=../../../fonts/]
\newfontfamily\popsemi{PlusJakartaSans-SemiBold.ttf}[Path=../../../fonts/]
\newfontfamily\popxbold{PlusJakartaSans-ExtraBold.ttf}[Path=../../../fonts/]
\newfontfamily\popxboldls{PlusJakartaSans-ExtraBold.ttf}[Path=../../../fonts/,LetterSpace=3.0]

\setlist[enumerate]{leftmargin=1.7em,itemsep=5pt,topsep=4pt}
\setlist[itemize]{leftmargin=1.7em,itemsep=4pt,topsep=4pt,label={\color{poporange}\textbullet}}
\setlength{\parindent}{0pt}
\setlength{\parskip}{5pt}
\renewcommand{\arraystretch}{1.25}

% pgfplots : style Pop par défaut
\pgfplotsset{
  every axis/.append style={axis line style={popink,line width=1pt},tick style={popink},
    grid style={popline,line width=0.5pt},label style={font=\small},tick label style={font=\footnotesize}},
  popcurve/.style={poporange,line width=1.8pt,no markers,smooth},
}
% Même style disponible dans un \draw TikZ simple (hors pgfplots)
\tikzset{popcurve/.style={poporange,line width=1.8pt}}
\tikzset{popgrid/.style={popline,very thin}}
\colorlet{popcurve}{poporange} % le nom du style est parfois utilisé comme couleur

% ============================================================
% Pill / squiggle
% ============================================================
\newcommand{\poppill}[3]{%
  \tikz[baseline=-0.55ex]\node[draw=popink,line width=1.2pt,rounded corners=2.6pt,%
    fill=#2,inner xsep=6.5pt,inner ysep=3pt]{%
    \popxboldls\fontsize{7}{7}\selectfont\color{#3}\MakeUppercase{#1}};%
}
\newcommand{\popsquiggle}[1]{%
  \tikz[baseline=-0.4ex]{
    \draw[#1,line width=2.6pt,line cap=round]
      plot [smooth] coordinates {(0,0.09) (0.34,0.01) (0.68,0.21) (1.02,0.01) (1.36,0.10)};
  }%
}
% Mot-clé surligné (marqueur orange sous le texte)
\newcommand{\cle}[1]{{\popxbold\color{popdark}#1}}

% ============================================================
% En-tête / pied
% ============================================================
\pagestyle{fancy}
\fancyhf{}
\renewcommand{\headrulewidth}{0pt}
\renewcommand{\footrulewidth}{0pt}
\lhead{\raisebox{-0.15ex}{\poplogo\fontsize{13}{13}\selectfont\color{popink}OnBuch\color{poporange}+}}
\rhead{\poppill{\DOCMATIERE\ \textbullet\ \DOCNIVEAU\ \textbullet\ Leçon \DOCLECON}{white}{popink}}
\lfoot{\popxbold\fontsize{7.6}{7.6}\selectfont\color{popmuted}\MakeUppercase{Cours signé OnBuch+}\ \textbullet\ {\popsemi\DOCTITRE}}
\rfoot{\tikz[baseline=-0.6ex]\node[rounded corners=8.5pt,fill=popink,minimum height=17pt,minimum width=17pt,inner xsep=5pt,inner sep=0]{\popxbold\fontsize{8}{8}\selectfont\color{popcream}\thepage\,/\,\pageref*{LastPage}};}

% ============================================================
% Titres
% ============================================================
\newcommand{\chaptitre}[2]{%
  \begin{tcolorbox}[enhanced,colback=poporange,colframe=popink,boxrule=2.6pt,arc=11pt,
      drop shadow=popink,left=15pt,right=15pt,top=12pt,bottom=11pt,before skip=3pt,after skip=18pt]
    \poppill{#2}{popink}{popcream}\par\vspace{9pt}
    {\raggedright\hyphenpenalty=10000\exhyphenpenalty=10000\popdisplay\fontsize{22}{24}\selectfont\color{popcream}\MakeUppercase{#1}\par}
  \end{tcolorbox}
}
% Section numérotée : pastille orange avec le numéro + titre Archivo
\newcounter{popsec}
\newcommand{\coursec}[1]{%
  \stepcounter{popsec}\setcounter{popsub}{0}%
  \par\vspace{16pt}\noindent
  \begin{minipage}{\linewidth}
  \tikz[baseline=-0.7ex]\node[draw=popink,line width=1.6pt,fill=poporange,rounded corners=4pt,minimum size=22pt,inner sep=0]{\popdisplay\fontsize{12}{12}\selectfont\color{popcream}\thepopsec};%
  \hspace{8pt}{\popdisplay\fontsize{15}{17}\selectfont\color{popink}\MakeUppercase{#1}}\par
  \vspace{4pt}\noindent{\color{poporange}\rule{36pt}{3.6pt}}
  \end{minipage}\par\nopagebreak\vspace{8pt}
}
\newcounter{popsub}
\newcommand{\courssub}[1]{%
  \stepcounter{popsub}%
  \par\vspace{9pt}\noindent{\popxbold\fontsize{11.5}{13}\selectfont\color{popink}{\color{poporange}\thepopsec.\thepopsub}\hspace{6pt}#1}\par\nopagebreak\vspace{4pt}
}

% ============================================================
% Boîtes pédagogiques — même recette : fond blanc, bordure épaisse colorée,
% ombre dure encre, badge plein. Titre optionnel [#1].
% ============================================================
\tcbset{popbase/.style={breakable,enhanced,colback=white,boxrule=2.3pt,arc=9pt,
  shadow={3pt}{-3pt}{0pt}{popink},left=14pt,right=14pt,top=11pt,bottom=11pt,before skip=11pt,after skip=15pt,
  fontupper=\popsemi\fontsize{10.6}{14.5}\selectfont\color{popink},
  fontlower=\popsemi\fontsize{10.6}{14.5}\selectfont\color{popink}}}
% Coche dessinée (le glyphe ✓ n'existe pas dans les polices du document)
\newcommand{\popcheck}[1]{\tikz[baseline=-0.5ex]\draw[#1,line width=1.6pt,line cap=round,line join=round](-0.13,0.0)--(-0.04,-0.09)--(0.13,0.1);}
\providecommand{\coche}{\popcheck{popgreen}}
\providecommand{\resultat}[1]{\par\smallskip\noindent\fcolorbox{popgreen}{popgreenL}{\parbox{\dimexpr\linewidth-2\fboxsep-2\fboxrule\relax}{\textbf{Résultat.} #1}}\par\smallskip}
\newcommand{\popboxhead}[3]{\poppill{#1}{#2}{white}\ifx\relax#3\relax\else\hspace{6pt}{\popxbold\fontsize{10.6}{12}\selectfont\color{popink}#3}\fi\par\vspace{7pt}}

\newtcolorbox{aretenir}[1][]{popbase,colframe=popgreen,before upper={\popboxhead{À retenir}{popgreen}{#1}}}
\newtcolorbox{exemplebox}[1][]{popbase,colframe=poporange,before upper={\popboxhead{Exemple}{poporange}{#1}}}
\newtcolorbox{definition}[1][]{popbase,colframe=popblue,before upper={\popboxhead{Définition}{popblue}{#1}}}
\newtcolorbox{propriete}[1][]{popbase,colframe=poppurple,before upper={\popboxhead{Propriété}{poppurple}{#1}}}
\newtcolorbox{methode}[1][]{popbase,colframe=popgold,before upper={\popboxhead{Méthode}{popgold}{#1}}}
\newtcolorbox{attention}[1][]{popbase,colframe=poppink,before upper={\popboxhead{Attention}{poppink}{#1}}}
\newtcolorbox{experience}[1][]{popbase,colframe=popblue,colback=popblueL,before upper={\popboxhead{Expérience}{popblue}{#1}}}
% « Le savais-tu ? » : carte violette pleine (contexte, culture, Cameroun)
\newtcolorbox{savaistu}[1][]{popbase,colback=poppurple,colframe=popink,
  fontupper=\popsemi\fontsize{10.4}{14.2}\selectfont\color{white},
  before upper={\poppill{Le savais-tu\,?}{white}{poppurple}\ifx\relax#1\relax\else\hspace{6pt}{\popxbold\color{white}#1}\fi\par\vspace{7pt}}}
% Exercice résolu : énoncé en haut, \tcblower puis la solution sur fond crème
\newcounter{popexo}\newcounter{popexr}
\newtcolorbox{exoresolu}[1][]{popbase,colframe=poporange,colbacklower=popcream,
  segmentation style={popink,line width=1.2pt,dashed},
  before upper={\stepcounter{popexr}\popboxhead{Exercice résolu \thepopexr}{poporange}{#1}},
  before lower={\poppill{Solution}{popink}{popcream}\par\vspace{6pt}}}
% Exercice (non résolu en place) — la correction est dans la partie Corrigés
\newtcolorbox{exercice}[1][]{popbase,colframe=popink,
  before upper={\stepcounter{popexo}\popboxhead{Exercice \thepopexo}{popink}{#1}}}
% Corrigé
\newtcolorbox{corrige}[1][]{popbase,colframe=popgreen,colback=popgreenL,
  before upper={\popboxhead{Corrigé}{popgreen}{#1}}}
% Objectifs / prérequis / bilan
\newtcolorbox{objectifs}{popbase,colframe=popink,colback=poporangeL,
  before upper={\popboxhead{Objectifs de la leçon}{popink}{}}}
\newtcolorbox{prerequis}{popbase,colframe=popink,colback=popblueL,
  before upper={\popboxhead{Prérequis}{popblue}{}}}
\newtcolorbox{bilan}{popbase,colframe=popink,colback=white,boxrule=2.8pt,
  before upper={\popboxhead{Fiche bilan}{poporange}{L'essentiel à connaître pour l'examen}}}
% Figure encadrée avec légende
\newtcolorbox{popfigure}[1][]{enhanced,breakable=false,colback=white,colframe=popline,boxrule=1.4pt,arc=8pt,
  left=10pt,right=10pt,top=10pt,bottom=8pt,before skip=10pt,after skip=12pt,halign=center,
  lower separated=false,halign lower=center,
  fontlower=\popsemi\fontsize{9}{11.5}\selectfont\color{popmuted},
  #1}
\newcommand{\legende}[1]{\par\vspace{6pt}{\popsemi\fontsize{9}{11.5}\selectfont\color{popmuted}#1}}

% ============================================================
% Signature OnBuch+
% ============================================================
\newcommand{\poplogoinline}[1]{{\poplogo\fontsize{#1}{#1}\selectfont\color{popink}OnBuch\color{poporange}+}}
\newcommand{\signatureonbuch}{%
  \begin{tcolorbox}[enhanced,colback=popink,colframe=popink,boxrule=2.2pt,arc=10pt,
      drop shadow=poporange,left=14pt,right=14pt,top=10pt,bottom=10pt,before skip=0pt,after skip=0pt]
    \tikz[baseline=-0.7ex]\node[circle,fill=popgreen,draw=popcream,line width=1.3pt,minimum size=22pt,inner sep=0]{\popcheck{white}};%
    \hspace{9pt}\begin{minipage}[c]{0.62\linewidth}
      {\popxbold\fontsize{12}{13}\selectfont\color{popcream}Signé\ \textbullet\ {\poplogo OnBuch\color{poporange}+}}\par\vspace{2pt}
      {\raggedright\popsemi\fontsize{8}{10.5}\selectfont\color{popline}\MakeUppercase{Équipe pédagogique OnBuch+}\\\MakeUppercase{Conforme au programme officiel MINESEC}\par}
    \end{minipage}\hfill
    {\poplogo\fontsize{22}{22}\selectfont\color{popcream}OnBuch\color{poporange}+}
  \end{tcolorbox}%
}

% ============================================================
% Page de garde
% #1 titre · #2 matière · #3 niveau · #4 module · #5 n° leçon · #6 durée ·
% #7 description / objectifs (texte ou itemize)
% ============================================================
\newcommand{\pagegarde}[7]{%
  \thispagestyle{empty}
  \newgeometry{a4paper,margin=1.7cm,top=1.6cm,bottom=1.4cm}%
  \begin{tikzpicture}[remember picture,overlay]
    \fill[poporange!18] ([xshift=-3.2cm,yshift=-3.6cm]current page.north east) circle (4.6cm);
    \fill[poppurple!14] ([xshift=2.4cm,yshift=7.6cm]current page.south west) circle (3.8cm);
  \end{tikzpicture}%
  {\poplogo\fontsize{30}{30}\selectfont\color{popink}OnBuch\color{poporange}+}\hfill
  \raisebox{0.6ex}{\poppill{Cours signé \textbullet\ Premium}{popink}{popcream}}\par
  \vspace{1pt}\popsquiggle{poppurple}\par
  \vspace{8pt}
  \begin{tcolorbox}[enhanced,colback=poporange,colframe=popink,boxrule=2.8pt,arc=13pt,
      drop shadow=popink,left=18pt,right=18pt,top=12pt,bottom=13pt,after skip=10pt]
    \poppill{#2 \textbullet\ #3}{popink}{popcream}\hspace{5pt}\poppill{#4}{white}{popink}\par\vspace{8pt}
    {\popdisplay\fontsize{14}{14}\selectfont\color{popink}LEÇON #5}\par\vspace{4pt}
    {\raggedright\hyphenpenalty=10000\exhyphenpenalty=10000\popdisplay\fontsize{25}{27}\selectfont\color{popcream}\MakeUppercase{#1}\par}
    \vspace{8pt}
    {\popxbold\fontsize{9.5}{9.5}\selectfont\color{popink}\MakeUppercase{Durée conseillée : #6}}
  \end{tcolorbox}
  \begin{tcolorbox}[enhanced,colback=white,colframe=popink,boxrule=2.2pt,arc=10pt,
      drop shadow=popink,left=16pt,right=16pt,top=10pt,bottom=10pt,after skip=8pt]
    \poppill{Dans cette leçon}{poppurple}{white}\par\vspace{5pt}
    \popsemi\fontsize{9.3}{12.6}\selectfont\color{popink}#7
  \end{tcolorbox}
  \noindent\begin{minipage}[t]{0.487\linewidth}
    \begin{tcolorbox}[enhanced,colback=popgreen,colframe=popink,boxrule=2.2pt,arc=10pt,
        drop shadow=popink,left=11pt,right=11pt,top=10pt,bottom=10pt,height=3.1cm,valign=center]
      \tcbox[colback=white,colframe=popink,boxrule=1.3pt,arc=5pt,boxsep=0pt,nobeforeafter,
        left=3pt,right=3pt,top=3pt,bottom=3pt,valign=center]{\qrcode[height=1.9cm]{https://play.google.com/store/apps/details?id=com.luvvix.onbuch&hl=fr}}%
      \hspace{8pt}\begin{minipage}[c]{0.5\linewidth}
        {\popdisplay\fontsize{11}{12.5}\selectfont\color{white}\MakeUppercase{Télécharge OnBuch}}\par\vspace{4pt}
        {\raggedright\popsemi\fontsize{8.4}{11}\selectfont\color{white}Gratuit \textbullet\ Google Play\\Tuteur IA, annales, concours, résultats\par}
      \end{minipage}
    \end{tcolorbox}
  \end{minipage}\hfill
  \begin{minipage}[t]{0.487\linewidth}
    \begin{tcolorbox}[enhanced,colback=poppurple,colframe=popink,boxrule=2.2pt,arc=10pt,
        drop shadow=popink,left=11pt,right=11pt,top=10pt,bottom=10pt,height=3.1cm,valign=center]
      \tikz[baseline=-0.7ex]\node[circle,fill=white,draw=popink,line width=1.3pt,minimum size=1.6cm,inner sep=0]{\popdisplay\fontsize{24}{24}\selectfont\color{poppurple}+};%
      \hspace{8pt}\begin{minipage}[c]{0.6\linewidth}
        {\popdisplay\fontsize{11}{12.5}\selectfont\color{white}\MakeUppercase{Passe à OnBuch+}}\par\vspace{4pt}
        {\raggedright\popsemi\fontsize{8.4}{11}\selectfont\color{white}Tous les cours signés, Tuteur IA illimité, fiches PDF — onglet OnBuch+ de l'app\par}
      \end{minipage}
    \end{tcolorbox}
  \end{minipage}
  \vspace{8pt}
  \popcontenu
  \vfill
  \signatureonbuch
  \restoregeometry
  \newpage
}

% Rangée « Ce cours contient » (page de garde)
\newcommand{\poptile}[3]{% #1 couleur #2 gros texte #3 légende
  \begin{tcolorbox}[enhanced,colback=#1,colframe=popink,boxrule=2pt,arc=9pt,shadow={2.5pt}{-2.5pt}{0pt}{popink},
      left=6pt,right=6pt,top=9pt,bottom=9pt,halign=center,valign=center,height=1.75cm,width=\linewidth,nobeforeafter]
    {\popdisplay\fontsize{12.5}{13}\selectfont\color{white}\MakeUppercase{#2}}\par\vspace{3pt}
    {\centering\popsemi\fontsize{7.8}{9.5}\selectfont\color{white}#3\par}
  \end{tcolorbox}}
\newcommand{\popcontenu}{%
  \par\noindent{\popxboldls\fontsize{7.5}{7.5}\selectfont\color{popmuted}CE COURS SIGNÉ CONTIENT}\par\vspace{7pt}
  \noindent\begin{minipage}[t]{0.235\linewidth}\poptile{popblue}{Cours}{complet, illustré, pas à pas}\end{minipage}\hfill
  \begin{minipage}[t]{0.235\linewidth}\poptile{poporange}{Méthodes}{et exercices résolus}\end{minipage}\hfill
  \begin{minipage}[t]{0.235\linewidth}\poptile{poppink}{Exercices}{type examen + corrigés}\end{minipage}\hfill
  \begin{minipage}[t]{0.235\linewidth}\poptile{popgreen}{Fiche bilan}{l'essentiel à retenir}\end{minipage}\par}

% Sommaire
\newtcolorbox{sommaire}{enhanced,colback=white,colframe=popink,boxrule=2.2pt,arc=10pt,
  drop shadow=popink,left=18pt,right=18pt,top=13pt,bottom=13pt,before skip=6pt,after skip=20pt,
  before upper={\poppill{Sommaire}{poppurple}{white}\par\vspace{9pt}\popsemi\fontsize{10.6}{14.5}\selectfont\color{popink}}}

% Signature de fin de cours — poussée tout en bas de la dernière page
\newcommand{\findecours}{%
  \par\vfill\nopagebreak
  \begin{center}\popsquiggle{poporange}\end{center}\vspace{4pt}
  \signatureonbuch
}
\AtBeginDocument{\def\llbracket{\mathopen{[\mkern-3.5mu[}}\def\rrbracket{\mathclose{]\mkern-3.5mu]}}} % intervalles d'entiers (glyphe ⟦ absent de la police)
% Glyphes absents de Fira Math : redéfinis à partir de glyphes disponibles
\AtBeginDocument{%
  \def\setminus{\mathbin{\backslash}}\let\smallsetminus\setminus
  \def\vdots{\mathord{\vbox{\baselineskip3.2pt\lineskiplimit0pt\kern4pt\hbox{.}\hbox{.}\hbox{.}}}}%
  \def\ddots{\mathinner{\mkern1mu\raise6.5pt\hbox{.}\mkern2mu\raise3.6pt\hbox{.}\mkern2mu\raise0.7pt\hbox{.}\mkern1mu}}%
  \def\triangle{\mathord{\scalebox{0.9}{$\symup{\Delta}$}}}%
  \def\nabla{\mathord{\raisebox{\depth}{\scalebox{1}[-1]{$\symup{\Delta}$}}}}%
  \def\checkmark{\popcheck{popdark}}%
  \def\star{\mathbin{\ast}}\def\bigstar{\mathord{\ast}}%
  \def\diamond{\mathbin{\scalebox{0.6}{$\Diamond$}}}%
  \def\bigcup{\mathop{\mathchoice{\raisebox{-0.3em}{\scalebox{1.7}{$\cup$}}}{\scalebox{1.25}{$\cup$}}{\cup}{\cup}}\displaylimits}%
  \def\bigcap{\mathop{\mathchoice{\raisebox{-0.3em}{\scalebox{1.7}{$\cap$}}}{\scalebox{1.25}{$\cap$}}{\cap}{\cap}}\displaylimits}%
  \def\lesseqqgtr{\mathrel{\lessgtr}}\def\bigoplus{\mathop{\oplus}\displaylimits}%
}
\providecommand{\ssi}{si et seulement si} % abréviation fréquente
\AtBeginDocument{\providecommand{\wideparen}{\overparen}} % arc de cercle

%% FIGFIX-BEGIN v5 — correctifs globaux des figures (docs/onbuch-generation/tools/figfix)
\usepackage{adjustbox}\usepackage{etoolbox}\tcbuselibrary{raster}
% toute figure TikZ/pgfplots plus large que la ligne est réduite à la largeur disponible
\BeforeBeginEnvironment{tikzpicture}{\begin{adjustbox}{max width=\linewidth}}
\AfterEndEnvironment{tikzpicture}{\end{adjustbox}}
% légendes pgfplots lisibles par-dessus les courbes
\pgfplotsset{every axis/.append style={legend style={fill=white,fill opacity=0.92,text opacity=1,draw=black!15,inner sep=3pt}}}
% étiquettes d'arêtes (syntaxe edge["..."]) avec fond blanc
\tikzset{every edge quotes/.append style={fill=white,inner sep=1.2pt}}
%% FIGFIX-END
\begin{document}

\pagegarde{Modéliser un phénomène : exemples de modèles utilisés en physique et en chimie}{Physique}{1ère TI}{Module 1 : Mesures et incertitudes}{N02}{4 h}{Cette leçon initie les élèves à la démarche de modélisation en physique et en chimie : formuler des hypothèses simplificatrices, établir une loi mathématique à partir de mesures et tester son domaine de validité. À travers la loi d'Ohm et la loi des gaz parfaits, les élèves apprennent à travailler comme des scientifiques en collectant et en traitant des données expérimentales.
\vspace{4pt}
\begin{multicols}{2}\small
\begin{itemize}
\item À la fin de cette leçon, je sais définir ce qu'est un modèle en physique et expliquer son rôle
\item À la fin de cette leçon, je sais formuler des hypothèses simplificatrices pour modéliser une situation concrète
\item À la fin de cette leçon, je sais énoncer la loi d'Ohm et préciser ses conditions de validité
\item À la fin de cette leçon, je sais énoncer la loi des gaz parfaits et identifier les hypothèses du modèle du gaz parfait
\item À la fin de cette leçon, je sais tracer une caractéristique U = f(I) ou P = f(1/V) et en déduire une loi mathématique simple
\item À la fin de cette leçon, je sais tester la validité d'une hypothèse ou d'une loi par la mesure et l'analyse d'un graphique
\end{itemize}\end{multicols}}

\begin{sommaire}
\begin{multicols}{2}
\begin{enumerate}
\item Qu'est-ce qu'un modèle en physique ?
\item Établir une loi à partir de mesures : la démarche expérimentale
\item Exemple 1 : la loi d'Ohm
\item Exemple 2 : la loi des gaz parfaits
\item Contraintes d'une loi : domaine de validité et limites des modèles
\item Synthèse : travailler comme un scientifique
\item Activité d'intégration
\item Méthodes et exercices résolus
\item Exercices
\item Corrigés des exercices
\item Fiche bilan
\end{enumerate}
\end{multicols}
\end{sommaire}

\begin{objectifs}
\begin{itemize}
\item À la fin de cette leçon, je sais définir ce qu'est un modèle en physique et expliquer son rôle
\item À la fin de cette leçon, je sais formuler des hypothèses simplificatrices pour modéliser une situation concrète
\item À la fin de cette leçon, je sais énoncer la loi d'Ohm et préciser ses conditions de validité
\item À la fin de cette leçon, je sais énoncer la loi des gaz parfaits et identifier les hypothèses du modèle du gaz parfait
\item À la fin de cette leçon, je sais tracer une caractéristique U = f(I) ou P = f(1/V) et en déduire une loi mathématique simple
\item À la fin de cette leçon, je sais tester la validité d'une hypothèse ou d'une loi par la mesure et l'analyse d'un graphique
\item À la fin de cette leçon, je sais identifier le domaine de validité d'une loi et repérer quand un modèle cesse de s'appliquer
\item À la fin de cette leçon, je sais exploiter une loi de modélisation pour prévoir la valeur d'une grandeur (calcul, extrapolation raisonnée)
\end{itemize}
\end{objectifs}

\begin{prerequis}
\begin{itemize}
\item Grandeurs physiques, unités du Système international et conversions (Seconde, rappels de début d'année)
\item Notion de proportionnalité et équation d'une droite y = ax (mathématiques, Seconde et Première)
\item Lecture et tracé de graphiques à partir d'un tableau de mesures
\item Notions de tension, intensité, résistance, pression, volume et température (Seconde)
\item Incertitudes et chiffres significatifs (Leçon N01 du Module 1)
\end{itemize}
\end{prerequis}

\newpage

\coursec{Qu'est-ce qu'un modèle en physique ?}

\courssub{Une situation d'accroche : prévoir sans casser}

À Yaoundé, un technicien de la SONEL constate qu'en branchant plusieurs climatiseurs sur la même ligne, la tension chute et les appareils chauffent : comment prévoir ce phénomène sans tout casser ? De même, un chauffeur de taxi-moto sait que la pression de ses pneus augmente après un long trajet sous le soleil de Douala. Derrière ces observations se cachent des lois simples — la loi d'Ohm, la loi des gaz parfaits — qui permettent de prévoir, calculer et dimensionner. Mais ces lois sont-elles toujours vraies ? Quelles hypothèses se cachent derrière elles ? C'est tout l'art du physicien : \cle{modéliser} la réalité pour la comprendre et la prévoir.

\textbf{Problématique.} Comment le physicien transforme-t-il une situation réelle, toujours compliquée, en un objet simple sur lequel on peut calculer ? Quelles garanties avons-nous que ce calcul décrira bien la réalité ?

\courssub{Définition d'un modèle}

Commençons par l'intuition. Quand tu regardes la carte routière du Cameroun pour préparer un voyage Douala--Bafoussam, tu ne vois ni les manguiers, ni les marchés, ni les nids-de-poule : la carte est une \emph{représentation simplifiée} du territoire. Elle est « fausse » au sens où elle a oublié presque tout, mais elle est \emph{utile} car elle permet de prévoir la distance et le temps de trajet. Un modèle en physique fonctionne exactement ainsi.

\begin{definition}[Modèle]
Un \cle{modèle} est une représentation simplifiée d'un objet ou d'un phénomène réel, construite à partir d'\cle{hypothèses} explicites, qui permet de \textbf{décrire}, d'\textbf{expliquer} et de \textbf{prévoir} son comportement.
\end{definition}

Trois verbes sont à retenir : un bon modèle \emph{décrit} ce qui se passe, \emph{explique} pourquoi cela se passe, et surtout \emph{prévoit} ce qui se passera dans une situation nouvelle. C'est cette capacité de prévision qui donne toute sa valeur au modèle : le technicien de la SONEL peut calculer la chute de tension \emph{avant} de brancher les climatiseurs.

\begin{attention}[Modèle et loi : ne pas confondre]
Une \cle{loi} est une \textbf{relation mathématique entre des grandeurs mesurables} (par exemple $U = R\,I$ pour la loi d'Ohm). Un \cle{modèle} est plus large : c'est l'ensemble formé par les hypothèses simplificatrices, l'objet idéalisé et les lois qui s'y appliquent. Ainsi, le \emph{modèle du gaz parfait} (hypothèses : molécules ponctuelles, sans interaction) contient la \emph{loi des gaz parfaits} $PV = nRT$. La loi est la « formule » ; le modèle est tout l'édifice intellectuel qui la justifie et précise quand on peut l'utiliser.
\end{attention}

\courssub{Les trois étapes de la démarche de modélisation}

Modéliser n'est pas un acte magique : c'est une démarche organisée, que tout scientifique applique, souvent sans s'en rendre compte.

\textbf{Étape 1 : observer le phénomène.} On commence par regarder la réalité et par identifier les \cle{grandeurs pertinentes}. Pour la chute d'une mangue : sa masse, sa hauteur, sa vitesse, le temps. Pour la ligne électrique : la tension, l'intensité, la résistance des câbles.

\textbf{Étape 2 : formuler des hypothèses simplificatrices.} On décide \emph{explicitement} de ce que l'on néglige : frottements de l'air, résistance des fils de connexion, variation de température, etc. Ces hypothèses doivent être \textbf{énoncées clairement} : c'est elles qui définissent le modèle.

\textbf{Étape 3 : confronter le modèle à l'expérience.} Le modèle fournit des prévisions chiffrées ; on les compare à des mesures. Si l'écart est faible, le modèle est validé \emph{dans ce domaine}. Sinon, on revient aux hypothèses : on en modifie une, et on recommence. La modélisation est donc un \textbf{processus itératif}.

\begin{popfigure}
\begin{center}
\begin{tikzpicture}[font=\small,
  bloc/.style={draw=popblue, fill=popblueL, rounded corners, thick, minimum height=1.1cm, minimum width=2.6cm, align=center},
  fleche/.style={-{Stealth[length=3mm]}, very thick, popdark},
  etiq/.style={midway, above, font=\footnotesize\itshape, popdark, align=center}]
\node[bloc, align=center] (R) at (0,0){\textbf{Réalité}\\ phénomène observé};
\node[bloc, fill=poporangeL, draw=poporange, align=center] (M) at (5.2,0){\textbf{Modèle}\\ objet idéalisé};
\node[bloc, fill=popgreenL, draw=popgreen, align=center] (P) at (10.4,0){\textbf{Prévisions}\\ résultats calculés};
\draw[fleche] (R) -- (M) node[etiq, align=center]{hypothèses\\ simplificatrices};
\draw[fleche] (M) -- (P) node[etiq, align=center]{loi\\ mathématique};
\draw[fleche, poppink, dashed] (P.south) .. controls (7,-2.2) and (3.4,-2.2) .. (R.south)
  node[midway, below, font=\footnotesize\itshape, poppink, align=center] {confrontation à l'expérience\\ (mesures) : validation ou retour aux hypothèses};
\end{tikzpicture}
\end{center}
\legende{Le cycle de la modélisation : de la réalité au modèle, du modèle aux prévisions, puis retour à la réalité par la mesure.}
\end{popfigure}

\begin{methode}[Modéliser une situation concrète]
\begin{enumerate}
\item \textbf{Identifier} les grandeurs pertinentes du phénomène (celles qui varient et qu'on peut mesurer).
\item \textbf{Formuler} des hypothèses simplificatrices explicites (que néglige-t-on ? que suppose-t-on constant ?).
\item \textbf{Établir} (ou utiliser) une relation mathématique entre ces grandeurs : c'est la loi du modèle.
\item \textbf{Tester} la loi par la mesure : comparer prévisions et résultats expérimentaux, puis conclure sur la validité du modèle.
\end{enumerate}
\end{methode}

\courssub{Quelques modèles célèbres que tu connais déjà}

Tu as déjà utilisé des modèles sans toujours le savoir. En voici quatre, que nous retrouverons tout au long de l'année.

\textbf{Le modèle du \cle{point matériel}.} Pour étudier le mouvement d'une moto-taxi entre Mokolo et le centre-ville, on assimile la moto (et son pilote !) à un simple point doté d'une masse. On néglige sa taille, sa forme, ses roues qui tournent. Ce modèle est excellent pour décrire la trajectoire, mais il serait absurde pour étudier le capotage de la moto dans un virage : là, la répartition de la masse compte.

\textbf{Le modèle du \cle{gaz parfait}.} On suppose que les molécules d'un gaz sont des points sans volume propre, sans interaction entre elles, n'échangeant de l'énergie que par chocs. Ce modèle conduit à la loi $PV = nRT$, remarquablement efficace pour l'air des pneus du taxi-moto à pression ordinaire.

\textbf{Le modèle du \cle{dipôle ohmique} (résistor).} On suppose qu'un conducteur s'oppose au passage du courant avec une résistance $R$ constante, quelle que soit la tension : $U = RI$. Ce modèle oublie que $R$ dépend en fait de la température, donc de l'échauffement du conducteur.

\textbf{Le modèle \cle{corpusculaire} de la matière.} La matière est faite de particules (atomes, molécules, ions) en mouvement incessant. Ce modèle explique les états solide, liquide, gazeux et les changements d'état étudiés en chimie.

\courssub{Le rôle des hypothèses simplificatrices}

Les hypothèses simplificatrices sont le cœur du modèle. En voici quelques-unes que tu rencontreras souvent :

\begin{itemize}
\item \emph{« On néglige les frottements »} : l'air, les contacts, les pivots sont supposés parfaits ;
\item \emph{« La température est supposée constante »} : transformation isotherme ;
\item \emph{« Les fils de connexion ont une résistance nulle »} : toute la tension se retrouve aux bornes des dipôles ;
\item \emph{« Le solide est indéformable »} : sa forme ne change pas sous l'action des forces ;
\item \emph{« Le champ de pesanteur est uniforme »} : $\vec{g}$ garde même direction, même sens et même norme à l'échelle de l'expérience.
\end{itemize}

\begin{attention}[Simplifier ne veut pas dire fausser]
Erreur fréquente : croire qu'une hypothèse simplificatrice « rend le problème faux ». Non ! Une \textbf{bonne} hypothèse néglige des effets \emph{petits devant le phénomène étudié}, jamais le phénomène lui-même. Négliger les frottements de l'air sur une mangue qui tombe de 3 m est légitime (l'effet est de l'ordre de quelques pourcents) ; négliger le poids de la mangue serait absurde, car c'est précisément le poids qui provoque la chute. Un bon modèle simplifie la réalité sans la trahir.
\end{attention}

\courssub{Exemple guidé : la chute d'une mangue, avec et sans modèle}

Une mangue de masse $m = \SI{0,15}{kg}$ se détache d'une branche située à $h = \SI{3,0}{m}$ du sol. Comparons la situation réelle et la situation modélisée.

\begin{experience}[Du réel au modèle : la chute de la mangue]
\textbf{Situation réelle.} La mangue subit son poids $\vec{P}$, la poussée d'Archimède de l'air $\vec{\Pi}$ et la force de frottement de l'air $\vec{f}$, qui dépend de la vitesse. L'équation du mouvement est compliquée.

\textbf{Hypothèses du modèle.} (H1) La mangue est un point matériel $M$ de masse $m$ ; (H2) on néglige la poussée d'Archimède et les frottements de l'air ; (H3) $\vec{g}$ est uniforme, $g = \SI{9,8}{N.kg^{-1}}$.

\textbf{Conséquence.} La mangue n'est soumise qu'à son poids : c'est le modèle de la \cle{chute libre}. La deuxième loi de Newton donne $m\vec{a} = m\vec{g}$, soit $\vec{a} = \vec{g}$ : mouvement rectiligne uniformément accéléré. En intégrant deux fois avec une vitesse initiale nulle :
\[ z(t) = h - \tfrac{1}{2}g\,t^2 \qquad\text{et}\qquad v(t) = g\,t . \]
\end{experience}

\begin{exoresolu}[Durée de chute de la mangue]
Énoncé. Dans le modèle de la chute libre, calculer la durée de chute $t_c$ de la mangue ($h = \SI{3,0}{m}$, $g = \SI{9,8}{N.kg^{-1}}$) et sa vitesse d'arrivée au sol. Commenter.
\tcblower
Solution. La mangue touche le sol quand $z(t_c) = 0$, c'est-à-dire $h = \tfrac{1}{2}g\,t_c^2$. On en tire :
\[ t_c = \sqrt{\dfrac{2h}{g}} = \sqrt{\dfrac{2 \times 3,0}{9,8}} = \sqrt{0,612} \approx \SI{0,78}{s}. \]
Vérification dimensionnelle : $\sqrt{\dfrac{[L]}{[L][T]^{-2}}} = \sqrt{[T]^2} = [T]$ : le résultat est bien un temps.
La vitesse au sol vaut :
\[ v_c = g\,t_c = 9,8 \times 0,78 \approx \SI{7,7}{m.s^{-1}} \approx \SI{28}{km.h^{-1}}. \]
Commentaire. La mesure réelle donnerait un temps légèrement supérieur (les frottements ralentissent la chute) : l'écart, de quelques pourcents, reste faible car la mangue est dense et la hauteur modeste. Le modèle est donc \emph{validé} dans cette situation. Pour une feuille de papier, le même modèle échouerait spectaculairement : il faudrait changer d'hypothèses.
\end{exoresolu}

\begin{popfigure}
\begin{center}
\begin{tikzpicture}[font=\footnotesize, scale=0.92]
% Colonne gauche : réalité
\node[draw=poppink, fill=poppinkL, rounded corners, thick, minimum width=5.6cm, minimum height=0.9cm] at (0,3.4) {\textbf{Situation réelle}};
\draw[thick, popdark] (-2.6,0) -- (2.6,0) node[below left]{sol};
\draw[thick, popgreen] (0,2.2) circle (0.28);
\node at (0,2.2) {mangue};
\draw[-{Stealth}, very thick, popblue] (0,2.2) -- (0,0.9) node[midway, right]{$\vec{P}$};
\draw[-{Stealth}, thick, poporange] (0.35,2.4) -- (0.35,2.9) node[right]{$\vec{\Pi}$};
\draw[-{Stealth}, thick, poppink] (-0.35,2.4) -- (-0.35,3.0) node[left]{$\vec{f}$ (frottements)};
\node[align=center, popmuted] at (0,-0.9) {3 forces, force de frottement\\ dépendant de la vitesse :\\ calcul difficile};
% Colonne droite : modèle
\node[draw=popgreen, fill=popgreenL, rounded corners, thick, minimum width=5.6cm, minimum height=0.9cm] at (9,3.4) {\textbf{Situation modélisée (chute libre)}};
\draw[thick, popdark] (6.4,0) -- (11.6,0) node[below left]{sol};
\fill[popgreen] (9,2.2) circle (0.09);
\node[right] at (9.15,2.2) {point $M$ de masse $m$};
\draw[-{Stealth}, very thick, popblue] (9,2.2) -- (9,0.9) node[midway, right]{$\vec{P} = m\vec{g}$};
\node[align=center, popmuted] at (9,-0.9) {1 seule force : $\vec{a} = \vec{g}$,\\ $t_c = \sqrt{2h/g} \approx \SI{0,78}{s}$};
\end{tikzpicture}
\end{center}
\legende{Situation réelle (trois forces) contre situation modélisée (une seule force) : l'hypothèse « frottements négligés » rend le calcul possible.}
\end{popfigure}

\courssub{Un modèle n'est jamais « vrai » en absolu}

Voici peut-être l'idée la plus importante de cette leçon. Un modèle n'est ni vrai ni faux : il est \textbf{plus ou moins adapté} à la situation étudiée. Le modèle du point matériel est parfait pour la moto sur l'axe Douala--Yaoundé, inutilisable pour étudier son équilibre dans un virage. Le modèle du gaz parfait décrit bien l'air d'un pneu, mais devient mauvais pour un gaz très comprimé ou très froid, où les molécules ne peuvent plus être considérées comme des points sans interaction.

C'est pourquoi tout énoncé de physique sérieux précise ses hypothèses : « on néglige les frottements », « le gaz est supposé parfait », « la résistance du fil est négligeable ». Ces phrases ne sont pas des détails : elles définissent le \cle{domaine de validité} du raisonnement qui suit. Nous y reviendrons en détail dans la section 5.

\begin{savaistu}[« Tous les modèles sont faux, mais certains sont utiles »]
Cette phrase célèbre est du statisticien britannique George Box (1919--2013). Elle résume la philosophie de la modélisation : puisque tout modèle simplifie la réalité, aucun ne la reproduit exactement — tous sont donc « faux » au sens strict. Mais la valeur d'un modèle ne se juge pas à sa « vérité » : elle se juge à la \textbf{qualité de ses prédictions}. Les ingénieurs qui ont dimensionné le barrage de Lom Pangar ont utilisé des modèles simplifiés de l'écoulement de l'eau et de la résistance des matériaux : ces modèles étaient « faux », mais suffisamment bons pour que l'ouvrage tienne. C'est exactement la démarche que tu apprends cette année.
\end{savaistu}

\begin{aretenir}[L'essentiel de la section]
\begin{itemize}
\item Un \cle{modèle} est une représentation simplifiée de la réalité, fondée sur des hypothèses explicites, qui permet de décrire, expliquer et prévoir un phénomène.
\item La démarche comporte trois étapes : \textbf{observer}, \textbf{formuler des hypothèses simplificatrices}, \textbf{confronter le modèle à l'expérience} (cycle itératif).
\item Une \cle{loi} est une relation mathématique entre grandeurs ; elle vit à l'intérieur d'un modèle qui précise ses conditions d'application.
\item Une bonne hypothèse néglige des effets petits devant le phénomène étudié, jamais le phénomène lui-même.
\item Un modèle n'est jamais vrai en absolu : il est adapté ou non à la situation, et sa valeur se mesure à la qualité de ses prédictions.
\end{itemize}
\end{aretenir}

\coursec{Établir une loi à partir de mesures : la démarche expérimentale}

Dans la section précédente, nous avons vu qu'un \cle{modèle} est une représentation simplifiée de la réalité, construite à partir d'hypothèses, et qu'une \cle{loi physique} est une relation mathématique entre grandeurs mesurables. Mais une question demeure : comment un scientifique passe-t-il concrètement d'une série de mesures réalisées au laboratoire à une formule comme $U = R\,I$ ou $P\,V = \text{constante}$ ? C'est toute la \cle{démarche expérimentale}, véritable colonne vertébrale du métier de physicien, que nous allons détailler ici. Cette démarche est exigible au Probatoire : sais-tu qu'au XIX\textsuperscript{e} siècle, c'est exactement ainsi que Georg Ohm, professeur de lycée à Cologne, a découvert la loi qui porte son nom, en mesurant patiemment tensions et courants sur des fils de cuivre ?

\courssub{Le protocole général : isoler les grandeurs en jeu}

\textbf{Situation concrète.} Un élève du lycée de Ngaoundéré veut savoir comment la tension $U$ aux bornes d'un conducteur dépend de l'intensité $I$ qui le traverse. Il remarque vite que $U$ pourrait aussi dépendre de la température du fil, de sa longueur, de son diamètre... S'il fait tout varier en même temps, il sera incapable de conclure quoi que ce soit !

\begin{definition}[Variable contrôlée, variable mesurée, paramètres fixés]
Dans une étude expérimentale de la relation entre deux grandeurs $x$ et $y$ :
\begin{itemize}
\item la \cle{variable contrôlée} (ou variable indépendante) $x$ est la grandeur que l'expérimentateur fait varier volontairement, de valeurs connues en valeurs connues ;
\item la \cle{variable mesurée} (ou variable dépendante) $y$ est la grandeur que l'on mesure pour chaque valeur de $x$ ;
\item toutes les autres grandeurs susceptibles d'influencer $y$ sont maintenues \textbf{constantes} : ce sont les \cle{paramètres fixés}.
\end{itemize}
\end{definition}

\begin{methode}[Protocole général d'établissement d'une loi]
\begin{enumerate}
\item \textbf{Identifier} les grandeurs en jeu et formuler une hypothèse : « $y$ dépend probablement de $x$ ».
\item \textbf{Choisir} la variable contrôlée $x$ et fixer tous les autres paramètres (température, longueur du fil, nature du gaz, etc.).
\item \textbf{Faire varier} $x$ sur une plage aussi large que possible (au moins 5 à 6 valeurs régulièrement espacées, en incluant si possible $x = 0$).
\item \textbf{Mesurer} $y$ pour chaque valeur de $x$, en répétant si possible chaque mesure, et consigner les résultats dans un tableau avec leurs unités et leurs incertitudes.
\item \textbf{Tracer} le graphique $y = f(x)$, identifier la forme de la courbe, en déduire la loi mathématique.
\item \textbf{Valider} la loi sur de nouvelles mesures non utilisées pour l'établir.
\end{enumerate}
\end{methode}

\begin{attention}[La règle d'or : une seule grandeur varie à la fois]
Si deux paramètres varient simultanément, il est impossible d'attribuer la variation de $y$ à l'un ou à l'autre. Par exemple, pour étudier $P$ en fonction de $V$ pour un gaz, la température $T$ et la quantité de matière $n$ doivent rester rigoureusement constantes : on utilise une seringue fermée, manipulée lentement pour éviter l'échauffement du gaz. Oublier cette précaution est la cause première des « lois » fausses trouvées en TP !
\end{attention}

\courssub{Le tableau de mesures : unités et incertitudes}

Toute mesure expérimentale est entachée d'une \cle{incertitude} : l'appareil a une précision limitée, l'opérateur lit le cadran avec une marge d'erreur, la grandeur elle-même peut fluctuer. On écrit donc chaque résultat sous la forme
\[ y = y_{\text{mes}} \pm \Delta y \]
où $\Delta y$ est l'incertitude absolue, exprimée dans la même unité que $y$.

\begin{exemplebox}[Tableau de mesures : tension et intensité aux bornes d'un conducteur]
Un groupe d'élèves mesure la tension $U$ aux bornes d'un conducteur ohmique pour différentes intensités $I$, à température constante. Le voltmètre a une incertitude $\Delta U = \SI{0,1}{V}$, l'ampèremètre $\Delta I = \SI{0,01}{A}$.
\begin{center}
\begin{tabularx}{\linewidth}{|l|X|X|X|X|}
\hline
\rowcolor{popblueL} Mesure n\textdegree & 1 & 2 & 3 & 4 \\
\hline
$I$ (A) & $0 \pm 0{,}01$ & $0{,}10 \pm 0{,}01$ & $0{,}20 \pm 0{,}01$ & $0{,}30 \pm 0{,}01$ \\
\hline
$U$ (V) & $0 \pm 0{,}1$ & $2{,}1 \pm 0{,}1$ & $4{,}0 \pm 0{,}1$ & $6{,}2 \pm 0{,}1$ \\
\hline
\end{tabularx}
\end{center}
\end{exemplebox}

\begin{attention}[Rédaction du tableau]
Chaque colonne ou ligne de grandeurs doit porter le \textbf{symbole de la grandeur et son unité} : on écrit « $U$ (V) » et non « $U$ » seul. Les valeurs d'un tableau sont des nombres sans unité, puisque l'unité est indiquée dans l'en-tête. Oublier l'unité dans l'en-tête est sanctionné au Probatoire.
\end{attention}

\courssub{Le tracé du graphique : axes, échelle, barres d'incertitude}

Le graphique est l'outil central de la démarche : l'œil humain reconnaît instantanément une droite, alors qu'un tableau de nombres ne dit rien d'évident.

\begin{methode}[Règles de tracé d'un graphique scientifique]
\begin{enumerate}
\item Porter la \textbf{variable contrôlée} $x$ en abscisse et la \textbf{variable mesurée} $y$ en ordonnée.
\item \textbf{Graduer} chaque axe, indiquer la grandeur et son unité : « $U$ (V) ».
\item Choisir l'\textbf{échelle} de sorte que les points occupent la majeure partie de la feuille (pas un petit nuage tassé dans un coin).
\item Placer chaque \textbf{point expérimental} par une croix fine.
\item Tracer les \cle{barres d'incertitude} : un segment horizontal de demi-longueur $\Delta x$ et un segment vertical de demi-longueur $\Delta y$ centrés sur le point. Le rectangle ainsi délimité est la zone où se trouve « vraiment » le point.
\end{enumerate}
\end{methode}

\begin{popfigure}
\begin{center}
\begin{tikzpicture}
\begin{axis}[
width=0.85\linewidth, height=7cm,
xlabel={$I$ (A)}, ylabel={$U$ (V)},
xmin=0, xmax=0.35, ymin=0, ymax=7,
grid=both, grid style={popline!40},
legend style={at={(0.03,0.97)}, anchor=north west, font=\small},
]
\addplot[popcurve, domain=0:0.35, samples=50] {20*x};
\addlegendentry{Droite modèle $U = 20\,I$}
\addplot[only marks, mark=+, mark size=3pt, poporange, thick,
error bars/.cd, x dir=both, x explicit, y dir=both, y explicit,
error bar style={popdark}]
coordinates {
(0,0) +- (0.01,0.1)
(0.10,2.1) +- (0.01,0.1)
(0.20,4.0) +- (0.01,0.1)
(0.30,6.2) +- (0.01,0.1)
};
\addlegendentry{Points expérimentaux}
\end{axis}
\end{tikzpicture}
\end{center}
\legende{Tension $U$ en fonction de l'intensité $I$ : les points expérimentaux, munis de leurs barres d'incertitude, sont alignés sur une droite passant par l'origine.}
\end{popfigure}

\courssub{Identifier la forme de la relation}

Une fois le graphique tracé, trois cas se présentent :

\begin{aretenir}[Les trois formes fondamentales]
\begin{itemize}
\item \textbf{Droite passant par l'origine} : $y$ et $x$ sont \cle{proportionnelles} ; la loi s'écrit $y = a\,x$.
\item \textbf{Droite ne passant pas par l'origine} : relation \cle{affine} $y = a\,x + b$ ; $b$ est l'ordonnée à l'origine, souvent porteuse d'un sens physique (force électromotrice d'un générateur, par exemple).
\item \textbf{Courbe} : relation \cle{non linéaire} ; il faut alors chercher un changement de variable pour la « linéariser » (voir plus loin).
\end{itemize}
\end{aretenir}

\begin{propriete}[Critère de proportionnalité (admis)]
Si les points expérimentaux sont alignés sur une droite passant par l'origine, \textbf{aux incertitudes de mesure près} (c'est-à-dire si la droite traverse le rectangle d'incertitude de chaque point), alors les grandeurs $x$ et $y$ sont proportionnelles et la loi s'écrit
\[ y = a \cdot x \]
où $a$ est le \cle{coefficient directeur} de la droite, constante caractéristique du système étudié.
\end{propriete}

\courssub{Détermination du coefficient directeur et interprétation physique}

Le coefficient directeur se calcule à partir de deux points \textbf{de la droite modèle} (et non de deux points expérimentaux), choisis éloignés l'un de l'autre pour minimiser l'erreur de lecture :
\[ a = \frac{y_2 - y_1}{x_2 - x_1} \]

\begin{attention}[Analyse dimensionnelle obligatoire]
Le coefficient directeur possède en général une \textbf{unité} : celle du quotient $\dfrac{[y]}{[x]}$. Ne jamais écrire « $a = 20$ » sans unité ! Cette unité est souvent la signature physique de la grandeur modélisée : ici, des volts par ampère, c'est-à-dire des ohms.
\end{attention}

\begin{exoresolu}[Établissement de la loi $U = R\,I$ à partir des mesures]
Énoncé. À partir du tableau de mesures ci-dessus, montrer que $U$ est proportionnelle à $I$, déterminer le coefficient directeur et écrire la loi. Vérifier l'homogénéité.
\tcblower
Solution détaillée.
\begin{enumerate}
\item \textbf{Graphique} : le tracé de $U$ en fonction de $I$ (figure ci-dessus) montre des points alignés sur une droite passant par l'origine ; la droite traverse toutes les barres d'incertitude. On conclut à la \textbf{proportionnalité} de $U$ et $I$.
\item \textbf{Coefficient directeur} : on choisit deux points éloignés de la droite modèle, par exemple l'origine $(0\,;\,0)$ et le point de la droite d'abscisse $\SI{0,30}{A}$ (soit une ordonnée de $\SI{6,0}{V}$). \emph{Attention : on ne prend pas la valeur brute $6,2\,\text{V}$ du tableau, mais la valeur lue sur la droite modèle.}
\[ a = \frac{6{,}0 - 0}{0{,}30 - 0} = \SI{20}{V.A^{-1}} = \SI{20}{\ohm} \]
car $1~\Omega = 1~\text{V}\cdot\text{A}^{-1}$ par définition de l'ohm.
\item \textbf{Loi} : $\boxed{U = 20 \times I}$ avec $U$ en volts et $I$ en ampères. Le coefficient $a = \SI{20}{\ohm}$ s'identifie à la \textbf{résistance} $R$ du conducteur.
\item \textbf{Homogénéité} : $[R\,I] = \Omega \times \text{A} = (\text{V}\cdot\text{A}^{-1}) \times \text{A} = \text{V} = [U]$. La relation est homogène. 
\item \textbf{Validation} : pour une nouvelle mesure $I = \SI{0,25}{A}$, la loi prédit $U = 20 \times 0{,}25 = \SI{5,0}{V}$ ; une mesure donnant $U = \SI{5,1}{V} \pm \SI{0,1}{V}$ confirme la loi dans les limites des incertitudes.
\end{enumerate}
\end{exoresolu}

\begin{attention}[Erreur fréquente : la droite « au jugé »]
Il est interdit de tracer la droite en la forçant à passer par \textbf{tous} les points expérimentaux, ni en reliant les points en segments brisés (« courbe en dents de scie »). Les points sont entachés d'incertitudes : la droite \textbf{modèle} est la droite qui passe \emph{au plus près} de l'ensemble des points, en traversant si possible tous les rectangles d'incertitude. De même, on ne calcule jamais le coefficient directeur avec deux points expérimentaux bruts : on utilise deux points \emph{de la droite tracée}.
\end{attention}

\courssub{Cas d'une relation non linéaire : la linéarisation}

Que faire quand les points forment une courbe ? Prenons l'exemple historique de la loi de Boyle-Mariotte : à température constante, on mesure la pression $P$ d'un gaz pour différents volumes $V$ (seringue fermée). Le tracé de $P$ en fonction de $V$ donne une branche d'hyperbole : impossible d'en tirer directement une loi simple.

L'astuce du scientifique consiste à effectuer un \cle{changement de variable} : si l'on soupçonne que $P\,V = \text{constante}$, alors $P = \dfrac{\text{constante}}{V}$, donc $P$ est proportionnelle à $\dfrac{1}{V}$. On trace alors $P$ en fonction de $x = \dfrac{1}{V}$ : si les points s'alignent sur une droite passant par l'origine, l'hypothèse est validée.

\begin{popfigure}
\begin{center}
\begin{tikzpicture}
\begin{axis}[
width=0.48\linewidth, height=6cm,
xlabel={$V$ (L)}, ylabel={$P$ (kPa)},
xmin=0, xmax=2.6, ymin=0, ymax=210,
title={$P = f(V)$ : hyperbole},
grid=both, grid style={popline!40},
]
\addplot[popcurve, domain=0.5:2.5, samples=100] {100/x};
\addplot[only marks, mark=+, mark size=3pt, poporange, thick] coordinates {
(0.5,200) (1.0,100) (1.5,66.7) (2.0,50) (2.5,40)
};
\end{axis}
\end{tikzpicture}
\hfill
\begin{tikzpicture}
\begin{axis}[
width=0.48\linewidth, height=6cm,
xlabel={$1/V$ (L$^{-1}$)}, ylabel={$P$ (kPa)},
xmin=0, xmax=1.15, ymin=0, ymax=115,
title={$P = f(1/V)$ : droite},
grid=both, grid style={popline!40},
]
\addplot[popcurve, domain=0:1.1, samples=50] {100*x};
\addplot[only marks, mark=+, mark size=3pt, poporange, thick] coordinates {
(0.4,40) (0.5,50) (0.667,66.7) (1.0,100)
};
\end{axis}
\end{tikzpicture}
\end{center}
\legende{Linéarisation : à gauche, $P$ en fonction de $V$ donne une hyperbole difficile à exploiter ; à droite, $P$ en fonction de $1/V$ donne une droite passant par l'origine, ce qui prouve que $P\,V = \text{constante}$ (loi de Boyle-Mariotte).}
\end{popfigure}

\begin{exemplebox}[Linéarisation numérique]
Pour un gaz à température constante, on mesure : $V = \SI{1,0}{L} \Rightarrow P = \SI{100}{kPa}$ ; $V = \SI{2,0}{L} \Rightarrow P = \SI{50}{kPa}$. On calcule $1/V$ : $\SI{1,0}{L^{-1}}$ et $\SI{0,50}{L^{-1}}$. Le coefficient directeur du graphe $P = f(1/V)$ vaut :
\[ a = \frac{100 - 50}{1{,}0 - 0{,}50} = \SI{100}{kPa \cdot L} \]
d'où la loi $P = \dfrac{100}{V}$, c'est-à-dire $P\,V = \SI{100}{kPa \cdot L}$ : le produit $P\,V$ est constant, conformément à la loi de Boyle-Mariotte.
\end{exemplebox}

\begin{methode}[Résumé : établir une loi entre deux grandeurs mesurées]
\begin{enumerate}
\item Tableau de mesures $(x_i, y_i)$ avec unités et incertitudes, une seule grandeur variable à la fois.
\item Graphique $y = f(x)$ : axes gradués et légendés, points avec barres d'incertitude.
\item Identification de la forme : droite par l'origine, droite affine ou courbe (auquel cas on linéarise par changement de variable).
\item Calcul du coefficient directeur sur la droite modèle, avec son unité, et interprétation physique.
\item Écriture de la loi $y = a\,x$ (ou $y = a\,x + b$) et vérification de l'homogénéité.
\item Validation de la loi sur une mesure nouvelle, non utilisée pour l'établir.
\end{enumerate}
\end{methode}

\begin{experience}[TP : Vérification de la loi d'Ohm sur un conducteur]
\textbf{Matériel.} Générateur de tension variable (0--12 V) ou pile + rhéostat, ampèremètre (calibre 0,5 A ou 1 A), voltmètre (calibre 10 V ou 15 V), fils de connexion, résistance inconnue (ex. fil de nichrome de 1 m), thermomètre.
\textbf{Protocole.}
\begin{enumerate}
\item Monter le circuit en série : générateur -- interrupteur -- ampèremètre -- résistance -- voltmètre aux bornes de la résistance.
\item Fixer la température (ne pas laisser le courant chauffer le fil : ouvrir l'interrupteur entre les mesures).
\item Faire varier la tension du générateur (ou le curseur du rhéostat) pour obtenir au moins 6 valeurs d'intensité $I$ régulièrement espacées de 0 au max.
\item Relever $U$ et $I$ pour chaque point, noter les incertitudes $\Delta U$, $\Delta I$ (classe de l'appareil ou demi-division).
\item Tracer $U = f(I)$, barres d'incertitude, droite modèle, coefficient directeur $R$.
\end{enumerate}
\textbf{Interprétation.} Si la droite passe par l'origine et traverse les barres d'incertitude, la loi d'Ohm est vérifiée \emph{dans les conditions de l'expérience} (température constante). La pente donne la résistance $R$ en ohms.
\end{experience}

\begin{savaistu}[Boyle, Mariotte et la naissance de la démarche expérimentale]
En 1662, le physicien irlandais Robert Boyle publie ses mesures de pression et de volume d'air enfermé dans un tube en U rempli de mercure : il constate que le produit $P\,V$ reste constant. En France, l'abbé Edme Mariotte retrouve indépendamment le même résultat en 1676 en précisant « à température constante » — un bel exemple de \textbf{paramètre fixé} ! Cette loi, première loi quantitative de l'histoire de la physique, illustre parfaitement la démarche que tu viens d'apprendre : mesures soignées, tableau, recherche d'une relation simple, énoncé d'une loi assortie de ses conditions de validité.
\end{savaistu}

\begin{exercice}[Entraînement : Linéarisation du pendule simple]
On étudie la période $T$ d'un pendule simple en fonction de sa longueur $l$. On obtient les mesures suivantes (incertitudes négligeables pour cet exercice) :

\begin{center}
\begin{tabular}{|c|c|c|c|c|c|}\hline
$l$ (m) & 0,20 & 0,40 & 0,60 & 0,80 & 1,00 \\ \hline
$T$ (s) & 0,90 & 1,27 & 1,55 & 1,79 & 2,00 \\ \hline
\end{tabular}
\end{center}

\begin{enumerate}
\item Tracer $T = f(l)$. Quelle est la forme de la courbe ?
\item En supposant une loi de type $T = k \sqrt{l}$, quel graphique tracer pour linéariser ?
\item Tracer ce graphique, déterminer $k$ et comparer à la valeur théorique $2\pi/\sqrt{g}$ (avec $g = \SI{9,81}{m.s^{-2}}$).
\end{enumerate}
\end{exercice}

\begin{corrige}[Exercice : Linéarisation du pendule simple]
\begin{enumerate}
\item Le graphique $T = f(l)$ est une courbe concave (racine carrée), ne passant pas par l'origine de manière linéaire.
\item Si $T = k \sqrt{l}$, alors $T$ est proportionnelle à $\sqrt{l}$. On trace $T = f(\sqrt{l})$.
\item Calcul des $\sqrt{l}$ : $0,447 ; 0,632 ; 0,775 ; 0,894 ; 1,000$ (en $\text{m}^{1/2}$).
Le graphique $T = f(\sqrt{l})$ est une droite passant par l'origine.
Coefficient directeur $k = \dfrac{2,00 - 0}{1,000 - 0} = \SI{2,00}{s.m^{-1/2}}$ (en prenant le point extrême de la droite).
Valeur théorique : $k_{\text{th}} = \dfrac{2\pi}{\sqrt{9,81}} \approx \SI{2,006}{s.m^{-1/2}}$.
L'expérience est en excellent accord avec la théorie.
\end{enumerate}
\end{corrige}

Cette démarche générale étant désormais claire, mettons-la en œuvre sur le premier grand exemple du programme : la loi d'Ohm.

\coursec{Exemple 1 : la loi d'Ohm}

Dans la section précédente, nous avons appris la démarche expérimentale qui permet d'établir une loi : on formule des hypothèses, on réalise des mesures, on trace une représentation graphique, puis on cherche une relation mathématique simple entre les grandeurs. Il est temps de mettre cette démarche en application sur un premier exemple historique et fondamental : la \cle{loi d'Ohm}. Cette loi, que tu utilises depuis la classe de Seconde sans toujours en connaître les limites, est un modèle : elle décrit parfaitement certains dipôles... et échoue pour d'autres. Comprendre \emph{pourquoi} et \emph{quand} elle s'applique, c'est exactement travailler comme un scientifique.

\courssub{Le dipôle ohmique : présentation et conventions}

\textbf{Une situation concrète.} Tu veux alimenter une petite diode électroluminescente (DEL) avec une batterie de \SI{9}{\volt}. Si tu branches la DEL directement, elle grille instantanément. Un électronicien te dira : « mets une résistance de protection en série ». Quelle valeur choisir ? Pour répondre, il faut comprendre ce qu'est un \cle{résistor} et quelle loi gouverne son comportement.

\begin{definition}[Dipôle ohmique]
Un \cle{dipôle ohmique} (appelé aussi \cle{résistor}) est un dipôle dont la \cle{caractéristique} $U = f(I)$ est une \textbf{droite passant par l'origine}. Autrement dit, la tension à ses bornes est proportionnelle à l'intensité du courant qui le traverse. Son symbole normalisé est un rectangle ; on note $R$ sa \cle{résistance}, exprimée en ohms ($\si{\ohm}$).
\end{definition}

\textbf{La convention récepteur.} Pour étudier un dipôle, il faut choisir une convention d'orientation des grandeurs. En \cle{convention récepteur}, la flèche représentant la tension $U$ aux bornes du dipôle et la flèche représentant le courant $I$ qui le traverse sont en \textbf{sens opposés}. C'est la convention naturelle pour tout dipôle qui \emph{reçoit} de l'énergie électrique (résistor, lampe, moteur). Pour un générateur, on utilise au contraire la convention générateur (flèches dans le même sens).

\begin{attention}[Convention récepteur : un réflexe à acquérir]
Si tu orientes $U$ et $I$ dans le même sens pour un résistor, la loi d'Ohm s'écrit $U = -R\,I$ : le signe moins est une pure conséquence du choix d'orientation, pas une nouvelle physique. Au Probatoire, précise toujours sur ton schéma le sens de $I$ et la flèche de $U$ avant d'écrire la moindre relation.
\end{attention}

\begin{popfigure}
\begin{center}
\begin{circuitikz}[european, scale=1.0, transform shape]
\draw (0,0) to[battery1, l_={$E$}] (0,3)
      to[ammeter, l={$A$}] (3,3)
      to[R, l={$R$}, v={$U$}, i>_=$I$] (6,3)
      -- (6,0) -- (0,0);
\draw (6,3) -- (8,3) to[voltmeter, l={$V$}] (8,0) -- (6,0);
\end{circuitikz}
\end{center}
\legende{Circuit de mesure : le générateur alimente le résistor $R$ ; l'ampèremètre $A$ (en série) mesure $I$, le voltmètre $V$ (en dérivation) mesure $U$. Les flèches $U$ et $I$ sont en sens opposés : convention récepteur.}
\end{popfigure}

\courssub{Énoncé de la loi d'Ohm}

Reprenons la démarche expérimentale de la section précédente. Avec le montage ci-dessus, on fait varier la tension délivrée par le générateur et on relève, pour chaque réglage, le couple de valeurs $(I, U)$. On obtient par exemple :

\begin{center}
\begin{tabularx}{\linewidth}{|l|X|X|X|X|X|X|}
\hline
\rowcolor{popblueL}
$I$ (\si{\milli\ampere}) & 0 & 21 & 43 & 64 & 85 & 106 \\
\hline
$U$ (\si{\volt}) & 0 & 1,0 & 2,0 & 3,0 & 4,0 & 5,0 \\
\hline
\end{tabularx}
\end{center}

Le rapport $\dfrac{U}{I}$ est constant (ici $\approx \SI{47}{\ohm}$) : les points sont alignés sur une droite passant par l'origine. C'est la signature d'une \textbf{proportionnalité}.

\begin{propriete}[Loi d'Ohm]
Pour un dipôle ohmique, la tension $U$ à ses bornes est proportionnelle à l'intensité $I$ du courant qui le traverse :
\[ U = R \times I \]
avec $U$ en volts (\si{\volt}), $I$ en ampères (\si{\ampere}) et $R$ en ohms (\si{\ohm}). Cette loi est \textbf{admise} : elle a été établie expérimentalement en activité ; aucune démonstration n'est exigible au Probatoire.
\end{propriete}

\textbf{Vérification par analyse dimensionnelle.} L'ohm est défini par cette loi : $1\ \si{\ohm} = 1\ \si{\volt} \cdot \si{\ampere}^{-1}$. Le produit $R \times I$ s'exprime donc en $\si{\ohm} \times \si{\ampere} = \si{\volt}$ : l'homogénéité est respectée. Prends l'habitude de vérifier systématiquement l'unité du résultat : une tension calculée en ampères trahit immédiatement une erreur de formule.

\begin{aretenir}[Les trois formes de la loi d'Ohm]
\[ U = R\,I \qquad\Longleftrightarrow\qquad I = \dfrac{U}{R} \qquad\Longleftrightarrow\qquad R = \dfrac{U}{I} \]
Connaissant deux des trois grandeurs, on calcule la troisième. La grandeur $G = \dfrac{1}{R}$, appelée \cle{conductance} (en siemens, \si{\siemens}), exprime au contraire la facilité avec laquelle le dipôle laisse passer le courant : $I = G\,U$.
\end{aretenir}

\courssub{Caractéristique $U = f(I)$ et interprétation physique de la résistance}

La \cle{caractéristique} d'un dipôle est la courbe $U = f(I)$. Pour un résistor, c'est une droite passant par l'origine, dont la \textbf{pente est la résistance} $R$ :
\[ R = \dfrac{\Delta U}{\Delta I} = \dfrac{U_B - U_A}{I_B - I_A} \]
Plus la pente est forte, plus le dipôle « résiste » au passage du courant.

\textbf{Interprétation microscopique (intuition).} Dans un conducteur, le courant est un déplacement ordonné de porteurs de charge (les électrons libres dans un métal). Ces électrons subissent des collisions avec les ions du réseau cristallin, qui freinent leur mouvement : c'est cette « friction » microscopique qui constitue la résistance. L'énergie perdue dans les collisions est transférée au réseau sous forme d'agitation thermique : c'est l'\cle{effet Joule}, $P = U\,I = R\,I^2$. Un résistor traversé par un courant chauffe donc toujours — ce point sera décisif pour comprendre les limites du modèle.

\begin{exemplebox}[Deux applications directes de la loi d'Ohm]
Un résistor de résistance $R = \SI{47}{\ohm}$ est traversé par un courant d'intensité $I = \SI{0,15}{\ampere}$. La tension à ses bornes vaut :
\[ U = R\,I = 47 \times 0,15 = \SI{7,05}{\volt} \approx \SI{7,1}{\volt} \]
Réciproquement, si on soumet ce même résistor à une tension $U = \SI{12}{\volt}$, l'intensité qui le traverse est :
\[ I = \dfrac{U}{R} = \dfrac{12}{47} \approx \SI{0,26}{\ampere} \]
Vérification d'homogénéité : $\si{\volt}/\si{\ohm} = \si{\ampere}$. Le résultat est cohérent.
\end{exemplebox}

\courssub{Les hypothèses simplificatrices du modèle}

Voici le cœur de cette leçon : la loi d'Ohm n'est pas une vérité absolue, c'est un \cle{modèle} construit sur des hypothèses simplificatrices qu'il faut savoir expliciter.

\begin{itemize}
\item \textbf{Hypothèse 1 : la résistance est constante.} On suppose que $R$ ne dépend ni de $U$, ni de $I$, ni du temps. Or $R$ dépend en réalité de la température : pour un métal, $R$ augmente quand la température augmente. Le modèle suppose donc implicitement que le résistor \textbf{ne chauffe pas sensiblement} pendant les mesures, c'est-à-dire que l'effet Joule reste modéré.
\item \textbf{Hypothèse 2 : le dipôle est symétrique.} La caractéristique est la même quel que soit le sens du courant : la droite obtenue pour $I > 0$ se prolonge pour $I < 0$ (caractéristique impaire). Ce n'est pas le cas de tous les dipôles : une diode laisse passer le courant dans un seul sens.
\item \textbf{Hypothèse 3 : le régime est permanent.} On néglige tout effet capacitif ou inductif : la relation $U = R\,I$ est supposée valable instantanément.
\end{itemize}

\begin{attention}[Erreur fréquente : appliquer $U = R\,I$ à n'importe quel dipôle]
La loi d'Ohm ne s'applique \textbf{que si la caractéristique est une droite passant par l'origine}. Avant d'écrire $U = R\,I$, vérifie cette condition (graphiquement ou par le rapport $U/I$ constant). L'appliquer à une diode, à une lampe à incandescence ou à un électrolyseur est une faute de physique, même si le calcul numérique « fonctionne ».
\end{attention}

\courssub{Domaine de validité : le contre-exemple de la lampe à filament}

\textbf{Expérience décisive.} Reprenons le même montage, mais remplaçons le résistor par une lampe à incandescence (à filament de tungstène). Les mesures donnent des points qui \emph{ne sont plus alignés} : la caractéristique se courbe, \textbf{sa pente augmente quand $I$ augmente}.

\textbf{Interprétation.} Quand l'intensité croît, l'effet Joule porte le filament à très haute température (environ \SI{2500}{\celsius} en fonctionnement normal : c'est ce qui le rend lumineux). Or la résistivité du tungstène augmente fortement avec la température : la résistance du filament à chaud peut être \textbf{dix fois plus grande} qu'à froid. L'hypothèse 1 ($R$ constante) est violée, et la loi d'Ohm cesse de s'appliquer.

\begin{popfigure}
\begin{center}
\begin{tikzpicture}
\begin{axis}[
    width=0.95\linewidth, height=7cm,
    xlabel={$I$ (\si{\milli\ampere})}, ylabel={$U$ (\si{\volt})},
    xmin=0, xmax=120, ymin=0, ymax=6,
    xtick={0,20,40,60,80,100,120},
    ytick={0,1,2,3,4,5,6},
    grid=major, grid style={popline!40},
    legend style={at={(0.03,0.97)}, anchor=north west, font=\small},
    clip=false
]
% Zone de validité de la loi d'Ohm (pour le résistor)
\fill[popgreenL, opacity=0.5] (axis cs:0,0) rectangle (axis cs:110,5.2);
\node[font=\small\itshape, text=popgreen!60!black, align=center] at (axis cs:78,4.6) {domaine de validité\\ de la loi d'Ohm};
% Résistor : droite U = 47 I (I en A)
\addplot[popcurve, popblue, very thick, domain=0:0.105, samples=50]{47*x};
\addlegendentry{résistor $R = \SI{47}{\ohm}$}
% Lampe à filament : courbe convexe (pente croissante)
\addplot[popcurve, poporange, very thick, domain=0:0.115, samples=100]{40*x + 1500*x^3};
\addlegendentry{lampe à filament}
% Point de mesure exemple
\addplot[only marks, mark=*, popdark] coordinates {(64,3.0)};
\end{axis}
\end{tikzpicture}
\end{center}
\legende{Caractéristiques $U = f(I)$. Le résistor : droite passant par l'origine, de pente $R$ (loi d'Ohm applicable). La lampe à filament : courbe convexe dont la pente augmente avec $I$ (le filament chauffe, $R$ augmente) : la loi d'Ohm n'est plus valable.}
\end{popfigure}

\begin{aretenir}[Domaine de validité de la loi d'Ohm]
La loi d'Ohm $U = R\,I$ s'applique :
\begin{itemize}
\item aux \textbf{dipôles ohmiques} (résistors, fils conducteurs métalliques) ;
\item tant que leur \textbf{température reste pratiquement constante} (courants modérés).
\end{itemize}
Elle ne s'applique pas aux dipôles \textbf{non ohmiques} : lampes à incandescence (caractéristique courbe, pente croissante), diodes (caractéristique non symétrique et non linéaire), électrolyseurs, moteurs (la tension ne passe pas par l'origine : $U = E' + r\,I$).
\end{aretenir}

\courssub{Application au dimensionnement : résistance de protection}

Revenons à notre problème de départ : protéger une DEL alimentée par une batterie. C'est une application typique du savoir-faire « établir et exploiter une loi pour dimensionner un circuit ».

\begin{exoresolu}[Choix d'une résistance de protection pour une DEL]
\textbf{Énoncé.} Une DEL rouge fonctionne nominalement sous une tension $U_{\mathrm{DEL}} = \SI{2,0}{\volt}$ avec un courant $I = \SI{20}{\milli\ampere}$. On l'alimente avec une batterie de fem $E = \SI{9,0}{\volt}$, en plaçant un résistor $R$ en série avec la DEL. Déterminer la valeur de $R$ à choisir, puis vérifier que la puissance dissipée par le résistor reste acceptable (le résistor courant supporte au maximum $P_{\max} = \SI{0,25}{\watt}$).
\tcblower
\textbf{Solution détaillée.}

\textbf{Étape 1 : schéma et loi des mailles.} Le générateur, le résistor et la DEL sont en série : le même courant $I$ les traverse. La loi d'additivité des tensions s'écrit :
\[ E = U_R + U_{\mathrm{DEL}} \qquad\Longrightarrow\qquad U_R = E - U_{\mathrm{DEL}} = 9,0 - 2,0 = \SI{7,0}{\volt} \]

\textbf{Étape 2 : application de la loi d'Ohm au résistor.} Le résistor est un dipôle ohmique, parcouru par $I = \SI{20}{\milli\ampere} = \SI{0,020}{\ampere}$ :
\[ R = \dfrac{U_R}{I} = \dfrac{7,0}{0,020} = \SI{350}{\ohm} \]
On choisira la valeur normalisée la plus proche, par exemple $R = \SI{330}{\ohm}$ ou $\SI{390}{\ohm}$.

\textbf{Étape 3 : vérification thermique (test de validité du modèle !).} Pour que la loi d'Ohm reste valable, le résistor ne doit pas trop chauffer. La puissance dissipée par effet Joule vaut :
\[ P = R\,I^2 = 350 \times (0,020)^2 = 350 \times 4,0 \times 10^{-4} = \SI{0,14}{\watt} \]
Or $P = \SI{0,14}{\watt} < P_{\max} = \SI{0,25}{\watt}$ : le résistor chauffe modérément, sa résistance reste pratiquement constante, et notre modèle (loi d'Ohm avec $R$ fixe) est \textbf{cohérent}. Si le calcul avait donné $P > P_{\max}$, il aurait fallu choisir un résistor de puissance supérieure — ou revoir les hypothèses, car $R$ aurait varié avec la température.

\textbf{Conclusion.} $R \approx \SI{350}{\ohm}$ (valeur normalisée : $\SI{330}{\ohm}$ ou $\SI{390}{\ohm}$), puissance dissipée $\SI{0,14}{\watt}$, compatible avec un résistor standard $\SI{0,25}{\watt}$.
\end{exoresolu}

\begin{methode}[Exploiter la loi d'Ohm dans un circuit]
\begin{enumerate}
\item Faire le schéma du circuit et y placer les flèches $U$ et $I$ en convention récepteur pour chaque dipôle passif.
\item Vérifier que le dipôle est bien \textbf{ohmique} (caractéristique droite passant par l'origine, ou mention explicite dans l'énoncé).
\item Écrire les lois du circuit (additivité des tensions, unicité du courant en série) pour exprimer $U$ et $I$ aux bornes du résistor.
\item Appliquer $U = R\,I$ sous la forme adaptée, en convertissant toutes les grandeurs en unités SI (mA $\to$ A).
\item Vérifier la cohérence : homogénéité, ordre de grandeur, et si besoin puissance dissipée (validité de l'hypothèse « $R$ constante »).
\end{enumerate}
\end{methode}

\begin{savaistu}[Georg Ohm : une loi d'abord rejetée]
Le physicien allemand \textbf{Georg Simon Ohm} (1789--1854) publia sa loi en 1827 dans un ouvrage intitulé \emph{Die galvanische Kette, mathematisch bearbeitet} (« La chaîne galvanique traitée mathématiquement »). Loin d'être célébrée, sa théorie fut d'abord \textbf{rejetée} par la communauté scientifique allemande, qui lui reprochait son approche trop mathématique ; Ohm, alors simple professeur de collège, dut même démissionner de son poste. Il fallut attendre une quinzaine d'années pour que ses travaux soient reconnus : la Royal Society de Londres lui décerne la prestigieuse \textbf{médaille Copley} en 1841, et l'unité de résistance porte son nom depuis 1881. Une leçon d'histoire des sciences : une loi n'est acceptée que lorsque l'expérience, répétée et vérifiable, la confirme — et lorsque son domaine de validité est clairement établi. Aujourd'hui, chaque technicien d'ENEO qui dimensionne une ligne, chaque réparateur de téléphones à Douala qui teste une carte mère, applique sans y penser la loi de Georg Ohm.
\end{savaistu}

\begin{exercice}[Caractéristique expérimentale d'un dipôle inconnu]
On relève pour un dipôle inconnu les mesures suivantes : $(I\ \text{en mA} ; U\ \text{en V})$ : $(0;0)$, $(10;0{,}6)$, $(20;1{,}4)$, $(30;2{,}6)$, $(40;4{,}2)$.
\begin{enumerate}
\item Ce dipôle est-il ohmique ? Justifier en calculant le rapport $U/I$.
\item Tracer l'allure de sa caractéristique et proposer une nature possible pour ce dipôle.
\end{enumerate}
\end{exercice}

\begin{corrige}[Exercice : caractéristique d'un dipôle inconnu]
\begin{enumerate}
\item Calculons les rapports : $\dfrac{0,6}{0,010} = \SI{60}{\ohm}$ ; $\dfrac{1,4}{0,020} = \SI{70}{\ohm}$ ; $\dfrac{2,6}{0,030} \approx \SI{87}{\ohm}$ ; $\dfrac{4,2}{0,040} = \SI{105}{\ohm}$. Le rapport $U/I$ n'est \textbf{pas constant} : il augmente avec $I$. Le dipôle n'est donc \textbf{pas ohmique} et la loi d'Ohm ne s'y applique pas.
\item La caractéristique est une courbe passant par l'origine dont la pente croît avec $I$ : le dipôle chauffe et sa résistance augmente avec la température. Il s'agit typiquement d'une \textbf{lampe à incandescence} (filament de tungstène).
\end{enumerate}
\end{corrige}

En résumé, la loi d'Ohm illustre parfaitement ce qu'est un modèle en physique : une relation mathématique simple, $U = R\,I$, obtenue expérimentalement, extrêmement utile pour dimensionner les circuits, mais construite sur des hypothèses ($R$ constante, dipôle symétrique) qui définissent son \textbf{domaine de validité}. Le scientifique — et l'élève de Première C — doit toujours se poser la question : « les hypothèses du modèle sont-elles satisfaites ici ? ». Nous retrouverons exactement cette logique dans la section suivante avec la loi des gaz parfaits.

\coursec{Exemple 2 : la loi des gaz parfaits}

Dans la section précédente, nous avons établi la loi d'Ohm $U = R\,I$ à partir de mesures : une loi simple, valable dans un domaine précis (conducteur ohmique, température constante). Nous allons maintenant appliquer exactement la même démarche scientifique à un système d'apparence très différente : un \cle{gaz} enfermé dans une enceinte. Là encore, des mesures soigneuses vont conduire à une loi mathématique simple, la \cle{loi des gaz parfaits}, dont nous préciserons les hypothèses et le domaine de validité. Cette loi gouverne des situations que tu connais bien : le pneu d'une mototaxi qui chauffe sur la route de Yaoundé à Douala, la bouteille de gaz domestique exposée au soleil, le ballon de football qu'on regonfle avant le match.

\courssub{Les grandeurs d'état d'un gaz}

\begin{experience}[Observer pour modéliser : un gaz dans un piston]
Enfermons de l'air dans une seringue dont on bouche l'extrémité. Si l'on pousse le piston, le volume diminue et il faut forcer de plus en plus : la pression augmente. Si l'on chauffe la seringue (avec précaution) à volume bloqué, le piston pousse plus fort : la pression augmente encore. Si l'on injecte davantage d'air, la pression augmente aussi. Trois grandeurs semblent donc piloter l'état du gaz : le \textbf{volume}, la \textbf{pression} et la \textbf{température}, auxquelles s'ajoute la \textbf{quantité de matière} (le nombre de particules contenues dans l'enceinte). Ce sont les \cle{grandeurs d'état} du gaz.
\end{experience}

\begin{definition}[Grandeurs d'état d'un gaz]
L'état macroscopique d'un gaz est décrit par quatre grandeurs :
\begin{itemize}
\item la \cle{pression} $P$, en pascals (Pa) : elle traduit les chocs des particules sur les parois ; $\SI{1}{Pa} = \SI{1}{N.m^{-2}}$ ;
\item le \cle{volume} $V$ occupé par le gaz, en mètres cubes ($\text{m}^3$) ; rappel : $\SI{1}{L} = \SI{e-3}{m^3}$ ;
\item la \cle{température absolue} $T$, en kelvins (K) : elle mesure l'agitation thermique des particules ;
\item la \cle{quantité de matière} $n$, en moles (mol) : une mole contient $N_A = \num{6,02e23}$ particules.
\end{itemize}
\end{definition}

\begin{attention}[La température doit TOUJOURS être en kelvins]
L'erreur la plus fréquente — et la plus sanctionnée au Probatoire — est d'utiliser la température en degrés Celsius dans la loi des gaz parfaits. La conversion est \textbf{indispensable} :
\[ T(\text{K}) = \theta(^\circ\text{C}) + 273,15 \]
En exercice, on utilise la valeur approchée : $T(\text{K}) = \theta(^\circ\text{C}) + 273$. Ainsi $27\,^\circ\text{C}$ correspond à $T = \SI{300}{K}$. Le kelvin est une échelle \emph{absolue} : $\SI{0}{K}$ (zéro absolu) correspond à l'absence totale d'agitation thermique, et il n'existe pas de température négative en kelvins. C'est précisément pour cela que les lois des gaz, qui font intervenir des rapports de températures, exigent cette échelle.
\end{attention}

\begin{popfigure}
\begin{center}
\begin{tikzpicture}[scale=1.0]
% Enceinte
\draw[fill=popblueL, draw=popdark, thick] (0,0) rectangle (5,3);
% Piston
\draw[fill=popmuted!40, draw=popdark, thick] (5,-0.15) rectangle (5.5,3.15);
\draw[popdark, thick] (5.5,1.5) -- (7.2,1.5);
\draw[->, popink, very thick] (7.6,1.5) -- (6.6,1.5) node[midway, above]{\small force};
% Particules et vitesses
\foreach \p/\a in {{0.7,0.6}/25, {1.5,2.2}/-40, {2.3,1.1}/60, {3.1,2.5}/10, {3.9,0.8}/-70, {4.4,2.0}/45, {1.0,1.7}/-15, {2.8,0.4}/80, {3.6,1.6}/-30, {0.9,2.6}/55, {2.0,1.9}/-55, {4.2,1.2}/20}
{
  \fill[poporange] (\p) circle (0.06);
  \draw[->, poporange!80!black, thick] (\p) -- ++(\a:0.45);
}
% Annotations
\node[popdark] at (2.5,-0.45) {\small gaz : $n$ mol de particules en mouvement désordonné};
\node[popblue, font=\bfseries] at (0.45,2.6) {$P$};
\node[popblue, font=\bfseries] at (4.6,0.35) {$T$};
\draw[<->, poppurple, thick] (0,-0.15) -- (5,-0.15) node[midway, below, yshift=-14pt]{\small volume $V$};
\end{tikzpicture}
\end{center}
\legende{Un gaz enfermé dans une enceinte fermée par un piston mobile. Les particules, en mouvement désordonné (agitation thermique), frappent les parois : c'est l'origine microscopique de la pression $P$.}
\end{popfigure}

\courssub{Le modèle du gaz parfait : des hypothèses simplificatrices}

Comme pour la loi d'Ohm, le physicien commence par \emph{simplifier la réalité} pour la rendre calculable. Un gaz réel contient environ $\num{e25}$ molécules par mètre cube : suivre chacune d'elles est impossible. On remplace donc le gaz réel par un gaz idéalisé, le \cle{gaz parfait}.

\begin{definition}[Hypothèses du modèle du gaz parfait]
Le modèle du gaz parfait repose sur trois hypothèses simplificatrices :
\begin{enumerate}
\item les particules sont \textbf{ponctuelles} : leur volume propre est négligeable devant le volume $V$ de l'enceinte ;
\item il n'y a \textbf{aucune interaction} entre les particules, en dehors des chocs (pas de forces d'attraction à distance) ;
\item les chocs entre particules et avec les parois sont \textbf{élastiques} : l'énergie cinétique totale se conserve.
\end{enumerate}
Entre deux chocs, chaque particule est donc un système pseudo-isolé : d'après le principe d'inertie, elle se déplace en ligne droite à vitesse constante.
\end{definition}

Ces hypothèses sont d'autant mieux vérifiées que le gaz est \textbf{dilué} (faible pression, donc particules éloignées les unes des autres) et \textbf{chaud} (agitation thermique intense, qui domine les faibles attractions résiduelles). Nous y reviendrons au moment de définir le domaine de validité.

\courssub{Énoncé de la loi des gaz parfaits}

Les mesures expérimentales (pression, volume, température, quantité de matière) effectuées sur des gaz dilués montrent que le rapport $\dfrac{P\,V}{n\,T}$ est le même pour tous les gaz : c'est une constante universelle, notée $R$.

\begin{propriete}[Loi des gaz parfaits (admise)]
Pour tout gaz dont le comportement est correctement décrit par le modèle du gaz parfait, les quatre grandeurs d'état sont reliées par :
\[ \boxed{\ P\,V = n\,R\,T\ } \]
où $P$ est en pascals (Pa), $V$ en mètres cubes ($\text{m}^3$), $n$ en moles (mol), $T$ en kelvins (K), et $R$ la \cle{constante des gaz parfaits} :
\[ R = \SI{8,314}{\joule\per\kelvin\per\mole} \]
Cette loi est \textbf{admise} en Première : sa démonstration n'est pas exigible au Probatoire, mais tu dois savoir l'appliquer et en déduire les lois à paramètre constant.
\end{propriete}

Vérifions l'homogénéité, comme tout bon scientifique : $[nRT] = \text{mol}\times\text{J}\cdot\text{K}^{-1}\cdot\text{mol}^{-1}\times\text{K} = \text{J}$, et $[PV] = \text{Pa}\times\text{m}^3 = \text{N}\cdot\text{m}^{-2}\times\text{m}^3 = \text{N}\cdot\text{m} = \text{J}$. Les deux membres sont bien des énergies : la loi est dimensionnellement cohérente.

\begin{exemplebox}[Pression de l'air dans un pneu de mototaxi]
Un pneu contient $n = \SI{1,2}{mol}$ d'air dans un volume $V = \SI{30}{L} = \SI{0,030}{m^3}$, à la température $\theta = \SI{27}{\celsius}$, soit $T = \SI{300}{K}$. La loi des gaz parfaits donne :
\begin{align*}
P &= \frac{n\,R\,T}{V} = \frac{1,2 \times 8,314 \times 300}{0,030} \\
  &\approx \SI{1,0e5}{Pa}
\end{align*}
soit environ la pression atmosphérique. Après un long trajet sous le soleil, la température atteint $\theta' = \SI{57}{\celsius}$, soit $T' = \SI{330}{K}$. À volume quasi constant :
\[ P' = \frac{1,2 \times 8,314 \times 330}{0,030} \approx \SI{1,1e5}{Pa} \]
La pression a augmenté d'environ $10\ \%$ : c'est pourquoi les vulcanisateurs de Douala recommandent de vérifier la pression des pneus \emph{à froid}, et pourquoi une bouteille de gaz domestique ne doit jamais rester exposée en plein soleil.
\end{exemplebox}

\courssub{Les lois déduites à paramètre constant}

Toute la puissance de la loi $PV = nRT$ est de contenir, comme cas particuliers, les lois historiques découvertes expérimentalement aux XVII\textsuperscript{e} et XVIII\textsuperscript{e} siècles. La méthode est toujours la même : on bloque deux grandeurs, et on étudie la relation entre les deux autres.

\begin{experience}[Démonstration guidée : la loi de Boyle-Mariotte]
Considérons une quantité $n$ de gaz maintenue à température $T$ constante (enceinte thermostatée). Rédigeons le raisonnement étape par étape :
\begin{enumerate}
\item Écris la loi des gaz parfaits : $P\,V = n\,R\,T$.
\item Justifie que le membre de droite est constant : $n$ est fixé (enceinte fermée), $T$ est fixé (thermostat), $R$ est une constante universelle. Donc $nRT = \text{constante}$.
\item Conclus que $P\,V = \text{constante}$.
\item Applique ce résultat entre deux états (1) et (2) du même gaz : $P_1 V_1 = nRT$ et $P_2 V_2 = nRT$, d'où :
\[ P_1\,V_1 = P_2\,V_2 \]
\end{enumerate}
C'est la \cle{loi de Boyle-Mariotte} : à température et quantité de matière constantes, le produit $P\,V$ se conserve. Si l'on comprime le gaz d'un facteur 2, sa pression double.
\end{experience}

\begin{propriete}[Les trois lois à paramètre constant]
À partir de $PV = nRT$, avec $n$ constant (enceinte fermée) :
\begin{itemize}
\item \textbf{Loi de Boyle-Mariotte} ($T$ constante) : $P\,V = \text{constante}$, soit $P_1V_1 = P_2V_2$ ;
\item \textbf{Loi de Charles} ($P$ constante) : $\dfrac{V}{T} = \text{constante}$, soit $\dfrac{V_1}{T_1} = \dfrac{V_2}{T_2}$ ;
\item \textbf{Loi de Gay-Lussac} ($V$ constant) : $\dfrac{P}{T} = \text{constante}$, soit $\dfrac{P_1}{T_1} = \dfrac{P_2}{T_2}$.
\end{itemize}
Dans toutes ces relations, $T$ est impérativement en kelvins.
\end{propriete}

\begin{popfigure}
\begin{center}
\begin{tikzpicture}
\begin{axis}[
  width=0.85\linewidth, height=6.5cm,
  xlabel={$V$ (m$^3$)}, ylabel={$P$ (Pa)},
  xmin=0, xmax=0.06, ymin=0, ymax=250000,
  xtick={0.01,0.02,0.03,0.04,0.05},
  scaled y ticks=false,
  yticklabel style={/pgf/number format/fixed},
  grid=major, grid style={popline!50},
  legend style={at={(0.97,0.97)}, anchor=north east, font=\small}
]
% T1 = 280 K : P = nRT/V avec n=1.2 -> 2793.9/V
\addplot[popcurve, domain=0.015:0.06, samples=100]{2793.9/x};
\addlegendentry{$T_1 = 280\ \text{K}$}
% T2 = 330 K
\addplot[poporange, very thick, domain=0.015:0.06, samples=100]{3292.9/x};
\addlegendentry{$T_2 = 330\ \text{K}$}
% Zone basse pression
\draw[popgreen, dashed, thick] (axis cs:0.015,100000) -- (axis cs:0.06,100000);
\node[popgreen!60!black, font=\small, anchor=south west] at (axis cs:0.016,102000) {modèle valide à basse pression};
\end{axis}
\end{tikzpicture}
\end{center}
\legende{Isothermes $P = \dfrac{nRT}{V}$ pour $n = \SI{1,2}{mol}$ à deux températures $T_1 < T_2$ : ce sont des hyperboles (loi de Boyle-Mariotte). À volume donné, la pression est d'autant plus grande que la température est élevée.}
\end{popfigure}

\courssub{Domaine de validité : quand le gaz parfait devient un gaz réel}

Comme la loi d'Ohm cesse de s'appliquer quand le conducteur chauffe, la loi des gaz parfaits a ses \cle{contraintes} : elle n'est valable que si ses hypothèses restent raisonnablement vérifiées.

\begin{aretenir}[Domaine de validité du modèle du gaz parfait]
Le modèle du gaz parfait décrit correctement un gaz réel :
\begin{itemize}
\item à \textbf{faible pression} (typiquement $P$ de l'ordre de la pression atmosphérique ou inférieure) : les particules sont loin les unes des autres, leur volume propre et leurs interactions deviennent négligeables ;
\item à \textbf{température modérée ou élevée}, loin de la température de liquéfaction du gaz.
\end{itemize}
Le modèle \textbf{échoue} à très haute pression (le volume propre des molécules n'est plus négligeable devant $V$) et près de la liquéfaction (les attractions entre molécules deviennent importantes) : on parle alors de \cle{gaz réel}, décrit par des lois plus complexes.
\end{aretenir}

\begin{attention}[Pièges fréquents dans l'application de $PV = nRT$]
\begin{itemize}
\item \textbf{Température en Celsius} : convertis toujours en kelvins avant tout calcul ($T = \theta + 273$).
\item \textbf{Volume en litres} : convertis en mètres cubes ($\SI{1}{L} = \SI{e-3}{m^3}$). Avec $R = \SI{8,314}{J.K^{-1}.mol^{-1}}$, toutes les grandeurs doivent être en unités SI.
\item \textbf{Pression en bars} : $\SI{1}{bar} = \SI{e5}{Pa}$ ; pense à convertir.
\item \textbf{Oublier que $n$ doit être constant} dans les lois de Boyle-Mariotte, Charles et Gay-Lussac : si l'enceinte fuit, ces lois ne s'appliquent plus directement.
\end{itemize}
\end{attention}

\begin{exoresolu}[Bouteille de gaz domestique exposée au soleil]
Une bouteille de butane contient $n = \SI{40}{mol}$ de gaz dans un volume utile $V = \SI{50}{L}$. Le matin, la température est $\theta_1 = \SI{22}{\celsius}$ et la pression $P_1 = \SI{2,0e5}{Pa}$. L'après-midi, la bouteille oubliée au soleil atteint $\theta_2 = \SI{52}{\celsius}$. On suppose le volume constant et le gaz parfait. Calculer la pression $P_2$ atteinte et commenter.
\tcblower
\textbf{Solution détaillée.} Le volume et la quantité de matière sont constants : on applique la loi de Gay-Lussac, après conversion des températures en kelvins :
\[ T_1 = 22 + 273 = \SI{295}{K} \qquad T_2 = 52 + 273 = \SI{325}{K} \]
La loi de Gay-Lussac s'écrit $\dfrac{P_1}{T_1} = \dfrac{P_2}{T_2}$, d'où :
\begin{align*}
P_2 &= P_1 \times \frac{T_2}{T_1} = \SI{2,0e5}{Pa} \times \frac{325}{295} \\
    &\approx \SI{2,2e5}{Pa}
\end{align*}
\textbf{Commentaire.} La pression augmente d'environ $10\ \%$ pour une élévation de température de $30\,^\circ\text{C}$. Cette surpression, répétée jour après jour, fatigue les soudures et les robinets des bouteilles : c'est la raison physique des consignes de sécurité (stockage à l'ombre, loin des sources de chaleur) affichées chez les revendeurs de gaz à Yaoundé comme à Garoua. Remarquons aussi que si l'on avait utilisé les degrés Celsius, on aurait trouvé $P_2 = \SI{2,0e5}{Pa} \times \dfrac{52}{22} \approx \SI{4,7e5}{Pa}$ : un résultat physiquement faux, d'où l'importance absolue de la conversion en kelvins.
\end{exoresolu}

\begin{savaistu}[L'air est (presque) un gaz parfait]
Dans les conditions ambiantes ($P \approx \SI{e5}{Pa}$, $\theta \approx \SI{25}{\celsius}$), l'air que nous respirons vérifie la loi des gaz parfaits à mieux que $1\ \%$ près : ses molécules (diazote, dioxygène) sont petites, très éloignées les unes des autres et n'interagissent presque pas. C'est pourquoi la loi $PV = nRT$ s'applique si bien aux pneus, aux ballons de football et aux poumons ! En revanche, dans les bouteilles de gaz comprimé à très haute pression ou dans les systèmes de réfrigération où le fluide frôle la liquéfaction, les ingénieurs doivent utiliser des modèles de gaz réels. La loi des gaz parfaits a été établie progressivement : Boyle et Mariotte au XVII\textsuperscript{e} siècle pour le produit $PV$, Charles et Gay-Lussac vers 1800 pour le rôle de la température, avant la synthèse définitive au XIX\textsuperscript{e} siècle.
\end{savaistu}

\begin{methode}[Résoudre un problème de gaz en quatre étapes]
\begin{enumerate}
\item \textbf{Identifier} les grandeurs connues et inconnues ($P$, $V$, $n$, $T$) et préciser ce qui reste constant.
\item \textbf{Convertir} toutes les données en unités SI : $T$ en kelvins, $V$ en $\text{m}^3$, $P$ en pascals.
\item \textbf{Choisir} la bonne relation : $PV = nRT$ en général, ou la loi à paramètre constant adaptée (Boyle-Mariotte, Charles, Gay-Lussac).
\item \textbf{Calculer}, vérifier l'ordre de grandeur et l'homogénéité du résultat, puis commenter en lien avec la situation concrète.
\end{enumerate}
\end{methode}

\begin{exercice}[Le ballon du match]
Un ballon de football de volume $V = \SI{5,6}{L}$ est gonflé avec $n = \SI{0,50}{mol}$ d'air à $\theta = \SI{30}{\celsius}$. 
\begin{enumerate}
\item Calculer la pression de l'air dans le ballon.
\item Le ballon est laissé la nuit à $\SI{18}{\celsius}$ : que devient la pression, le volume étant supposé constant ?
\item Justifier que le modèle du gaz parfait est applicable ici.
\end{enumerate}
\end{exercice}

\begin{corrige}[Exercice : le ballon du match]
\begin{enumerate}
\item Conversion : $V = \SI{5,6e-3}{m^3}$, $T = 30 + 273 = \SI{303}{K}$. La loi des gaz parfaits donne :
\[ P = \frac{nRT}{V} = \frac{0,50 \times 8,314 \times 303}{5,6\times10^{-3}} \approx \SI{2,25e5}{Pa} \]
On arrondit à $\SI{2,2e5}{Pa}$ (deux chiffres significatifs).
\item À volume constant, loi de Gay-Lussac avec $T' = 18 + 273 = \SI{291}{K}$ :
\[ P' = P\,\frac{T'}{T} = \SI{2,25e5}{Pa}\times\frac{291}{303} \approx \SI{2,16e5}{Pa} \]
La pression diminue légèrement : le ballon paraît un peu « mou » le matin.
\item La pression est de l'ordre de $\SI{2e5}{Pa}$ (faible) et la température est ordinaire, très loin de la liquéfaction de l'air ($\approx \SI{-190}{\celsius}$) : les hypothèses du gaz parfait sont excellentes.
\end{enumerate}
\end{corrige}

Nous disposons maintenant de deux exemples complets de lois physiques : la loi d'Ohm et la loi des gaz parfaits. Dans les deux cas, une loi simple a été obtenue moyennant des hypothèses simplificatrices, et elle possède un domaine de validité qu'il faut connaître. La section suivante généralise cette réflexion : qu'est-ce qui fait la force — et la fragilité — d'une loi physique ?

\coursec{Contraintes d'une loi : domaine de validité et limites des modèles}

Dans la section précédente, nous avons établi la loi des gaz parfaits $PV = nRT$ à partir de mesures réalisées sur de l'air à pression modérée. Mais que se passerait-il si l'on comprimait ce gaz jusqu'à \SI{200}{bar}, comme dans les bouteilles de plongée ou les bonbonnes de gaz industrielles que l'on voit chez les soudeurs de Douala ? La loi continuerait-elle à décrire fidèlement la réalité ? De même, la loi d'Ohm $U = RI$, vérifiée avec soin sur un résistor, s'applique-t-elle au filament d'une lampe à incandescence porté à plus de \SI{2500}{\celsius} ?

Ces questions ne sont pas des détails : elles touchent au c\oe ur du travail scientifique. \textbf{Toute loi physique est un modèle}, c'est-à-dire une représentation simplifiée de la réalité, construite sur des hypothèses. Or une simplification qui est acceptable dans certaines conditions devient intenable dans d'autres. Cette section te donne les outils pour répondre à la question fondamentale : \emph{jusqu'où peut-on faire confiance à une loi ?}

\courssub{Le domaine de validité d'une loi}

\textbf{Intuition.} Imagine la carte routière d'une ville : elle est parfaite pour se déplacer à Yaoundé, mais inutilisable pour naviguer en pirogue sur le fleuve Wouri. La carte n'est pas « fausse » : elle est simplement conçue pour un certain usage, dans certaines conditions. Une loi physique fonctionne exactement pareil.

\begin{definition}[Domaine de validité]
Le \cle{domaine de validité} d'une loi est l'ensemble des conditions expérimentales (valeurs des grandeurs mises en jeu, hypothèses vérifiées) pour lesquelles les prévisions de la loi sont en accord avec les mesures, compte tenu des incertitudes expérimentales.
\end{definition}

Ce domaine est délimité par deux types de frontières :

\begin{itemize}
\item des \textbf{frontières numériques} : intervalles de valeurs des grandeurs (pression, température, intensité, vitesse...) à l'intérieur desquels la loi a été testée avec succès ;
\item des \textbf{frontières conceptuelles} : les hypothèses simplificatrices sur lesquelles repose le modèle (frottements négligés, résistance constante, molécules sans volume propre...).
\end{itemize}

\begin{aretenir}[Chaque hypothèse est une limite]
Chaque hypothèse simplificatrice formulée lors de la construction d'un modèle définit une limite de ce modèle. Dès qu'une hypothèse cesse d'être vérifiée expérimentalement, la loi risque de sortir de son domaine de validité. Connaître une loi, c'est donc aussi connaître ses hypothèses.
\end{aretenir}

Illustrons ce principe sur nos deux exemples de référence :

\begin{center}
\begin{tabularx}{\linewidth}{|l|X|X|}
\hline
\rowcolor{popblueL} \textbf{Loi} & \textbf{Hypothèses simplificatrices} & \textbf{Limite correspondante} \\
\hline
Loi d'Ohm $U = RI$ & La résistance $R$ est constante (température du conducteur fixe) & Invalide si le conducteur chauffe fortement (lampe à incandescence) \\
\hline
Gaz parfaits $PV = nRT$ & Molécules ponctuelles, sans interaction entre elles & Invalide à haute pression ou basse température (gaz réel) \\
\hline
Chute libre $v = gt$ & Frottements de l'air négligés & Invalide pour une feuille de papier ou un parachutiste \\
\hline
\end{tabularx}
\end{center}

\courssub{Tester la validité d'un modèle par la mesure}

Comment savoir, en pratique, si l'on se trouve encore dans le domaine de validité ? Il n'existe qu'un seul juge : \textbf{l'expérience}. La démarche consiste à comparer quantitativement les prévisions du modèle à des mesures réalisées dans des conditions variées.

\begin{definition}[Écart relatif]
L'\cle{écart relatif} entre une mesure et la prévision d'un modèle est le rapport :
\[
\varepsilon = \frac{|\text{mesure} - \text{modèle}|}{\text{mesure}}
\]
Il s'exprime sans unité, le plus souvent en pourcentage. Vérifions l'homogénéité : le numérateur et le dénominateur ont la même unité, donc $\varepsilon$ est bien un nombre sans dimension.
\end{definition}

\begin{methode}[Tester la validité d'une hypothèse par la mesure]
\begin{enumerate}
\item \textbf{Énoncer la prévision du modèle} : calculer la valeur de la grandeur prédite par la loi dans les conditions de l'expérience.
\item \textbf{Mesurer la grandeur correspondante} dans ces mêmes conditions, en estimant l'incertitude de mesure.
\item \textbf{Calculer l'écart relatif} $\varepsilon$ entre la mesure et la prévision.
\item \textbf{Conclure} : si $\varepsilon$ est inférieur au seuil fixé (par exemple 5 \%) ou compatible avec les incertitudes, le modèle est validé dans ces conditions ; sinon, on est sorti du domaine de validité et il faut réviser les hypothèses.
\end{enumerate}
\end{methode}

\begin{propriete}[Critère quantitatif de validité]
Un modèle est jugé \cle{valide} dans des conditions données si l'écart relatif entre ses prévisions et les mesures reste inférieur à un seuil fixé à l'avance (par exemple 5 \%), ou si l'écart est compatible avec les incertitudes de mesure. Le choix du seuil dépend de la précision exigée par la situation : 1 \% pour un dosage médical, 10 \% peuvent suffire pour une estimation artisanale.
\end{propriete}

\begin{exoresolu}[Sortie du domaine de validité de la loi d'Ohm]
Un élève de Première C a établi, à partir de mesures à faible courant, le modèle $U = 20\,I$ pour un résistor métallique (avec $U$ en volts et $I$ en ampères). Il alimente ensuite ce résistor avec un courant $I = \SI{10}{A}$.

1. Quelle tension le modèle prévoit-il ?
2. La mesure réelle donne $U_{\text{mes}} = \SI{215}{V}$. Calculer l'écart relatif.
3. Le seuil de validité est fixé à 5 \%. Conclure et proposer une explication physique.
\tcblower
\textbf{Solution détaillée.}

1. Prévision du modèle :
\[
U_{\text{mod}} = 20 \times 10 = \SI{200}{V}.
\]

2. Écart relatif :
\[
\varepsilon = \frac{|U_{\text{mes}} - U_{\text{mod}}|}{U_{\text{mes}}} = \frac{|215 - 200|}{215} = \frac{15}{215} \approx \num{0,0698} \approx 7,0\,\%.
\]

3. Comparaison au seuil : $\varepsilon \approx 7,0\,\% > 5\,\%$. \textbf{Le modèle $U = 20\,I$ n'est plus valide} dans ces conditions.

\textbf{Interprétation physique} : un courant de \SI{10}{A} dégage par effet Joule une puissance $P = U_{\text{mes}} I \approx \SI{2,15}{kW}$, ce qui échauffe fortement le résistor. Or la résistivité d'un conducteur métallique augmente avec la température : la résistance réelle n'est plus \SI{20}{\ohm} mais $R_{\text{réelle}} = \dfrac{215}{10} = \SI{21,5}{\ohm}$. L'hypothèse « $R$ constante égale à \SI{20}{\ohm} » est mise en défaut par l'échauffement. La leçon essentielle : \textbf{l'hypothèse « résistance constante » cesse d'être vérifiée dès que l'échauffement devient important}.
\end{exoresolu}

\begin{attention}[Extrapoler sans vérifier : l'erreur classique]
L'erreur la plus fréquente — et la plus sanctionnée au Probatoire — consiste à \textbf{extrapoler} une loi établie sur une petite gamme de mesures à des valeurs très différentes, sans vérification expérimentale. Avoir vérifié $U = RI$ entre \SI{0}{A} et \SI{0,5}{A} ne prouve rien quant à son comportement à \SI{10}{A}. De même, $PV = nRT$ vérifiée entre 1 et \SI{3}{bar} ne dit rien à \SI{200}{bar}. L'\cle{extrapolation} (prévoir hors de la gamme mesurée) est toujours risquée : elle doit être annoncée comme telle et justifiée, ou confirmée par de nouvelles mesures. L'\emph{interpolation} (prévoir entre deux mesures) est en général beaucoup plus sûre.
\end{attention}

\courssub{Visualiser la sortie du domaine de validité}

Le graphique ci-dessous représente une situation typique : tant que l'on reste dans la gamme de fonctionnement « doux », les points de mesure s'alignent sur la droite du modèle ; au-delà d'une valeur seuil, ils s'en écartent progressivement.

\begin{popfigure}
\begin{center}
\begin{tikzpicture}
\begin{axis}[
    width=0.9\linewidth,
    height=7cm,
    xlabel={Intensité $I$ (A)},
    ylabel={Tension $U$ (V)},
    xmin=0, xmax=12,
    ymin=0, ymax=260,
    grid=both,
    grid style={popline!40},
    legend style={at={(0.03,0.97)}, anchor=north west, font=\small},
]
% Zone de validité
\fill[popgreen!18] (axis cs:0,0) rectangle (axis cs:6,260);
% Zone hors domaine
\fill[poporange!18] (axis cs:6,0) rectangle (axis cs:12,260);
% Ligne du modèle U = 20 I
\addplot[popcurve, domain=0:12, samples=50] {20*x};
\addlegendentry{Modèle $U = 20\,I$}
% Mesures réelles qui s'écartent (courbure vers le haut pour R qui augmente)
\addplot[only marks, mark=*, mark size=2pt, poppurple] coordinates {
    (0.5,10.2) (1,19.8) (2,40.5) (3,59.5) (4,80.8) (5,99.2)
    (6,121) (7,142) (8,164) (9,186) (10,215) (11,245)
};
\addlegendentry{Mesures réelles}
% Ligne de seuil
\draw[dashed, popdark, thick] (axis cs:6,0) -- (axis cs:6,260);
\node[popdark, font=\small, anchor=south, rotate=90] at (axis cs:6,130) {seuil};
% Annotations des zones
\node[popgreen!60!black, font=\small\bfseries] at (axis cs:3,235) {Domaine de validité};
\node[poporange!70!black, font=\small\bfseries, align=center] at (axis cs:9.2,235) {Modèle\\inadapté};
\end{axis}
\end{tikzpicture}
\end{center}
\legende{Caractéristique $U(I)$ d'un dipôle ohmmique à température constante (droite) et d'un résistor métallique s'échauffant (points) : les mesures suivent la droite dans la zone verte, puis s'en écartent vers le haut (zone orange) lorsque l'échauffement augmente la résistance.}
\end{popfigure}

La lecture de ce graphique est un savoir-faire exigible : l'écart entre la courbe réelle et la droite du modèle croît avec $I$, et c'est précisément cet écart, rapporté à la mesure, qui donne l'écart relatif $\varepsilon$ permettant de fixer la frontière du domaine de validité (ici autour de $I \approx \SI{6}{A}$ si le seuil est 5 \%).

\courssub{Deux exemples emblématiques de dépassement de domaine}

\textbf{La lampe à incandescence : la loi d'Ohm hors domaine.} Le filament de tungstène d'une lampe à incandescence (celles que l'on trouve encore dans de nombreuses maisons lors des délestages, alimentées par des groupes électrogènes) est porté à environ \SI{2700}{\celsius} en fonctionnement. Sa résistance à chaud est environ \textbf{quinze fois plus grande} qu'à froid. Il est donc impossible de décrire cette lampe par une loi $U = RI$ avec $R$ constante : la caractéristique $U(I)$ est courbe. La loi d'Ohm n'est pas « fausse » : son hypothèse implicite (température constante) est violemment violée. Le modèle adapté fait intervenir une résistance $R(T)$ dépendant de la température.

\textbf{Le gaz comprimé : la loi des gaz parfaits hors domaine.} Pour de l'air à \SI{1}{bar} et \SI{20}{\celsius}, la loi $PV = nRT$ est excellente (écarts inférieurs à 1 \%). Mais dans une bonbonne d'air comprimé à \SI{200}{bar}, le volume propre des molécules n'est plus négligeable devant le volume offert au gaz, et les forces d'attraction entre molécules deviennent sensibles : les écarts atteignent plusieurs dizaines de pourcents. Les ingénieurs utilisent alors des modèles plus riches, comme l'équation de Van der Waals, qui ajoute deux termes correctifs. Là encore, la loi des gaz parfaits n'est pas fausse : elle est \emph{inadaptée} hors de son domaine (pressions modérées, températures éloignées de la liquéfaction).

\begin{savaistu}[Pourquoi les pétroliers ne se contentent pas des gaz parfaits]
Dans les installations gazières de Kribi ou de Logbaba, le gaz naturel est transporté et stocké sous fortes pressions. Les ingénieurs qui dimensionnent les pipelines et les réservoirs ne peuvent pas utiliser $PV = nRT$ sans correction : ils emploient des « facteurs de compressibilité » mesurés expérimentalement. Une erreur de modèle à ce niveau, ce sont des millions de francs CFA de matériel sous-dimensionné — ou des risques d'accident. Choisir le bon modèle dans son domaine de validité est un enjeu économique et de sécurité, pas seulement un exercice scolaire.
\end{savaistu}

\courssub{L'attitude scientifique face aux limites d'un modèle}

\begin{aretenir}[Une loi hors domaine n'est pas « fausse »]
Une loi utilisée hors de son domaine de validité n'est pas \emph{fausse} : elle est simplement \cle{inadaptée} à la situation. La démarche scientifique consiste alors à identifier l'hypothèse mise en défaut, puis à construire un \textbf{modèle plus riche} qui la remplace. L'histoire de la physique avance ainsi : la mécanique de Newton, inadaptée aux vitesses proches de celle de la lumière, a été enrichie — et non détruite — par la relativité d'Einstein, qui la contient comme cas particulier aux faibles vitesses.
\end{aretenir}

Cette attitude se résume par l'arbre de décision suivant, que tu dois pouvoir appliquer face à n'importe quelle situation-problème :

\begin{popfigure}
\begin{center}
\begin{tikzpicture}[
    font=\small,
    decision/.style={diamond, draw=popblue, fill=popblueL, text width=3.2cm, align=center, aspect=2, inner sep=1pt},
    action/.style={rectangle, rounded corners, draw=popgreen, fill=popgreenL, text width=3.4cm, align=center},
    stop/.style={rectangle, rounded corners, draw=poporange, fill=poporangeL, text width=3.4cm, align=center},
    fleche/.style={-{Stealth[length=2.5mm]}, thick, popdark},
    node distance=1.6cm and 2.2cm
]
\node[decision] (hyp) {Les hypothèses du modèle sont-elles vérifiées ?};
\node[decision, below=of hyp] (dom) {Les grandeurs sont-elles dans le domaine mesuré ?};
\node[action, below=of dom] (ok) {Appliquer le modèle et calculer la prévision};
\node[stop, right=of hyp] (rev) {Réviser les hypothèses : chercher un modèle plus riche};
\node[stop, right=of dom] (mes) {Prudence : extrapolation. Faire de nouvelles mesures};
\node[decision, below=of ok] (test) {Écart relatif $<$ seuil fixé ?};
\node[action, below=of test] (fin) {Modèle validé dans ces conditions};
\node[stop, right=of test] (lim) {Sortie du domaine : réviser le modèle};

\draw[fleche] (hyp) -- node[left]{oui} (dom);
\draw[fleche] (hyp) -- node[above]{non} (rev);
\draw[fleche] (dom) -- node[left]{oui} (ok);
\draw[fleche] (dom) -- node[above]{non} (mes);
\draw[fleche] (ok) -- (test);
\draw[fleche] (test) -- node[left]{oui} (fin);
\draw[fleche] (test) -- node[above]{non} (lim);
\end{tikzpicture}
\end{center}
\legende{Arbre de décision : le modèle est-il applicable ? On vérifie successivement les hypothèses, le domaine des grandeurs, puis la concordance entre prévisions et mesures.}
\end{popfigure}

\begin{exercice}[Gaz parfait ou gaz réel ?]
Une bouteille de volume $V = \SI{50}{L}$ contient $n = \SI{40}{mol}$ d'air à $\theta = \SI{27}{\celsius}$. On donne $R = \SI{8,31}{J.K^{-1}.mol^{-1}}$.

1. Calculer la pression prévue par la loi des gaz parfaits (convertir la température en kelvins et le volume en mètres cubes).
2. La pression mesurée est $P_{\text{mes}} = \SI{1,90e6}{Pa}$. Calculer l'écart relatif.
3. Avec un seuil de 5 \%, la loi des gaz parfaits est-elle valide ici ? Justifier en citant l'hypothèse mise en défaut.
\end{exercice}

\begin{corrige}[Exercice : gaz parfait ou gaz réel ?]
1. Conversion des unités SI : $T = 27 + 273 = \SI{300}{K}$ et $V = \SI{50}{L} = \SI{0,050}{m^3}$. La loi des gaz parfaits donne :
\[
P_{\text{mod}} = \frac{nRT}{V} = \frac{40 \times 8{,}31 \times 300}{0{,}050} = \SI{1,99e6}{Pa} \approx \SI{20}{bar}.
\]

2. Écart relatif :
\[
\varepsilon = \frac{|1{,}90 - 1{,}99| \times 10^{6}}{1{,}90 \times 10^{6}} = \frac{0{,}09}{1{,}90} \approx 4,7\,\%.
\]

3. $\varepsilon \approx 4,7\,\% < 5\,\%$ : la loi des gaz parfaits reste \textbf{acceptable} à cette pression, mais on est proche de la limite. L'écart, déjà sensible, provient des interactions entre molécules et de leur volume propre, qui ne sont plus tout à fait négligeables à \SI{20}{bar}. À \SI{200}{bar}, l'écart dépasserait largement le seuil et un modèle de gaz réel s'imposerait.
\end{corrige}

\begin{aretenir}[L'essentiel de la section]
\begin{itemize}
\item Le \textbf{domaine de validité} d'une loi est l'ensemble des conditions où ses prévisions s'accordent avec les mesures ; chaque hypothèse simplificatrice en trace une frontière.
\item On teste la validité par la mesure : prévision $\rightarrow$ mesure $\rightarrow$ écart relatif $\varepsilon = \dfrac{|\text{mesure} - \text{modèle}|}{\text{mesure}}$ $\rightarrow$ comparaison à un seuil (souvent 5 \%).
\item Hors de son domaine, une loi n'est pas fausse mais \textbf{inadaptée} : on construit un modèle plus riche (lampe à incandescence, gaz comprimé).
\item \textbf{Extrapoler} hors de la gamme mesurée est risqué et doit toujours être justifié ou vérifié expérimentalement.
\end{itemize}
\end{aretenir}

Tu es maintenant capable de juger de la pertinence d'un modèle face à des données réelles : c'est exactement la compétence que la synthèse de la prochaine section mettra en scène, en te plaçant dans la peau d'un scientifique confronté à des documents expérimentaux complets.

\coursec{Synthèse : travailler comme un scientifique}

Nous venons de voir que toute loi physique possède des \cle{contraintes} : un domaine de validité, des hypothèses simplificatrices, des limites au-delà desquelles elle cesse de décrire correctement la réalité. La loi d'Ohm ne s'applique qu'aux conducteurs ohmiques dont la température reste sensiblement constante ; la loi des gaz parfaits ne décrit bien les gaz qu'à faible pression et à température modérée. Cette constatation n'est pas un aveu de faiblesse de la science : c'est au contraire sa plus grande force. Un physicien ne dit jamais « c'est vrai », il dit « c'est vrai \emph{dans telles conditions, à telle précision près} ». Cette section fait le bilan de la démarche complète que tu as suivie dans cette leçon, afin que tu puisses la réutiliser toute l'année — et au Probatoire.

\courssub{Le cycle de la modélisation : une démarche en boucle}

\begin{aretenir}[La démarche de modélisation en six étapes]
Modéliser un phénomène physique, c'est suivre un cycle que l'on peut résumer ainsi :
\begin{enumerate}
\item \textbf{Observer} un phénomène et identifier les \cle{grandeurs} pertinentes (celles qui semblent influencer le phénomène) ;
\item \textbf{Formuler des hypothèses simplificatrices} : on néglige volontairement certains effets (frottements, masse d'un fil, échanges thermiques...) pour rendre le problème traitable ;
\item \textbf{Construire un modèle}, c'est-à-dire une représentation simplifiée de la réalité (le résistor idéal, le gaz parfait, le point matériel...) ;
\item \textbf{Établir une loi mathématique} reliant les grandeurs (souvent par analyse des mesures : tableau, graphe, recherche de proportionnalité) ;
\item \textbf{Tester la loi} : comparer les prévisions du modèle à de \emph{nouvelles} mesures, en tenant compte des incertitudes ;
\item \textbf{Valider ou réviser} : si l'accord est bon dans le domaine prévu, la loi est validée (provisoirement !) ; sinon, on révise les hypothèses et on recommence.
\end{enumerate}
\end{aretenir}

Le point essentiel est que cette démarche est une \cle{boucle} : un modèle n'est jamais définitif. Il est valable tant qu'aucune mesure suffisamment précise ne le met en défaut. C'est exactement ainsi qu'a progressé la physique : le modèle du gaz parfait a été complété par l'équation de Van der Waals pour les gaz réels ; la mécanique de Newton a été complétée par la relativité d'Einstein aux grandes vitesses.

\begin{popfigure}
\begin{center}
\begin{tikzpicture}[>=Stealth, node distance=0pt]
% Nœuds disposés en cercle
\node[draw=popblue, fill=popblueL, rounded corners, thick, minimum width=2.6cm, minimum height=0.9cm, align=center] (obs) at (90:3.1) {\textbf{1. Observer}\\ grandeurs pertinentes};
\node[draw=poporange, fill=poporangeL, rounded corners, thick, minimum width=2.6cm, minimum height=0.9cm, align=center] (hyp) at (30:3.1) {\textbf{2. Hypothèses}\\ simplificatrices};
\node[draw=poppurple, fill=poppurpleL, rounded corners, thick, minimum width=2.6cm, minimum height=0.9cm, align=center] (mod) at (-30:3.1) {\textbf{3. Modèle}\\ loi mathématique};
\node[draw=popgreen, fill=popgreenL, rounded corners, thick, minimum width=2.6cm, minimum height=0.9cm, align=center] (prev) at (-90:3.1) {\textbf{4. Prévision}\\ calcul};
\node[draw=poppink, fill=poppinkL, rounded corners, thick, minimum width=2.6cm, minimum height=0.9cm, align=center] (mes) at (-150:3.1) {\textbf{5. Mesure}\\ avec incertitudes};
\node[draw=popgold, fill=popgoldL, rounded corners, thick, minimum width=2.6cm, minimum height=0.9cm, align=center] (test) at (150:3.1) {\textbf{6. Accord ?}\\ comparaison};
% Flèches du cycle
\draw[->, very thick, popdark, bend left=18] (obs) to (hyp);
\draw[->, very thick, popdark, bend left=18] (hyp) to (mod);
\draw[->, very thick, popdark, bend left=18] (mod) to (prev);
\draw[->, very thick, popdark, bend left=18] (prev) to (mes);
\draw[->, very thick, popdark, bend left=18] (mes) to (test);
\draw[->, very thick, popgreen, bend left=18] (test) to node[midway, above left, align=center, font=\small]{oui : \textbf{valider}\\ (dans le domaine)} (obs);
% Retour de révision
\draw[->, very thick, poppink, dashed] (test.east) .. controls (2.2,2.6) .. node[midway, above, align=center, font=\small]{non : \textbf{réviser}\\ les hypothèses} (hyp.north west);
\end{tikzpicture}
\end{center}
\legende{Le cycle de la modélisation : tant que la mesure confirme la prévision, la loi est validée dans son domaine ; sinon, on révise les hypothèses et on reconstruit le modèle.}
\end{popfigure}

\begin{methode}[Fiche réflexe « modéliser »]
Face à une situation nouvelle (au Probatoire comme en TP), pose-toi systématiquement ces cinq questions :
\begin{enumerate}
\item \textbf{Quelles grandeurs ?} Liste les grandeurs physiques en jeu, avec leurs unités SI. Lesquelles puis-je faire varier ? Lesquelles dois-je maintenir constantes ?
\item \textbf{Quelles hypothèses ?} Qu'est-ce que je néglige (frottements, température variable, interactions entre molécules...) et pourquoi est-ce raisonnable ?
\item \textbf{Quelle loi ?} Quelle relation mathématique simple (proportionnalité, produit constant...) relie les grandeurs ? Vérifie-la par analyse dimensionnelle.
\item \textbf{Quel domaine de validité ?} Dans quelle plage de valeurs la loi a-t-elle été vérifiée ? Où risque-t-elle de cesser d'être valable ?
\item \textbf{Quel test expérimental ?} Quelle mesure nouvelle, avec quelle incertitude, permettrait de confirmer ou d'invalider la loi ?
\end{enumerate}
\end{methode}

\courssub{Tableau récapitulatif des deux lois étudiées}

Rassemblons dans un tableau à double entrée tout ce que nous avons établi sur la loi d'Ohm et la loi des gaz parfaits. Ce type de synthèse est un excellent outil de révision : il t'oblige à formuler chaque loi avec ses quatre « piliers » — relation, hypothèses, domaine, application.

\begin{popfigure}
\begin{center}
\begin{tabularx}{\linewidth}{|l|X|X|}
\hline
\rowcolor{popblue}
\textbf{\color{white} Critère} & \textbf{\color{white} Loi d'Ohm} & \textbf{\color{white} Loi des gaz parfaits} \\
\hline
\textbf{Relation} &
$U = R\,I$ \quad $U$ en V, $I$ en A, $R$ en $\Omega$ &
$P\,V = n\,R\,T$ \quad $P$ en Pa, $V$ en m$^3$, $n$ en mol, $T$ en K, $R = \num{8,314}$ J\,K$^{-1}$\,mol$^{-1}$ \\
\hline
\textbf{Hypothèses simplificatrices} &
Conducteur ohmique ; température constante ; $R$ indépendante de $I$ et de $U$ &
Molécules ponctuelles (volume propre nul) ; aucune interaction entre molécules hors chocs \\
\hline
\textbf{Domaine de validité} &
Résistors métalliques, tant que l'échauffement reste faible ; exclut diodes, lampes à incandescence en régime chaud &
Gaz à faible pression ($P \lesssim$ quelques bars) et température modérée, loin de la liquéfaction \\
\hline
\textbf{Exemple d'application} &
Calibre d'un fusible, choix d'une résistance dans un circuit, prévision des chutes de tension sur une ligne &
Prévision de la pression d'une bouteille de gaz, météorologie, gonflage des pneus de mototaxi \\
\hline
\end{tabularx}
\end{center}
\legende{Synthèse comparative des deux lois-modèles étudiées dans cette leçon : chaque loi n'a de sens qu'accompagnée de ses hypothèses et de son domaine de validité.}
\end{popfigure}

\begin{attention}[Écrire une loi sans ses conditions d'application : l'erreur classique]
Beaucoup d'élèves écrivent « $U = RI$ » ou « $PV = nRT$ » comme s'il s'agissait de vérités absolues. C'est une erreur de physique \emph{et} une perte de points : au Probatoire, citer le \cle{domaine de validité} et les \cle{hypothèses} d'une loi est souvent explicitement valorisé dans le barème. Rédige ainsi : « le conducteur étant ohmique et sa température restant constante, la loi d'Ohm s'applique : $U = RI$ ». De même, n'oublie jamais que dans $PV = nRT$, la température doit être exprimée en \textbf{kelvins} et la pression en \textbf{pascals} : une température en degrés Celsius dans cette formule est une faute grave.
\end{attention}

\courssub{Le vocabulaire du scientifique : les mots justes}

\begin{definition}[Les mots-clés de la modélisation]
\begin{itemize}
\item Une \cle{grandeur} physique est une propriété mesurable d'un système (tension, pression, température, volume...). Elle s'exprime par une valeur numérique accompagnée d'une unité.
\item Une \cle{hypothèse simplificatrice} est un choix délibéré de négliger un effet jugé secondaire pour rendre le problème traitable (exemple : « la température du résistor reste constante »).
\item Un \cle{modèle} est une représentation simplifiée de la réalité, construite sur des hypothèses, qui permet de décrire et de prévoir un phénomène (exemples : le résistor idéal, le gaz parfait).
\item Une \cle{loi} physique est une relation mathématique entre grandeurs, établie et confirmée par l'expérience, valable dans un domaine précis.
\item Le \cle{domaine de validité} d'une loi est l'ensemble des conditions (plages de valeurs, nature du système) dans lesquelles elle décrit correctement les mesures.
\item Deux grandeurs sont \cle{proportionnelles} lorsque leur rapport est constant ; leur graphe est alors une droite passant par l'origine.
\end{itemize}
\end{definition}

\courssub{Le lien avec les incertitudes : le juge de paix de la validation}

Comment décider, à l'étape 6 du cycle, qu'il y a « accord » entre la prévision et la mesure ? C'est ici que la Leçon N01 (mesures et incertitudes) devient indispensable. Une mesure n'est jamais une valeur exacte : c'est un \emph{intervalle}. On compare donc la prévision du modèle au résultat expérimental sous la forme
\[ m_{\text{mes}} - U(m) \;\le\; m_{\text{prévu}} \;\le\; m_{\text{mes}} + U(m), \]
où $U(m)$ est l'incertitude élargie sur la mesure. Si la valeur prévue par la loi tombe dans cet intervalle, le modèle est \emph{compatible} avec l'expérience ; sinon, il doit être révisé — ou l'expérience doit être refaite plus soigneusement.

\begin{exoresolu}[Valider ou invalider la loi d'Ohm ?]
Énoncé. Un groupe d'élèves du lycée de Nkongsamba étudie un résistor étiqueté $R = \SI{100}{\Omega}$. Pour une tension mesurée $U = \SI{4,92}{V}$ (incertitude élargie $U(U) = \SI{0,06}{V}$), ils mesurent un courant $I = \SI{51,0}{mA}$ (incertitude élargie $U(I) = \SI{1,0}{mA}$). La loi d'Ohm avec $R = \SI{100}{\Omega}$ est-elle compatible avec cette mesure ?
\tcblower
Solution détaillée. La loi d'Ohm prévoit, pour $I = \SI{51,0}{mA} = \SI{0,0510}{A}$ :
\[ U_{\text{prévu}} = R\,I = 100 \times \num{0,0510} = \SI{5,10}{V}. \]
L'intervalle de mesure de la tension est :
\[ U_{\text{mes}} \in [\num{4,92} - \num{0,06}\,;\ \num{4,92} + \num{0,06}] = [\SI{4,86}{V}\,;\ \SI{4,98}{V}]. \]
La valeur prévue $\SI{5,10}{V}$ est \emph{en dehors} de cet intervalle. Deux conclusions sont possibles :
\begin{itemize}
\item soit la valeur affichée $R = \SI{100}{\Omega}$ est inexacte (les résistors courants ont une tolérance de 5 \%, donc $R$ peut valoir jusqu'à $\SI{105}{\Omega}$ ; avec $R = \SI{96,5}{\Omega}$, la prévision $U = RI = \SI{4,92}{V}$ serait parfaitement compatible) ;
\item soit le résistor s'est échauffé et sa résistance a varié : l'hypothèse « température constante » n'est plus respectée.
\end{itemize}
Le groupe reprend la mesure avec la tolérance du fabricant : $R = \SI{100}{\Omega} \pm 5\,\%$, soit $R \in [\SI{95}{\Omega}\,;\ \SI{105}{\Omega}]$. La prévision devient $U_{\text{prévu}} \in [\SI{4,85}{V}\,;\ \SI{5,36}{V}]$, qui recoupe largement l'intervalle de mesure. \textbf{Conclusion : la loi d'Ohm est validée} pour ce résistor, à condition de tenir compte de la tolérance du fabricant — c'est encore une hypothèse du modèle qui a dû être précisée.
\end{exoresolu}

Retiens la morale de cet exercice : \textbf{on ne valide jamais une loi en comparant deux nombres « secs »} ; on compare une prévision à un \emph{intervalle} de mesure. C'est toute la différence entre un calcul scolaire et un raisonnement de scientifique.

\courssub{Ouverture : les modèles que tu rencontreras cette année}

La démarche que tu viens d'apprendre va te servir dans presque tous les chapitres de l'année. En voici un avant-goût :
\begin{itemize}
\item En \textbf{optique}, le \cle{modèle de la lentille mince} (épaisseur négligeable, rayons proches de l'axe) conduira aux relations de conjugaison, valables seulement dans les conditions de Gauss ;
\item En \textbf{mécanique}, le \cle{modèle du point matériel} (toute la masse concentrée en un point) et l'hypothèse « frottements négligeables » permettront d'appliquer les lois de Newton au mouvement d'une mototaxi ou d'un ballon de football ;
\item En \textbf{chimie}, le \cle{modèle de la solution diluée} (les espèces dissoutes n'interagissent pas entre elles) justifiera la loi de Beer-Lambert utilisée en TP de dosage ;
\item En \textbf{électricité}, la loi d'Ohm sera généralisée aux circuits complets avec la loi de Pouillet.
\end{itemize}
À chaque fois, pose-toi les cinq questions de la fiche réflexe : tu verras que la structure du raisonnement est toujours la même.

\begin{savaistu}[Les modèles qui éclairent Douala et Yaoundé]
Les ingénieurs d'ENEO ne « bricolent » pas le réseau électrique : ils le \emph{modélisent}. Chaque ligne de transport est représentée par une résistance, une inductance et une capacité équivalentes ; chaque quartier est modélisé par sa consommation. Grâce à la loi d'Ohm généralisée et aux lois de Kirchhoff, des logiciels calculent les chutes de tension le long des lignes et prévoient où la tension risque de descendre sous le seuil acceptable aux heures de pointe (le soir, quand les quartiers de Douala et Yaoundé s'illuminent). Ces modèles guident les décisions : où installer un nouveau transformateur, quel câble choisir, comment raccorder un nouveau quartier. Et comme tout modèle, ils sont confrontés en permanence aux mesures réelles du réseau — exactement le cycle que tu viens d'apprendre. Les mêmes principes servent à dimensionner les barrages hydroélectriques de Lom Pangar et Nachtigal qui alimentent le pays.
\end{savaistu}

\begin{exercice}[Pour t'entraîner]
Un élève affirme : « J'ai vérifié la loi des gaz parfaits, donc elle est vraie pour tous les gaz dans toutes les conditions. »
\begin{enumerate}
\item Relever les deux erreurs de raisonnement contenues dans cette phrase.
\item Proposer une expérience simple permettant de mettre la loi en défaut (indication : penser à une bouteille de gaz fortement comprimé ou refroidi).
\item Rédiger en trois lignes une conclusion scientifiquement correcte à la place de l'affirmation de l'élève.
\end{enumerate}
\end{exercice}

\begin{corrige}[Exercice]
\begin{enumerate}
\item Première erreur : une vérification expérimentale, même réussie, ne « prouve » pas une loi de façon absolue ; elle la rend seulement compatible avec les mesures réalisées, compte tenu des incertitudes. Deuxième erreur : la validité d'une loi est limitée à son domaine ; la loi des gaz parfaits repose sur des hypothèses (molécules ponctuelles, sans interactions) qui échouent à forte pression ou basse température.
\item Comprimer fortement un gaz (par exemple mesurer le produit $PV$ d'une bouteille de gaz à plusieurs dizaines de bars) : on constate que $PV$ n'est plus proportionnel à $nT$, car le volume propre des molécules et leurs interactions ne sont plus négligeables. On peut aussi refroidir un gaz près de sa température de liquéfaction.
\item Rédaction correcte : « Dans le domaine exploré (pression inférieure à quelques bars, température ambiante), les mesures sont compatibles avec la loi $PV = nRT$ compte tenu des incertitudes ; la loi des gaz parfaits est donc validée dans ce domaine, sous réserve de ses hypothèses. »
\end{enumerate}
\end{corrige}

\begin{aretenir}[L'essentiel de la leçon]
\begin{itemize}
\item Modéliser, c'est suivre le cycle : \textbf{observer} $\rightarrow$ \textbf{hypothèses} $\rightarrow$ \textbf{modèle et loi} $\rightarrow$ \textbf{prévision} $\rightarrow$ \textbf{mesure} $\rightarrow$ \textbf{validation ou révision}.
\item Une loi physique n'existe jamais seule : elle s'énonce toujours avec ses \textbf{hypothèses} et son \textbf{domaine de validité}.
\item La validation d'un modèle repose sur la comparaison entre prévision et mesure \textbf{avec leurs incertitudes}.
\item Cette démarche sera réutilisée toute l'année : lentille mince, point matériel, solution diluée... Tu penses désormais comme un scientifique.
\end{itemize}
\end{aretenir}

\newpage

\coursec{Activité d'intégration}

\coursec{Modéliser un phénomène : exemples de modèles utilisés en physique et en chimie}

\courssub{Qu'est-ce qu'un modèle en physique ?}

\begin{definition}[Modèle physique]
Un \cle{modèle physique} est une représentation simplifiée et idéalisée d'un système réel, construite à l'aide de concepts et de lois mathématiques, dans le but d'expliquer des phénomènes observés et de faire des prévisions quantitatives. Le passage du réel au modèle repose sur des \cle{hypothèses simplificatrices} (négliger le frottement, considérer un fil comme indéformable, traiter un astre comme un point matériel, etc.).
\end{definition}

\begin{propriete}[Le cycle de la modélisation]
La démarche scientifique de modélisation suit un cycle itératif :
\begin{enumerate}
\item \textbf{Observation / Situation réelle} : collecte de données, identification des grandeurs pertinentes.
\item \textbf{Modélisation} : formulation d'hypothèses, choix d'un cadre théorique, écriture des équations du modèle.
\item \textbf{Prédiction / Simulation} : résolution des équations, calcul de valeurs prévisibles.
\item \textbf{Confrontation à l'expérience} : mesure, calcul d'incertitudes, comparaison prédiction/expérience (écart relatif).
\item \textbf{Validation / Raffinement} : si l'accord est satisfaisant dans le domaine de validité, le modèle est validé ; sinon, on modifie les hypothèses (ex. : ajouter un frottement, rendre une résistance dépendante de la température).
\end{enumerate}
\end{propriete}

\begin{methode}[Formuler des hypothèses simplificatrices]
Pour modéliser une situation concrète :
\begin{enumerate}
\item Identifier le \textbf{système étudié} et son \textbf{environnement}.
\item Lister les \textbf{grandeurs physiques} en jeu (entrées, sorties, paramètres).
\item Énoncer explicitement les \textbf{hypothèses} (ex. : « le transformateur est une source de tension idéale », « la température des câbles reste constante », « le régime est stationnaire »).
\item Justifier chaque hypothèse par un ordre de grandeur ou une référence (ex. : « la résistance interne du transformateur est négligeable devant $r$ »).
\end{enumerate}
\end{methode}

\courssub{Établir une loi à partir de mesures : la démarche expérimentale}

\begin{methode}[Déterminer une loi mathématique à partir de données expérimentales]
\begin{enumerate}
\item \textbf{Tracer le graphique} pertinent : $y = f(x)$ ou une forme linéarisée (ex. $\ln y = f(x)$, $y = f(1/x)$...). Choisir les échelles, reporter les points avec leurs barres d'incertitude si connues.
\item \textbf{Identifier la forme mathématique} : alignement (proportionnalité ou relation affine), courbe exponentielle, parabolique, hyperbolique...
\item \textbf{Déterminer les paramètres} : lecture graphique (ordonnée à l'origine, pente) ou régression linéaire (méthode des moindres carrés si exigé). Vérifier l'homogénéité dimensionnelle de la loi obtenue.
\item \textbf{Écrire la loi} sous la forme $y = ax + b$ ou $y = kx^n$, en précisant les unités de $a$, $b$, $k$, $n$.
\item \textbf{Tester la loi} sur un point de mesure non utilisé pour l'ajustement (point de contrôle) via l'écart relatif $\varepsilon = \frac{|y_{\text{prévu}} - y_{\text{mesuré}}|}{y_{\text{mesuré}}}$.
\end{enumerate}
\end{methode}

\courssub{Exemple 1 : La loi d'Ohm et la chute de tension dans une ligne électrique (Activité d'intégration)}

\courssub{Objectifs de l'activité}

À l'issue de cette activité d'intégration, tu seras capable de :
\begin{itemize}
\item \cle{formuler des hypothèses simplificatrices} pour modéliser une situation concrète (une ligne électrique réelle) ;
\item \cle{établir une loi mathématique simple} entre des grandeurs mesurées à partir d'un tableau de données et d'un graphique ;
\item \cle{tester la validité des hypothèses émises} par la mesure, en comparant une prévision du modèle à une donnée expérimentale ;
\item \cle{discuter le domaine de validité} d'un modèle et identifier les hypothèses susceptibles d'être remises en cause ;
\item rédiger un court texte argumenté, à la manière d'un scientifique, pour expliquer un phénomène de la vie quotidienne.
\end{itemize}

\courssub{Situation}

\begin{attention}[Le contexte]
Dans le quartier Nkolbisson à Yaoundé, les habitants se plaignent : chaque soir, vers 19~h, lorsque tout le monde allume la télévision, le fer à repasser et les climatiseurs, les lampes deviennent plus faibles et certains appareils s'éteignent. Un technicien de la SONATREL affirme dans un rapport : \og la tension délivrée en bout de ligne diminue lorsque l'intensité totale appelée par le quartier augmente ; nos mesures montrent que la chute est proportionnelle à l'intensité \fg. Un habitant sceptique rétorque : \og la SONATREL a tort, la tension devrait toujours être de 220~V ! \fg. Qui a raison ? À toi de trancher... comme un scientifique.
\end{attention}

Le quartier est alimenté par un transformateur situé à environ \SI{1,2}{km}. Le technicien a relevé, un mardi, la tension $U$ disponible au compteur général du quartier et l'intensité totale $I$ appelée par l'ensemble des habitations, à différentes heures de la journée :

\begin{center}
\begin{tabular}{|c|ccccccc|}\hline
\rowcolor{popblueL} Heure & 6~h & 10~h & 14~h & 17~h & 19~h & 20~h & 22~h \\ \hline
$I$ (A) & 5,0 & 10,0 & 15,0 & 20,0 & 25,0 & 30,0 & 35,0 \\ \hline
$U$ (V) & 219,5 & 219,0 & 218,5 & 218,0 & 217,5 & 217,0 & 216,5 \\ \hline
\end{tabular}
\end{center}

Une mesure supplémentaire, effectuée à 19~h~30 (heure de pointe maximale), donne : $I = \SI{40,0}{A}$ et $U = \SI{215,8}{V}$.

Le technicien propose le modèle suivant : le transformateur se comporte comme une \cle{source de tension constante} de force électromotrice $U_0$, et l'ensemble des câbles de la ligne (aller et retour) se comporte comme un \cle{unique résistor} de résistance $r$, placé en série entre la source et le quartier.

\courssub{Consignes}

\begin{enumerate}
\item \textbf{Formuler les hypothèses.} Reproduis le schéma électrique du modèle proposé (source idéale de tension $U_0$, résistor $r$, quartier parcouru par le courant $I$). Cite les trois hypothèses simplificatrices contenues dans ce modèle.
\item \textbf{Exploiter les mesures.} Trace sur papier millimétré (ou avec le graphique fourni) la courbe $U = f(I)$ pour les sept premières mesures. Que constates-tu ? Quelle relation mathématique simple cela suggère-t-il entre $U$ et $I$ ?
\item \textbf{Établir la loi.} À l'aide de la loi d'Ohm appliquée au modèle, montre que la tension en bout de ligne s'écrit $U = U_0 - r\,I$. Détermine graphiquement les valeurs de $U_0$ et de $r$, en précisant leurs unités.
\item \textbf{Tester le modèle.} Utilise la loi établie pour prévoir la tension $U_{\text{prévue}}$ lorsque $I = \SI{40,0}{A}$. Compare à la mesure réelle $U_{\text{mesurée}} = \SI{215,8}{V}$ en calculant l'écart relatif $\dfrac{|U_{\text{prévue}} - U_{\text{mesurée}}|}{U_{\text{mesurée}}}$ en pourcentage. Le modèle est-il validé à ce niveau de précision ?
\item \textbf{Discuter le domaine de validité.} La résistance d'un conducteur augmente avec sa température. Que se passe-t-il aux heures de pointe, lorsque la ligne est parcourue par un courant intense depuis longtemps ? Quelle hypothèse du modèle est alors remise en cause ? Comment la courbe $U = f(I)$ serait-elle modifiée ?
\item \textbf{Rédiger la conclusion.} Rédige un texte de 8 à 12 lignes, destiné à l'habitant sceptique, expliquant pourquoi la tension \og baisse \fg{} le soir et pourquoi le technicien de la SONATREL a (presque) raison.
\end{enumerate}

\courssub{Ressources à mobiliser}

\begin{aretenir}[Boîte à outils de l'activité]
\begin{itemize}
\item \textbf{Loi d'Ohm} : aux bornes d'un résistor, $U_R = r\,I$, avec $U_R$ en volts (V), $r$ en ohms ($\Omega$) et $I$ en ampères (A).
\item \textbf{Additivité des tensions} (maille série) : $U_0 = U_r + U$, donc $U = U_0 - r\,I$.
\item \textbf{Lecture graphique} : pour une droite $y = a\,x + b$, l'ordonnée à l'origine $b$ se lit à $x = 0$ et le coefficient directeur se calcule par $a = \dfrac{\Delta y}{\Delta x}$ entre deux points éloignés de la droite.
\item \textbf{Écart relatif} : $\varepsilon = \dfrac{|\text{valeur prévue} - \text{valeur mesurée}|}{\text{valeur mesurée}}$, exprimé en $\%$.
\item \textbf{Effet thermique} : pour un conducteur métallique, la résistance croît avec la température ; l'effet Joule ($P = r\,I^2$) échauffe la ligne.
\end{itemize}
\end{aretenir}

\courssub{Démarche et corrigé commenté}

\begin{exoresolu}[La SONATREL a-t-elle raison ?]
Énoncé : à partir des sept mesures $(I, U)$ du tableau et de la mesure de contrôle $(\SI{40,0}{A} ; \SI{215,8}{V})$, établir la loi $U = f(I)$ de la ligne, la tester, puis discuter son domaine de validité.
\tcblower
\textbf{Étape 1 : schéma et hypothèses.}

Le modèle proposé se schématise ainsi :

\begin{popfigure}
\begin{tikzpicture}[european]
\draw (0,0) to[V=$U_0$] (0,2.5) to[R=$r$, i={$I$}] (4,2.5) to[R, l_={$U$ (quartier)}] (4,0) -- (0,0);
\end{tikzpicture}
\legende{Modèle de la ligne : source idéale $U_0$, résistance de ligne $r$, quartier récepteur.}
\end{popfigure}

Trois hypothèses simplificatrices sont faites :
\begin{enumerate}
\item le transformateur délivre une tension \textbf{constante} $U_0$, quelle que soit l'intensité demandée (source idéale de tension) ;
\item les \SI{1,2}{km} de câbles (aller et retour) sont remplacés par un \textbf{unique résistor} de résistance $r$ (on néglige l'inductance et la capacitance de la ligne) ;
\item cette résistance $r$ est \textbf{constante}, en particulier indépendante de la température des câbles (hypothèse isotherme).
\end{enumerate}

\textbf{Étape 2 : tracé de $U = f(I)$.}

\begin{popfigure}
\begin{tikzpicture}
\begin{axis}[width=0.9\linewidth, height=7cm,
xlabel={$I$ (A)}, ylabel={$U$ (V)},
xmin=0, xmax=45, ymin=214.5, ymax=221,
grid=major, grid style={popline!50},
xtick={0,5,10,15,20,25,30,35,40},
ytick={215,216,217,218,219,220},
scaled y ticks=false,
tick label style={/pgf/number format/use comma, /pgf/number format/fixed, /pgf/number format/precision=1},
legend style={at={(0.98,0.02)}, anchor=south east, fill=white, draw=popline}]
\addplot[only marks, mark=*, mark size=2.5pt, poporange] coordinates {
(5,219.5) (10,219) (15,218.5) (20,218) (25,217.5) (30,217) (35,216.5)};
\addlegendentry{Mesures (6h--22h)}
\addplot[popcurve, domain=0:45, samples=2, thick] {220 - 0.1*x};
\addlegendentry{Ajustement affine $U = 220 - 0,10\,I$}
\addplot[only marks, mark=square*, mark size=3pt, poppurple] coordinates {(40,215.8)};
\addlegendentry{Mesure de contrôle (19h30)}
\end{axis}
\end{tikzpicture}
\legende{Tension en bout de ligne en fonction de l'intensité appelée. Les points ronds sont alignés ; le carré violet est la mesure de contrôle à \SI{40}{A}.}
\end{popfigure}

Les sept points sont \textbf{alignés sur une droite décroissante} : $U$ est une fonction affine décroissante de $I$, de la forme $U = U_0 - r\,I$.

\textbf{Étape 3 : établissement de la loi.}

D'après l'additivité des tensions dans la maille unique (loi des mailles) : $U_0 = U_r + U$. La chute de tension sur la ligne est $U_r = r\,I$ (loi d'Ohm). D'où
\[ U = U_0 - r\,I. \]
Vérification dimensionnelle : $[r\,I] = \Omega \times \text{A} = \text{V}$, homogène à une tension. La loi est cohérente.

\emph{Ordonnée à l'origine} : en prolongeant la droite jusqu'à $I = 0$, on lit $U_0 = \SI{220}{V}$. C'est la tension à vide du transformateur (force électromotrice).

\emph{Coefficient directeur} (pente) : on choisit deux points éloignés de la droite de régression, par exemple ceux correspondant aux mesures extrêmes $(\SI{5}{A} ; \SI{219,5}{V})$ et $(\SI{35}{A} ; \SI{216,5}{V})$ :
\[ r = -\frac{\Delta U}{\Delta I} = -\frac{216,5 - 219,5}{35 - 5} = \frac{3,0}{30} = \SI{0,10}{\Omega}. \]
(L'unité $\Omega = \text{V}\cdot\text{A}^{-1}$ est bien respectée.)

La loi de la ligne s'écrit donc : $\boxed{U = 220 - 0,10\,I}$ (avec $U$ en V et $I$ en A).

\textbf{Étape 4 : test du modèle.}

Pour $I = \SI{40,0}{A}$ : $U_{\text{prévue}} = 220 - 0,10 \times 40,0 = \SI{216,0}{V}$.

Écart relatif avec la mesure $U_{\text{mesurée}} = \SI{215,8}{V}$ :
\[ \varepsilon = \frac{|216,0 - 215,8|}{215,8} = \frac{0,2}{215,8} \approx 9,27 \times 10^{-4} \approx 0,093\,\%. \]

L'écart est inférieur à $0,1\,\%$ : le modèle linéaire est \textbf{validé} sur toute la plage de mesures, y compris aux heures de pointe. Le technicien a raison : la chute de tension est bien proportionnelle à l'intensité (relation affine).

\textbf{Étape 5 : domaine de validité.}

Aux heures de pointe, le courant intense échauffe les câbles par effet Joule : la puissance dissipée par la ligne est $P = r\,I^2$. À \SI{40}{A}, $P = 0,10 \times 40^2 = \SI{160}{W}$ sur \SI{1,2}{km}, ce qui reste modéré. Cependant, la résistance d'un métal \textbf{augmente avec la température} (coefficient de température positif pour le cuivre/aluminium). L'hypothèse \og $r$ constante \fg{} n'est donc qu'approchée : si l'échauffement devenait significatif (ligne plus longue, section plus faible, canicule, surcharge prolongée), la valeur effective de $r$ croîtrait avec $I$, la chute de tension serait \textbf{plus forte} que prévue par la loi affine, et la courbe $U = f(I)$ s'incurverait vers le bas (concave) au lieu de rester une droite. Le modèle reste excellent ici car l'échauffement est faible, mais il faudrait le corriger (modèle $r = r_0(1 + \alpha \Delta T)$) pour des lignes plus chargées.

\textbf{Étape 6 : texte de conclusion (exemple de rédaction attendue).}

\og Le soir, vers 19~h, tous les habitants du quartier branchent leurs appareils : l'intensité totale appelée sur la ligne augmente fortement, passant de \SI{5}{A} en journée à \SI{40}{A} en pointe. Or les câbles qui relient le transformateur au quartier ne sont pas parfaits : ils se comportent comme un résistor de résistance $r \approx \SI{0,10}{\Omega}$ (aller-retour). D'après la loi d'Ohm, ce câble provoque une chute de tension $r\,I$, d'autant plus grande que le courant est fort. La tension disponible dans les maisons vaut donc $U = U_0 - r\,I$ : elle passe de \SI{220}{V} à vide à environ \SI{216}{V} aux heures de pointe. Les mesures confirment cette loi à mieux de $0,1\,\%$ près : le technicien de la SONATREL a raison. Toutefois, ce modèle suppose que la résistance des câbles reste constante ; si la ligne chauffe trop, sa résistance augmente et la tension chuterait davantage que prévu. \fg
\end{exoresolu}

\courssub{Bilan de l'activité}

Cette activité t'a fait vivre la \cle{démarche scientifique complète} sur un problème authentique du quotidien camerounais :

\begin{itemize}
\item \textbf{Modéliser}, c'est d'abord \emph{simplifier} : remplacer \SI{1,2}{km} de câbles réels par un seul résistor et un transformateur réel par une source idéale, ce sont des choix délibérés, pas des évidences.
\item Une loi physique ne se devine pas : elle \textbf{s'établit à partir de mesures}. Le tracé de $U = f(I)$ a révélé une relation affine, que la théorie (loi d'Ohm + loi des mailles) a ensuite justifiée : $U = U_0 - r\,I$.
\item Un modèle n'a de valeur que s'il \textbf{fait des prévisions testables} : la prédiction $U = \SI{216,0}{V}$ à \SI{40}{A}, confirmée à $0,09\,\%$ près, valide le modèle sur la plage étudiée.
\item Toute loi possède un \textbf{domaine de validité} : l'hypothèse \og $r$ constante \fg{} cesserait d'être acceptable si l'effet Joule échauffait notablement la ligne. Savoir dire \emph{jusqu'où} un modèle est fiable est aussi important que de savoir l'utiliser.
\item Enfin, tu as rédigé une explication argumentée : communiquer clairement un résultat scientifique fait partie du métier de scientifique... et t'a permis de départager le technicien et l'habitant sceptique !
\end{itemize}

\begin{attention}[Ce qu'il faut retenir pour le Probatoire]
Face à un tableau de mesures, la démarche type est : (1) formuler des hypothèses et proposer un modèle ; (2) tracer le graphe pertinent ; (3) identifier la loi mathématique (proportionnalité, relation affine...) et déterminer ses paramètres ; (4) tester la loi sur une mesure indépendante par un écart relatif ; (5) discuter le domaine de validité et les hypothèses fragiles. Cette démarche en cinq étapes est transposable à tous les chapitres de physique et de chimie de l'année.
\end{attention}

\courssub{Exemple 2 : La loi des gaz parfaits}

\begin{experience}[Expérience de pensée : compression d'un gaz dans une seringue]
\textbf{Matériel :} Seringue en verre (volume variable), manomètre, thermomètre, masses étalons (pour varier la pression), bain thermostaté.
\textbf{Protocole :}
\begin{enumerate}
\item \textbf{Isotherme} ($T$ constant) : on comprime le gaz lentement (pour rester à $T$ ambiante) et on relève $(P, V)$. On trace $P = f(1/V)$ : droite passant par l'origine $\rightarrow P \propto 1/V$ (Loi de \textbf{Mariotte}, 1662).
\item \textbf{Isobare} ($P$ constant) : on chauffe le gaz dans le bain, piston mobile libre. On relève $(V, T)$. Attention : $T$ en Kelvin ! On trace $V = f(T)$ : droite passant par l'origine $\rightarrow V \propto T$ (Loi de \textbf{Gay-Lussac}, 1802).
\item \textbf{Isochore} ($V$ constant) : piston bloqué, on chauffe. On relève $(P, T)$. $P \propto T$ (Loi de \textbf{Charles}, 1787).
\end{enumerate}
\textbf{Interprétation :} La synthèse des trois lois donne l'équation d'état du \cle{gaz parfait} : $PV = nRT$.
\end{experience}

\begin{definition}[Gaz parfait (ou gaz idéal)]
Un gaz parfait est un modèle de gaz dont les molécules sont ponctuelles (volume propre négligeable devant le volume du contenant) et n'interagissent pas entre elles (pas de forces intermoléculaires), sauf lors de chocs parfaitement élastiques avec les parois. Il obéit à l'équation d'état :
\[ PV = nRT \]
avec $P$ en \SI{}{Pa}, $V$ en \SI{}{m^3}, $n$ en \SI{}{mol}, $T$ en \SI{}{K} et $R = \SI{8,314}{J\cdot mol^{-1}\cdot K^{-1}}$ (constante des gaz parfaits).
\end{definition}

\begin{propriete}[Domaine de validité du modèle du gaz parfait]
Le modèle est excellent pour les \textbf{faibles pressions} (quelques bar) et \textbf{hautes températures} (loin de la condensation). Il échoue près du point critique, à haute pression (volume propre non négligeable, forces attractives) ou à très basse température (liquéfaction). L'équation de \textbf{Van der Waals} $\left(P + a\frac{n^2}{V^2}\right)(V - nb) = nRT$ corrige ces deux défauts ($a$ : attraction, $b$ : volume propre).
\end{propriete}

\begin{savaistu}[Gaz parfait et industrie au Cameroun]
L'usine de \textbf{SONARA} (Limbe) et les centrales thermiques (Kribi, Dibamba) manipulent des hydrocarbures gazeux (butane, propane, méthane). Pour dimensionner les réservoirs sphériques (stockage GPL) ou les tuyauteries, les ingénieurs utilisent l'équation d'état. Aux pressions de stockage (\SI{10}{à 20}{bar}), le gaz parfait donne une première estimation, mais la correction de Van der Waals (ou l'équation de Peng-Robinson) est indispensable pour la sécurité : une erreur de $5\,\%$ sur le volume molaire peut signifier une surpression dangereuse en cas d'incendie. Les coefficients $a$ et $b$ sont tabulés pour chaque gaz (ex. pour le méthane $\ce{CH4}$ : $a = \SI{0,228}{Pa\cdot m^6\cdot mol^{-2}}$, $b = \SI{4,28e-5}{m^3\cdot mol^{-1}}$).
\end{savaistu}

\courssub{Contraintes d'une loi : domaine de validité et limites des modèles}

\begin{propriete}[Domaine de validité d'une loi physique]
Toute loi physique est établie dans un cadre précis : \textbf{hypothèses de base}, \textbf{gamme de variation des grandeurs} (température, pression, champ, fréquence...), \textbf{précision visée}. En dehors de ce domaine, la loi peut donner des résultats faux, voire aberrants (ex. : loi d'Ohm pour une diode, gaz parfait à \SI{1000}{bar}, mécanique newtonienne à $v \approx c$).
\end{propriete}

\begin{aretenir}[Critères de validité d'un modèle]
\begin{itemize}
\item \textbf{Cohérence interne} : homogénéité dimensionnelle, respect des principes (conservation énergie, charge...).
\item \textbf{Accord expérimental} : écart relatif compatible avec les incertitudes de mesure.
\item \textbf{Prédictibilité} : le modèle anticipe correctement de nouvelles mesures.
\item \textbf{Simplicité (Rasoir d'Ockham)} : entre deux modèles équivalents, on retient le plus simple.
\end{itemize}
\end{aretenir}

\begin{attention}[Pièges fréquents au Probatoire]
\begin{itemize}
\item Confondre \textbf{loi} (relation universelle, ex. $F=ma$) et \textbf{modèle} (application contextualisée avec hypothèses, ex. « le mobile est un point matériel sans frottement »).
\item Oublier de \textbf{préciser les unités} des paramètres déterminés graphiquement (pente en $\Omega$, ordonnée en V).
\item Ne pas \textbf{tester le modèle} sur un point indépendant (se contenter de dire « les points sont alignés »).
\item Dire « le modèle est faux » au lieu de « le modèle n'est plus valide car l'hypothèse X n'est plus vérifiée » (ex. : $r$ constant vs échauffement).
\item Utiliser une loi hors de son domaine (ex. : $PV=nRT$ pour de l'eau liquide).
\end{itemize}
\end{attention}

\courssub{Synthèse : Travailler comme un scientifique}

\begin{aretenir}[Compétences clés du Module 1 (Mesures et incertitudes)]
\begin{tabularx}{\linewidth}{|l|X|}\hline
\rowcolor{popblueL} \textbf{Savoir-faire} & \textbf{Illustration dans cette leçon} \\ \hline
Formuler des hypothèses simplificatrices & Source idéale $U_0$, résistance unique $r$, isothermie. \\ \hline
Choisir un modèle adapté & Modèle série $U_0, r$ pour une ligne électrique ; Gaz parfait $PV=nRT$. \\ \hline
Établir une loi à partir de mesures & Tracé $U=f(I)$, détection affine, détermination $U_0, r$. \\ \hline
Tester la validité par la mesure & Prédiction à \SI{40}{A}, écart relatif $0,09\,\%$. \\ \hline
Discuter le domaine de validité & Effet Joule $\rightarrow$ $r(T)$ $\rightarrow$ courbure de la droite. \\ \hline
Communiquer un résultat scientifique & Rédaction argumentée pour l'habitant sceptique. \\ \hline
\end{tabularx}
\end{aretenir}

\begin{exercice}[Entraînement : Modélisation d'une pile réelle]
On étudie une pile alcaline neuve. On mesure la tension $U$ à ses bornes pour différentes valeurs du courant $I$ débité (en variant la résistance de charge $R$).
\begin{center}
\begin{tabular}{|c|ccccc|}\hline
\rowcolor{popblueL} $I$ (mA) & 0 & 50 & 100 & 150 & 200 \\ \hline
$U$ (V) & 1,58 & 1,52 & 1,46 & 1,40 & 1,34 \\ \hline
\end{tabular}
\end{center}
\begin{enumerate}
\item Proposer un modèle simple pour cette pile (schéma + hypothèses).
\item Tracer $U = f(I)$. Déterminer la f.é.m. $E$ et la résistance interne $r$.
\item Prédire $U$ pour $I = \SI{250}{mA}$. La mesure donne \SI{1,27}{V}. Calculer l'écart relatif. Le modèle est-il encore valide ?
\item Expliquer physiquement pourquoi l'écart apparaît (hypothèse remise en cause).
\end{enumerate}
\end{exercice}

\begin{corrige}[Exercice : Modélisation d'une pile réelle]
\begin{enumerate}
\item \textbf{Modèle :} Pile = source de tension idéale $E$ (f.é.m.) en série avec une résistance interne $r$. Hypothèses : $E$ et $r$ constants (pas de polarisation, pas d'échauffement notable, régime stationnaire).
\item \textbf{Graphique :} Points alignés. Droite $U = E - rI$. $E = \SI{1,58}{V}$ (ordonnée à l'origine). Pente $r = -\frac{1,34 - 1,58}{0,200 - 0} = \frac{0,24}{0,200} = \SI{1,2}{\Omega}$.
\item \textbf{Prédiction :} $U_{\text{prév}} = 1,58 - 1,2 \times 0,250 = \SI{1,28}{V}$. Mesure $U_{\text{mes}} = \SI{1,27}{V}$. $\varepsilon = \frac{|1,28-1,27|}{1,27} \approx 0,8\,\%$. Modèle encore acceptable mais l'écart grandit.
\item \textbf{Interprétation :} À fort courant, la réaction chimique ne suit plus (polarisation de concentration) et l'échauffement interne augmente $r$. L'hypothèse « $r$ constant » est fragilisée. Le modèle affine est une approximation de premier ordre valable pour $I < \SI{200}{mA}$.
\end{enumerate}
\end{corrige}

\newpage

\coursec{Méthodes et exercices résolus}

\coursec{Qu'est-ce qu'un modèle en physique ?}

\begin{definition}[Modèle physique]
Un \cle{modèle} est une représentation simplifiée et idéalisée d'un système réel, construite pour en expliquer le comportement et en prévoir l'évolution. Il repose sur un ensemble d'\cle{hypothèses simplificatrices} qui délibérément ignorent les détails jugés secondaires pour ne garder que l'essentiel.
\end{definition}

\begin{exemplebox}[Du réel au modèle : le fil de cuivre]
\begin{itemize}
\item \textbf{Réel :} un fil de cuivre a une résistance qui dépend de la température, des impuretés, de la fréquence du courant, de sa géométrie exacte, etc.
\item \textbf{Modèle « conducteur ohmique » :} on suppose le fil à température constante, homogène, de section uniforme, et on ne retient qu'une seule grandeur : sa résistance $R$, constante. La loi $U = RI$ devient alors exacte \emph{dans le cadre du modèle}.
\end{itemize}
\end{exemplebox}

\begin{propriete}[Rôle des hypothèses simplificatrices]
Toute loi physique (loi d'Ohm, loi des gaz parfaits, lois de Newton...) n'est valable que \textbf{si les hypothèses qui la fondent sont respectées}. Ces hypothèses définissent le \cle{domaine de validité} de la loi. En dehors de ce domaine, la loi n'est plus applicable telle quelle.
\end{propriete}

\begin{savaistu}[Modélisation au Cameroun : le barrage de Lom Pangar]
La gestion du barrage de Lom Pangar (région de l'Est) repose sur des modèles hydrologiques complexes : modèles de précipitations, modèles d'écoulement dans la retenue, modèles de production turbine. Chaque modèle a son domaine de validité (ex. : modèle de crue centennale vs modèle de gestion courante). Les ingénieurs d'ENEO et du MINEE doivent connaître les limites de ces modèles pour garantir la sécurité du barrage et l'approvisionnement électrique de Yaoundé et du réseau interconnecté Sud (SONATREL).
\end{savaistu}

\coursec{Établir une loi à partir de mesures : la démarche expérimentale}

\begin{methode}[Démarche générale : de l'expérience à la loi]
\begin{enumerate}
\item \textbf{Questionner} : identifier les grandeurs susceptibles d'influencer le phénomène.
\item \textbf{Hypothétiser} : formuler une relation conjecturée (proportionnalité, proportionnalité inverse, relation affine, loi de puissance...).
\item \textbf{Expérimenter} : faire varier \textbf{une seule grandeur indépendante} à la fois, mesurer la grandeur dépendante correspondante. Répéter les mesures (au moins 5 à 6 points étalés sur le domaine).
\item \textbf{Analyser} : tracer le graphe, tester les allures (droite par l'origine, droite non passant par l'origine, hyperbole $\rightarrow$ tracer $y=f(1/x)$, etc.).
\item \textbf{Modéliser} : déterminer les paramètres (pente, ordonnée à l'origine) avec leurs unités par analyse dimensionnelle.
\item \textbf{Valider / Limiter} : vérifier la constance des rapports ou l'alignement aux incertitudes près ; expliciter le domaine de validité de la loi établie.
\end{enumerate}
\end{methode}

\coursec{Exemple 1 : la loi d'Ohm}

\courssub{Fiche méthode 1 : Formuler des hypothèses simplificatrices pour modéliser une situation}

\begin{methode}[Modéliser une situation concrète]
Pour transformer une situation réelle (souvent complexe) en un problème de physique traitable :
\begin{enumerate}
\item \textbf{Identifier le système étudié} et le délimiter clairement (exemple : le gaz contenu dans la seringue, et non la seringue entière).
\item \textbf{Repérer les grandeurs pertinentes} : celles que l'on peut mesurer ou qui interviennent dans le phénomène (pression, volume, température, tension, intensité...). Écarter les grandeurs dont l'influence est négligeable.
\item \textbf{Formuler des hypothèses simplificatrices explicites} : chaque hypothèse doit être écrite en une phrase claire. Exemples classiques :
\begin{itemize}
\item \og le gaz est supposé parfait \fg{} (molécules ponctuelles, sans interaction hors chocs) ;
\item \og les frottements sont négligés \fg{} ;
\item \og la température reste constante pendant l'expérience \fg{} ;
\item \og le fil de connexion a une résistance nulle \fg{}.
\end{itemize}
\item \textbf{Choisir le modèle adapté} : loi d'Ohm pour un conducteur ohmique, loi des gaz parfaits $PV=nRT$ pour un gaz dilué, etc.
\item \textbf{Annoncer le domaine de validité} : préciser dans quelles conditions les hypothèses sont raisonnables (faible pression, température proche de l'ambiante, tension pas trop élevée...).
\end{enumerate}
\textbf{Réflexe de rédaction} : au Probatoire, une hypothèse non écrite est une hypothèse perdue. Écris toujours : \og On modélise ... par ... en supposant que ... \fg{}.
\end{methode}

\begin{exoresolu}[Modéliser la détente de l'air dans une pompe à vélo]
Un élève de Première C à Ngaoundéré gonfle le pneu de sa moto avec une pompe à main. Il bouche la sortie de la pompe et enfonce lentement le piston : le volume d'air enfermé passe de $V_1 = \SI{300}{cm^3}$ à $V_2 = \SI{100}{cm^3}$, la pression initiale étant $P_1 = \SI{1,0e5}{Pa}$. On veut prévoir la pression finale $P_2$.
\begin{enumerate}
\item Formuler les hypothèses simplificatrices nécessaires pour modéliser cette situation.
\item Calculer $P_2$ dans le cadre de ce modèle.
\end{enumerate}
\tcblower
\textbf{1. Hypothèses simplificatrices.}

Le système étudié est \textbf{l'air enfermé dans le corps de la pompe} (quantité de matière $n$ constante, car la sortie est bouchée : pas de fuite).

Pour appliquer un modèle simple, on formule les hypothèses suivantes :
\begin{itemize}
\item \textbf{H1} : l'air est assimilé à un \cle{gaz parfait} (pression proche de la pression atmosphérique, température ordinaire : hypothèse raisonnable) ;
\item \textbf{H2} : la compression est \textbf{lente}, donc les échanges de chaleur avec l'extérieur maintiennent la température constante : transformation \cle{isotherme}, $T_1 = T_2$ ;
\item \textbf{H3} : le piston est étanche, donc $n$ est rigoureusement constant.
\end{itemize}

\textbf{2. Calcul de la pression finale.}

D'après la loi des gaz parfaits appliquée aux états 1 et 2 :
\[ P_1V_1 = nRT_1 \quad \text{et} \quad P_2V_2 = nRT_2. \]
Comme $n$ est constant et $T_1 = T_2$ (hypothèse H2), on obtient la \textbf{loi de Boyle-Mariotte} :
\[ P_1V_1 = P_2V_2 \quad \Longrightarrow \quad P_2 = P_1\,\dfrac{V_1}{V_2}. \]
Application numérique (les volumes étant dans le rapport, on peut garder les \si{cm^3}) :
\[ P_2 = \num{1,0e5} \times \dfrac{300}{100} = \SI{3,0e5}{Pa}. \]

\textbf{Conclusion} : en modélisant l'air comme un gaz parfait subissant une compression isotherme, la pression finale est $P_2 = \SI{3,0e5}{Pa}$, soit environ 3 fois la pression atmosphérique. Si l'on enfonçait le piston brutalement, l'hypothèse H2 serait fausse (l'air chauffe) et la pression réelle serait supérieure à cette valeur.
\end{exoresolu}

\courssub{Fiche méthode 2 : Tester la validité des hypothèses par la mesure}

\begin{methode}[Valider ou invalider un modèle expérimentalement]
Une hypothèse n'a de valeur que si l'expérience la confirme. Démarche :
\begin{enumerate}
\item \textbf{Prédire} : à partir du modèle, écrire la relation attendue entre les grandeurs (exemple : $U = RI$, donc $U$ proportionnelle à $I$).
\item \textbf{Mesurer} : réaliser une série de mesures $(x_i, y_i)$ en faisant varier une seule grandeur à la fois, les autres étant maintenues constantes.
\item \textbf{Tracer} le graphe $y = f(x)$ et observer son allure.
\item \textbf{Comparer à la prédiction} :
\begin{itemize}
\item si le modèle prévoit une proportionnalité, vérifier que les points sont \textbf{alignés sur une droite passant par l'origine} ;
\item calculer le rapport $y_i/x_i$ pour chaque mesure : il doit être \textbf{constant aux incertitudes près}.
\end{itemize}
\item \textbf{Conclure honnêtement} : préciser le \textbf{domaine de mesures} où le modèle est validé, et signaler les écarts éventuels (points qui s'écartent de la droite aux fortes valeurs, par exemple).
\end{enumerate}
\textbf{Réflexe} : un seul couple de mesures ne valide jamais une loi ; il faut au minimum 5 à 6 points étalés sur tout le domaine étudié.
\end{methode}

\begin{exoresolu}[Tester l'hypothèse \og conducteur ohmique \fg{}]
Au laboratoire du lycée, on mesure la tension $U$ aux bornes d'un dipôle inconnu pour différentes valeurs de l'intensité $I$ :

\begin{center}
\begin{tabular}{|c|cccccc|}
\hline
$I$ (mA) & 0 & 20 & 40 & 60 & 80 & 100 \\
\hline
$U$ (V) & 0 & 0,44 & 0,90 & 1,31 & 1,78 & 2,21 \\
\hline
\end{tabular}
\end{center}

\begin{enumerate}
\item Si le dipôle est un conducteur ohmique, quelle relation doit vérifier le couple $(U, I)$ ?
\item Tester cette hypothèse à l'aide des mesures et conclure.
\end{enumerate}
\tcblower
\textbf{1. Relation attendue.}

Si le dipôle est un conducteur ohmique, il vérifie la \textbf{loi d'Ohm} :
\[ U = R\,I, \]
c'est-à-dire que $U$ est \textbf{proportionnelle} à $I$ : le rapport $U/I$ doit être constant et égal à la résistance $R$, et la caractéristique $U = f(I)$ doit être une droite passant par l'origine.

\textbf{2. Test de l'hypothèse.}

Calculons le rapport $U/I$ pour chaque mesure (avec $I$ en ampères pour obtenir $R$ en ohms) :

\begin{center}
\begin{tabular}{|c|ccccc|}
\hline
$I$ (A) & 0,020 & 0,040 & 0,060 & 0,080 & 0,100 \\
\hline
$U/I$ ($\Omega$) & $\dfrac{0,44}{0,020}=22,0$ & $\dfrac{0,90}{0,040}=22,5$ & $\dfrac{1,31}{0,060}=21,8$ & $\dfrac{1,78}{0,080}=22,3$ & $\dfrac{2,21}{0,100}=22,1$ \\
\hline
\end{tabular}
\end{center}

Les valeurs du rapport $U/I$ sont toutes voisines de $\SI{22}{\Omega}$ : l'écart maximal est
\[ \dfrac{22,5 - 21,8}{22,1} \approx 0,03 \quad \text{soit environ } 3\,\%, \]
ce qui est compatible avec les incertitudes de lecture des appareils.

De plus, la mesure $(I=0 \Rightarrow U=0)$ montre que la caractéristique passe par l'origine.

\textbf{Conclusion} : le rapport $U/I$ est constant aux incertitudes expérimentales près : l'hypothèse \og conducteur ohmique \fg{} est \textbf{validée} sur le domaine de mesures $0 \leq I \leq \SI{100}{mA}$, avec $R \approx \SI{22}{\Omega}$. Attention : ce test ne dit rien au-delà de $\SI{100}{mA}$ ; pour des intensités plus grandes, l'échauffement du dipôle pourrait modifier sa résistance et invalider le modèle.
\end{exoresolu}

\courssub{Fiche méthode 4 : Exploiter la loi d'Ohm dans le cadre de son domaine de validité}

\begin{methode}[Appliquer correctement la loi d'Ohm]
\begin{enumerate}
\item \textbf{Vérifier les conditions d'application} : le dipôle est un conducteur ohmique (résistor), maintenu à \textbf{température pratiquement constante} (intensité modérée pour éviter l'effet Joule excessif).
\item \textbf{Écrire la loi} : $U = R\,I$, avec $U$ en volts (V), $I$ en ampères (A), $R$ en ohms ($\Omega$).
\item \textbf{Convertir les unités avant tout calcul} : mA $\rightarrow$ A ($\div 1000$), k$\Omega$ $\rightarrow$ $\Omega$ ($\times 1000$).
\item \textbf{Isoler la grandeur cherchée} littéralement : $R = \dfrac{U}{I}$ ou $I = \dfrac{U}{R}$.
\item \textbf{Faire l'application numérique} avec les unités SI, puis exprimer le résultat avec un nombre raisonnable de chiffres significatifs et son unité.
\item \textbf{Vérifier l'ordre de grandeur} : pour un résistor de laboratoire (quelques dizaines d'ohms) sous quelques volts, on attend des intensités de l'ordre de $0,01$ à $\SI{0,5}{A}$.
\end{enumerate}
\textbf{Analyse dimensionnelle} : $[R] = \dfrac{[U]}{[I]} = \dfrac{\si{V}}{\si{A}} = \si{\Omega}$. \checkmark
\end{methode}

\begin{exoresolu}[Installation électrique dans un quartier de Douala]
Pour alimenter un atelier, un technicien utilise un enrouleur de fil de cuivre dont la résistance totale est $R = \SI{2,5}{\Omega}$. L'atelier est alimenté sous une tension efficace que l'on modélise ici par une tension continue $U = \SI{220}{V}$ aux bornes de l'ensemble \{fil + appareil\}. Dans une première question, on s'intéresse au seul enrouleur.
\begin{enumerate}
\item Calculer l'intensité $I$ qui traverserait l'enrouleur s'il était seul branché sous $\SI{220}{V}$.
\item Cette situation est-elle réaliste pour un fil restant \og conducteur ohmique \fg{} ? Commenter en évoquant le domaine de validité de la loi d'Ohm.
\end{enumerate}
On donne : la puissance dissipée par effet Joule dans un résistor est $P = RI^2$.
\tcblower
\textbf{1. Calcul de l'intensité.}

L'enrouleur est un conducteur ohmique ; la loi d'Ohm s'écrit $U = RI$, d'où :
\[ I = \dfrac{U}{R} = \dfrac{220}{2,5} = \SI{88}{A}. \]

\textbf{2. Commentaire sur le domaine de validité.}

La puissance dissipée par effet Joule dans le fil serait :
\[ P = RI^2 = 2,5 \times 88^2 = 2,5 \times 7744 \approx \SI{1,9e4}{W} \approx \SI{19}{kW}. \]
Une telle puissance porterait le cuivre à très haute température en quelques secondes : le fil fondrait (c'est le principe du fusible). Or la loi d'Ohm $U = RI$ avec $R$ \textbf{constante} n'est valable que si la température du conducteur reste pratiquement constante ; à forte température, la résistivité du cuivre augmente, donc $R$ augmente et la relation $U = RI$ avec $R = \SI{2,5}{\Omega}$ cesse d'être exacte.

\textbf{Conclusion} : le calcul formel donne $I = \SI{88}{A}$, mais cette valeur n'est pas physiquement atteignable durablement : le modèle \og conducteur ohmique à $R$ constante \fg{} sort de son domaine de validité bien avant, car l'échauffement modifie la résistance. C'est pourquoi les installations réelles sont protégées par des disjoncteurs qui coupent le circuit dès que l'intensité dépasse une valeur de sécurité (par exemple $\SI{16}{A}$).
\end{exoresolu}

\coursec{Exemple 2 : la loi des gaz parfaits}

\courssub{Fiche méthode 3 : Établir une loi mathématique simple entre des grandeurs mesurées}

\begin{methode}[Dégager une loi à partir d'un tableau de mesures]
\begin{enumerate}
\item \textbf{Tracer le graphe} $y = f(x)$ avec des axes gradués, nommés et munis d'unités.
\item \textbf{Identifier l'allure} :
\begin{itemize}
\item points alignés sur une droite passant par l'origine $\Rightarrow$ \textbf{proportionnalité} : $y = kx$ ;
\item points alignés sur une droite ne passant pas par l'origine $\Rightarrow$ \textbf{relation affine} : $y = ax + b$ ;
\item courbe d'allure hyperbolique $\Rightarrow$ tester l'\textbf{inverse proportionnalité} : tracer $y$ en fonction de $1/x$ ; si les points s'alignent sur une droite passant par l'origine, alors $y = k/x$.
\end{itemize}
\item \textbf{Déterminer les paramètres} :
\begin{itemize}
\item pente $a = \dfrac{y_2 - y_1}{x_2 - x_1}$ en prenant deux points éloignés de la droite moyenne (pas forcément des points de mesure) ;
\item ordonnée à l'origine $b$ lue graphiquement ou calculée : $b = y_1 - a x_1$.
\end{itemize}
\item \textbf{Donner l'unité des paramètres} par analyse dimensionnelle : $[a] = [y]/[x]$.
\item \textbf{Écrire la loi} sous forme littérale puis numérique avec les unités, et préciser le domaine de mesures utilisé.
\end{enumerate}
\end{methode}

\begin{exoresolu}[Établir la loi de Boyle-Mariotte]
Un groupe d'élèves comprime de l'air dans une seringue étanche reliée à un manomètre, à température constante. Résultats :

\begin{center}
\begin{tabular}{|c|ccccc|}
\hline
$V$ (cm\textsuperscript{3}) & 60 & 40 & 30 & 24 & 20 \\
\hline
$P$ (kPa) & 100 & 150 & 200 & 250 & 300 \\
\hline
\end{tabular}
\end{center}

\begin{enumerate}
\item Montrer que $P$ n'est pas proportionnelle à $V$.
\item Établir la loi mathématique reliant $P$ et $V$.
\end{enumerate}
\tcblower
\textbf{1. Test de proportionnalité directe.}

Si $P$ était proportionnelle à $V$, le rapport $P/V$ serait constant. Or :
\[ \dfrac{P_1}{V_1} = \dfrac{100}{60} \approx 1,67 \ \si{kPa/cm^3} \quad \text{et} \quad \dfrac{P_2}{V_2} = \dfrac{150}{40} = 3,75 \ \si{kPa/cm^3}. \]
Le rapport n'est pas constant : $P$ \textbf{n'est pas proportionnelle} à $V$. De plus, $P$ augmente quand $V$ diminue : on soupçonne une \textbf{inverse proportionnalité}.

\textbf{2. Recherche de la loi.}

Testons l'hypothèse $P = \dfrac{k}{V}$, c'est-à-dire $PV = k$ constant. Calculons le produit $PV$ :

\begin{center}
\begin{tabular}{|c|ccccc|}
\hline
$V$ (cm\textsuperscript{3}) & 60 & 40 & 30 & 24 & 20 \\
\hline
$PV$ (kPa$\cdot$cm\textsuperscript{3}) & 6000 & 6000 & 6000 & 6000 & 6000 \\
\hline
\end{tabular}
\end{center}

Le produit $PV$ est \textbf{constant} : $PV = \SI{6000}{kPa.cm^3}$.

Vérification graphique équivalente : en traçant $P$ en fonction de $1/V$, les points s'alignent sur une droite passant par l'origine, de pente
\[ k = \dfrac{\Delta P}{\Delta(1/V)} = \SI{6000}{kPa.cm^3}. \]

\textbf{Conclusion} : la loi établie à partir des mesures est
\[ P \times V = \SI{6,0e3}{kPa.cm^3} \quad \text{ou} \quad P = \dfrac{\num{6,0e3}}{V} \ \text{(}P\text{ en kPa, }V\text{ en cm}^3\text{)}. \]
C'est la \textbf{loi de Boyle-Mariotte} : à température constante et quantité de gaz constante, le produit $PV$ est constant. Cette loi n'est valable que dans le domaine exploré (et plus généralement pour un gaz à faible pression, proche du gaz parfait).
\end{exoresolu}

\courssub{Fiche méthode 5 : Exploiter la loi des gaz parfaits}

\begin{methode}[Appliquer la loi des gaz parfaits $PV = nRT$]
\begin{enumerate}
\item \textbf{Vérifier les conditions} : gaz à \textbf{faible pression} (typiquement $P$ de l'ordre de la pression atmosphérique ou inférieure) et température éloignée de la liquéfaction du gaz.
\item \textbf{Convertir TOUTES les grandeurs en unités SI} avant le calcul :
\begin{itemize}
\item $P$ en pascals (Pa) : $P(\si{Pa}) = P(\si{bar}) \times 10^5$ ;
\item $V$ en mètres cubes : $V(\si{m^3}) = V(\si{L}) \times 10^{-3}$ ;
\item $T$ en kelvins : $T(\si{K}) = \theta(\si{\celsius}) + 273,15$ ;
\item $n$ en moles, $R = \SI{8,314}{J.K^{-1}.mol^{-1}}$.
\end{itemize}
\item \textbf{Isoler la grandeur inconnue} : $n = \dfrac{PV}{RT}$, $V = \dfrac{nRT}{P}$, etc.
\item \textbf{Pour une transformation à $n$ constant}, écrire les états initial et final :
\[ \dfrac{P_1V_1}{T_1} = \dfrac{P_2V_2}{T_2}, \]
puis simplifier selon la transformation (isotherme : $T_1 = T_2$ ; isobare : $P_1 = P_2$ ; isochore : $V_1 = V_2$).
\item \textbf{Vérifier le résultat} : analyse dimensionnelle et ordre de grandeur (une mole de gaz parfait occupe environ $\SI{24}{L}$ à $\SI{20}{\celsius}$ sous $\SI{1}{bar}$).
\end{enumerate}
\textbf{Analyse dimensionnelle} : $[PV] = \si{Pa.m^3} = \si{N.m^{-2}.m^3} = \si{N.m} = \si{J}$ et $[nRT] = \si{mol} \times \si{J.K^{-1}.mol^{-1}} \times \si{K} = \si{J}$. \checkmark
\end{methode}

\begin{exoresolu}[Bouteille d'oxygène à l'hôpital de Laquintinie]
À l'hôpital Laquintinie de Douala, une bouteille d'oxygène médical de volume $V = \SI{10}{L}$ contient du dioxygène sous une pression $P = \SI{150}{bar}$ à la température $\theta = \SI{27}{\celsius}$.
\begin{enumerate}
\item Calculer la quantité de matière $n$ de dioxygène contenue dans la bouteille, en assimilant le dioxygène à un gaz parfait.
\item Quel volume ce gaz occuperait-il à la pression atmosphérique $P_0 = \SI{1,0}{bar}$, à la même température ?
\item La pression de $\SI{150}{bar}$ est élevée : que peut-on dire de la validité du modèle du gaz parfait dans cette bouteille ?
\end{enumerate}
On donne : $R = \SI{8,314}{J.K^{-1}.mol^{-1}}$.
\tcblower
\textbf{1. Quantité de matière.}

Conversions en unités SI :
\[ P = 150 \times 10^5 = \SI{1,50e7}{Pa}, \quad V = 10 \times 10^{-3} = \SI{1,0e-2}{m^3}, \quad T = 27 + 273,15 = \SI{300,15}{K} \approx \SI{300}{K}. \]
D'après la loi des gaz parfaits :
\[ n = \dfrac{PV}{RT} = \dfrac{\num{1,50e7} \times \num{1,0e-2}}{8,314 \times 300,15} = \dfrac{\num{1,50e5}}{\num{2495,2}} \approx \SI{60,1}{mol} \approx \SI{60}{mol}. \]

\textbf{2. Volume à pression atmosphérique.}

À $n$ et $T$ constants (transformation isotherme), la loi de Boyle-Mariotte donne :
\[ P_1V_1 = P_2V_2 \quad \Longrightarrow \quad V_2 = V_1\,\dfrac{P_1}{P_2} = 10 \times \dfrac{150}{1,0} = \SI{1500}{L} = \SI{1,5}{m^3}. \]
(Le rapport des pressions étant sans dimension, on peut garder les pressions en bar et le volume en litres.)

\textbf{3. Validité du modèle.}

Le modèle du gaz parfait suppose des molécules ponctuelles sans interactions hors chocs : c'est une bonne approximation à \textbf{faible pression}. À $\SI{150}{bar}$, les molécules sont proches les unes des autres : leur volume propre n'est plus négligeable et des interactions attractives apparaissent. Le modèle du gaz parfait devient donc \textbf{approximatif} : la valeur $n \approx \SI{60}{mol}$ est une estimation légèrement entachée d'erreur (un modèle de gaz réel, comme celui de Van der Waals, serait plus adapté).

\textbf{Conclusion} : la bouteille contient environ $n \approx \SI{60}{mol}$ de dioxygène, qui occuperait environ $\SI{1,5}{m^3}$ à pression atmosphérique — d'où l'intérêt de la compression pour le stockage — mais ce résultat repose sur un modèle poussé ici aux limites de son domaine de validité.
\end{exoresolu}

\coursec{Contraintes d'une loi : domaine de validité et limites des modèles}

\courssub{Fiche méthode 6 : Identifier et exploiter le domaine de validité d'une loi}

\begin{methode}[Raisonner sur les contraintes d'une loi]
\begin{enumerate}
\item \textbf{Lister les hypothèses} sur lesquelles repose la loi (gaz parfait : molécules ponctuelles sans interaction ; loi d'Ohm : température constante ; loi de Boyle-Mariotte : $n$ et $T$ constants...).
\item \textbf{Traduire chaque hypothèse en contrainte concrète} : pression faible, intensité modérée, transformation lente, système fermé...
\item \textbf{Comparer la situation proposée aux contraintes} : si une condition est violée (pression de $\SI{150}{bar}$, filament porté à $\SI{2500}{\celsius}$, compression brutale...), la loi n'est plus applicable telle quelle.
\item \textbf{Prévoir qualitativement l'écart} : dans quel sens le comportement réel s'écarte-t-il du modèle ? (Exemple : la résistance d'un filament augmente avec la température, donc $I$ croît moins vite que $U/R$.)
\item \textbf{Conclure en deux temps} : \og dans le domaine ..., la loi ... s'applique car ... ; en dehors, ... \fg{}.
\end{enumerate}
\textbf{Réflexe Probatoire} : citer un domaine de validité dans la conclusion d'un exercice de modélisation est systématiquement valorisé.
\end{methode}

\begin{exoresolu}[La lampe à incandescence : un faux conducteur ohmique]
On relève la caractéristique d'une lampe à filament (lampe de poche) :

\begin{center}
\begin{tabular}{|c|cccc|}
\hline
$U$ (V) & 0 & 1,0 & 2,0 & 3,0 \\
\hline
$I$ (mA) & 0 & 120 & 190 & 240 \\
\hline
\end{tabular}
\end{center}

\begin{enumerate}
\item Calculer le rapport $U/I$ pour chaque point de fonctionnement. La lampe vérifie-t-elle la loi d'Ohm avec une résistance constante ?
\item Identifier l'hypothèse de la loi d'Ohm qui n'est plus respectée et expliquer physiquement l'écart observé.
\end{enumerate}
\tcblower
\textbf{1. Calcul des rapports $U/I$.}

\begin{center}
\begin{tabular}{|c|ccc|}
\hline
$U$ (V) & 1,0 & 2,0 & 3,0 \\
\hline
$U/I$ ($\Omega$) & $\dfrac{1,0}{0,120} \approx 8,3$ & $\dfrac{2,0}{0,190} \approx 10,5$ & $\dfrac{3,0}{0,240} = 12,5$ \\
\hline
\end{tabular}
\end{center}

Le rapport $U/I$ \textbf{n'est pas constant} : il passe de $\SI{8,3}{\Omega}$ à $\SI{12,5}{\Omega}$, soit une augmentation de $50\,\%$. La lampe à filament \textbf{ne vérifie pas} la loi d'Ohm avec une résistance constante : ce n'est pas un conducteur ohmique sur ce domaine de fonctionnement.

\textbf{2. Hypothèse en cause et interprétation.}

La loi d'Ohm $U = RI$ avec $R$ constante suppose que le conducteur reste à \textbf{température constante}. Or, quand la tension augmente, l'intensité augmente et l'effet Joule ($P = UI$) porte le filament de tungstène à plusieurs centaines de degrés : il devient incandescent (c'est son rôle !). La résistivité des métaux \textbf{augmente avec la température}, car l'agitation thermique des atomes du réseau freine davantage le déplacement des électrons. La résistance du filament augmente donc avec $U$, ce qui explique que $U/I$ croisse : l'intensité croît \textbf{moins vite} que ne le prévoirait la loi d'Ohm à $R$ constante.

\textbf{Conclusion} : la lampe à filament illustre les contraintes d'une loi : la loi d'Ohm à résistance constante n'est applicable que si la température du conducteur varie peu. Pour des tensions faibles (filament tiède), le modèle reste approximativement valable ; à tension nominale, il faut considérer une résistance dépendant de la température.
\end{exoresolu}

\coursec{Synthèse : travailler comme un scientifique}

\begin{aretenir}[Bilan des compétences : Modéliser un phénomène]
\begin{itemize}
\item \textbf{Modéliser} : identifier le système, choisir les grandeurs, écrire les hypothèses simplificatrices, choisir la loi adaptée, annoncer le domaine de validité.
\item \textbf{Établir une loi} : mesurer (5-6 points minimum), tracer, identifier l'allure (proportionnalité, affine, inverse), déterminer les paramètres (pente, origine) avec unités, écrire la loi littérale et numérique.
\item \textbf{Tester une loi} : prédire la relation, comparer mesures/prédiction (alignement, constance des rapports), conclure sur le domaine de validité.
\item \textbf{Appliquer une loi} : vérifier les conditions (température, pression, lenteur...), convertir en SI, calculer, vérifier l'ordre de grandeur, discuter la validité \emph{a posteriori}.
\item \textbf{Rédiger} : phrases complètes, unités systématiques, chiffres significatifs cohérents, conclusion synthétique mentionnant le domaine de validité.
\end{itemize}
\end{aretenir}

\begin{attention}[Les 5 erreurs qui coûtent des points au Probatoire]
\begin{enumerate}
\item \textbf{Oublier de convertir la température en kelvins} dans la loi des gaz parfaits : $T(\si{K}) = \theta(\si{\celsius}) + 273,15$. Utiliser $\SI{27}{\celsius}$ à la place de $\SI{300}{K}$ fausse tout le calcul et est sanctionné comme erreur de physique, pas seulement de calcul.
\item \textbf{Appliquer une loi sans vérifier ni citer ses conditions} : écrire $PV = nRT$ pour un gaz très comprimé, ou $U = RI$ pour une lampe à filament, sans discuter le domaine de validité. Le correcteur attend la mention des hypothèses (\og gaz supposé parfait \fg{}, \og température constante \fg{}).
\item \textbf{Confondre proportionnalité et relation affine} : une droite qui ne passe pas par l'origine ne traduit pas une proportionnalité. Avant d'écrire $y = kx$, vérifier que le point $(0\,;0)$ appartient à la droite moyenne.
\item \textbf{Calculer une pente avec des unités mélangées} : diviser des volts par des milliampères et donner le résultat en ohms sans conversion ($\si{V}/\si{mA} = \si{k\Omega}$, pas $\si{\Omega}$ !). Toujours convertir en unités SI avant l'application numérique.
\item \textbf{Conclure sans phrase de synthèse ni unité} : un résultat brut \og $n = 60$ \fg{} sans unité, sans chiffres significatifs cohérents et sans phrase de conclusion (\og la bouteille contient environ $\SI{60}{mol}$ de dioxygène \fg{}) fait perdre des points de rédaction, même si le calcul est juste.
\end{enumerate}
\end{attention}

\newpage

\coursec{Exercices}

\courssub{Je vérifie mes connaissances}

\begin{exercice}[QCM : modèles et lois]
Pour chaque question, une seule réponse est correcte. Entoure la bonne lettre.
\begin{enumerate}
\item Un modèle en physique est :
\begin{enumerate}
\item une représentation exacte de la réalité ;
\item une représentation simplifiée de la réalité, fondée sur des hypothèses, qui permet de faire des prévisions ;
\item une maquette à l'échelle d'un phénomène.
\end{enumerate}
\item La loi d'Ohm $U = R\,I$ s'applique :
\begin{enumerate}
\item à tout dipôle électrique ;
\item uniquement aux dipôles ohmiques (conducteurs ohmiques) ;
\item uniquement aux générateurs.
\end{enumerate}
\item Dans la loi des gaz parfaits $PV = nRT$, la température $T$ doit être exprimée en :
\begin{enumerate}
\item degrés Celsius ;
\item degrés Fahrenheit ;
\item kelvins.
\end{enumerate}
\item Une loi physique possède toujours :
\begin{enumerate}
\item un domaine de validité ;
\item une valeur infinie ;
\item aucune hypothèse.
\end{enumerate}
\end{enumerate}
\end{exercice}

\begin{exercice}[Vrai ou faux justifié]
Pour chaque affirmation, indique si elle est vraie ou fausse, puis justifie ta réponse en une ou deux phrases.
\begin{enumerate}
\item « La loi des gaz parfaits reste valable quelle que soit la pression du gaz. »
\item « Un modèle dont les prévisions ne sont pas confirmées par la mesure doit être modifié ou abandonné. »
\item « La caractéristique $U = f(I)$ d'un conducteur ohmique est une droite passant par l'origine. »
\item « Les hypothèses simplificatrices d'un modèle n'ont aucune importance pour la qualité des prévisions. »
\end{enumerate}
\end{exercice}

\begin{exercice}[Définitions et vocabulaire]
\begin{enumerate}
\item Définis les termes suivants : \cle{modèle}, \cle{hypothèse simplificatrice}, \cle{domaine de validité}.
\item Cite deux hypothèses du modèle du gaz parfait concernant les molécules.
\item Donne l'expression de la loi des gaz parfaits en précisant le nom et l'unité SI de chaque grandeur.
\end{enumerate}
\end{exercice}

\begin{exercice}[Calculs directs]
\begin{enumerate}
\item Un conducteur ohmique de résistance $R = \SI{47}{\Omega}$ est traversé par un courant d'intensité $I = \SI{0,20}{A}$. Calcule la tension $U$ à ses bornes.
\item Une enceinte de volume $V = \SI{2,0}{L}$ contient $n = \num{0,10}$ mol de gaz parfait à la température $T = \SI{300}{K}$. Calcule la pression $P$ du gaz en pascals, puis en bars. On donne $R = \SI{8,314}{J.K^{-1}.mol^{-1}}$.
\item Convertis en kelvins : $t_1 = \SI{27}{\celsius}$ ; $t_2 = \SI{-23}{\celsius}$.
\end{enumerate}
\end{exercice}

\courssub{Je m'entraîne}

\begin{exercice}[Caractéristique d'un dipôle inconnu]
Au laboratoire du lycée, un élève mesure la tension $U$ aux bornes d'un dipôle inconnu pour différentes valeurs de l'intensité $I$ :
\begin{center}
\begin{tabular}{|l|*{6}{c|}}\hline
$I$ (mA) & 0 & 20 & 40 & 60 & 80 & 100 \\ \hline
$U$ (V) & 0 & 0,44 & 0,90 & 1,32 & 1,78 & 2,21 \\ \hline
\end{tabular}
\end{center}
\begin{enumerate}
\item Trace la caractéristique $U = f(I)$ sur papier millimétré (échelles : 1 cm pour 10 mA ; 1 cm pour 0,2 V).
\item Montre que ce dipôle est un conducteur ohmique.
\item Détermine graphiquement la valeur de la résistance $R$ de ce dipôle.
\item Calcule l'intensité du courant qui traverserait ce dipôle sous une tension de $\SI{12}{V}$. Commente la valeur obtenue au regard du domaine des mesures.
\end{enumerate}
\end{exercice}

\begin{exercice}[Bouteille de gaz domestique]
Une bouteille de gaz domestique de volume $V = \SI{12,5}{L}$ contient $n = \num{5,0}$ mol de butane, assimilé à un gaz parfait, à la température $t_1 = \SI{27}{\celsius}$. On donne $R = \SI{8,314}{J.K^{-1}.mol^{-1}}$.
\begin{enumerate}
\item Convertis les données en unités SI, puis calcule la pression $P_1$ du gaz dans la bouteille.
\item La bouteille est exposée en plein soleil : la température atteint $t_2 = \SI{47}{\celsius}$. En supposant le volume constant, calcule la nouvelle pression $P_2$.
\item Commente le danger d'une telle exposition. Cite deux hypothèses du modèle du gaz parfait utilisées dans ce calcul.
\end{enumerate}
\end{exercice}

\begin{exercice}[Vérification de la loi de Boyle-Mariotte]
À température constante, on mesure la pression $P$ d'une masse d'air enfermée dans une seringue pour différents volumes $V$ :
\begin{center}
\begin{tabular}{|l|*{5}{c|}}\hline
$V$ (mL) & 60 & 50 & 40 & 30 & 20 \\ \hline
$P$ (hPa) & 500 & 600 & 750 & 1000 & 1500 \\ \hline
\end{tabular}
\end{center}
\begin{enumerate}
\item Complète le tableau en calculant $\dfrac{1}{V}$ (en mL$^{-1}$) et le produit $P \times V$ pour chaque mesure.
\item Trace la courbe $P = f\left(\dfrac{1}{V}\right)$. Quelle relation mathématique simple peut-on établir entre $P$ et $V$ ?
\item Énonce la loi de Boyle-Mariotte et précise ses conditions d'application.
\item Calcule le volume qu'occuperait cet air sous une pression de $\SI{2000}{hPa}$.
\item La loi reste-t-elle valable si la pression devient très élevée ? Justifie en évoquant les hypothèses du modèle.
\end{enumerate}
\end{exercice}

\begin{exercice}[Lampe à filament : la loi d'Ohm en défaut]
On mesure la tension $U$ aux bornes d'une lampe à filament pour plusieurs intensités $I$ :
\begin{center}
\begin{tabular}{|l|*{5}{c|}}\hline
$I$ (A) & 0,10 & 0,20 & 0,30 & 0,40 & 0,50 \\ \hline
$U$ (V) & 0,50 & 1,60 & 3,30 & 5,60 & 8,50 \\ \hline
\end{tabular}
\end{center}
\begin{enumerate}
\item Calcule le rapport $\dfrac{U}{I}$ pour chaque mesure. Ce rapport est-il constant ?
\item La loi d'Ohm s'applique-t-elle à cette lampe ? Justifie.
\item Explique pourquoi le filament ne vérifie pas la loi d'Ohm, en identifiant l'hypothèse du modèle du conducteur ohmique qui n'est pas respectée (on rappelle que la résistance d'un métal augmente avec la température).
\end{enumerate}
\end{exercice}

\courssub{Je me prépare au Probatoire}

\begin{exercice}[Étude expérimentale d'un dipôle ohmique]
Un professeur de physique du lycée de Ngaoundéré remet à ses élèves de Première C un dipôle $D$ de nature inconnue et leur demande de le modéliser. Les élèves réalisent le montage comportant un générateur de tension réglable, un ampèremètre en série et un voltmètre aux bornes de $D$. Ils obtiennent les mesures suivantes :
\begin{center}
\begin{tabular}{|l|*{7}{c|}}\hline
$U$ (V) & 0 & 1,5 & 3,0 & 4,5 & 6,0 & 7,5 & 9,0 \\ \hline
$I$ (mA) & 0 & 15 & 30 & 45 & 60 & 75 & 90 \\ \hline
\end{tabular}
\end{center}
\begin{enumerate}
\item Dessine le schéma du circuit électrique utilisé, en plaçant correctement l'ampèremètre et le voltmètre.
\item Trace la caractéristique $U = f(I)$ du dipôle (échelles : 1 cm pour 10 mA ; 1 cm pour 1 V).
\item À l'aide du graphe, montre que le dipôle $D$ est un conducteur ohmique. Écris la loi mathématique reliant $U$ et $I$.
\item Détermine graphiquement la résistance $R$ du dipôle, en précisant la méthode utilisée.
\item Calcule l'intensité du courant qui traverserait ce dipôle s'il était soumis à une tension de $\SI{12}{V}$.
\item En réalité, le dipôle chauffe lorsqu'il est parcouru par un courant intense. Explique pourquoi la prévision de la question 5 pourrait s'écarter de la mesure réelle, et précise le domaine de validité du modèle utilisé.
\end{enumerate}
\end{exercice}

\begin{exercice}[Gaz parfait et sécurité domestique]
Dans une cuisine de Yaoundé, une bouteille de gaz domestique de volume intérieur $V = \SI{12,5}{L}$ contient $n = \num{5,0}$ mol de butane. On assimilera le butane à un gaz parfait. On donne $R = \SI{8,314}{J.K^{-1}.mol^{-1}}$.
\begin{enumerate}
\item Rappelle l'expression de la loi des gaz parfaits en nommant chaque grandeur et en précisant son unité SI.
\item Cite deux hypothèses simplificatrices du modèle du gaz parfait.
\item La température de la cuisine est $t_1 = \SI{27}{\celsius}$. Calcule la pression $P_1$ du gaz dans la bouteille, en pascals puis en bars ($\SI{1}{bar} = \SI{e5}{Pa}$).
\item Par négligence, la bouteille est laissée près du feu : la température du gaz atteint $t_2 = \SI{47}{\celsius}$. Le volume de la bouteille étant supposé constant, montre que $\dfrac{P_1}{T_1} = \dfrac{P_2}{T_2}$, puis calcule la pression $P_2$.
\item La bouteille est conçue pour résister à une pression maximale de $\SI{1,5e6}{Pa}$. Y a-t-il danger dans la situation de la question 4 ? Justifie par un calcul.
\item En réalité, à forte pression, le butane se liquéfie partiellement et ne se comporte plus comme un gaz parfait. Explique, en t'appuyant sur les hypothèses du modèle, pourquoi la loi des gaz parfaits atteint alors ses limites.
\end{enumerate}
\end{exercice}

\begin{exercice}[Détermination expérimentale de la constante $R$ et discussion sur les modèles]
Lors d'une séance de travaux pratiques, un groupe d'élèves mesure la pression $P$, le volume $V$ et la température $T$ d'une quantité $n = \num{0,050}$ mol d'air enfermé dans une enceinte étanche. Résultats : $P = \SI{1,02e5}{Pa}$ ; $V = \SI{1,24}{L}$ ; $t = \SI{25}{\celsius}$.
\begin{enumerate}
\item Convertis le volume et la température en unités SI.
\item En assimilant l'air à un gaz parfait, déduis des mesures une valeur expérimentale $R_{\text{exp}}$ de la constante des gaz parfaits.
\item Calcule l'écart relatif $\dfrac{|R_{\text{exp}} - R|}{R}$ entre ta valeur et la valeur de référence $R = \SI{8,314}{J.K^{-1}.mol^{-1}}$. Exprime le résultat en pourcentage et conclus sur la qualité de la mesure.
\item Un élève affirme : « la loi des gaz parfaits est toujours vraie ». À l'aide de la notion de domaine de validité et d'un exemple concret (gaz fortement comprimé ou proche de la liquéfaction), rédige un paragraphe argumenté d'une dizaine de lignes pour discuter cette affirmation. Tu préciseras les hypothèses du modèle et ce qui se passe lorsqu'elles cessent d'être respectées.
\end{enumerate}
\end{exercice}

\newpage

\coursec{Corrigés des exercices}

\begin{corrige}[Exercice 1 — QCM : modèles et lois]
\begin{enumerate}
\item \textbf{b.} Un modèle en physique est une représentation simplifiée de la réalité, fondée sur des hypothèses, qui permet de faire des prévisions. Il n'est jamais une réplique exacte (réponse a fausse) ni une simple maquette (réponse c fausse).
\item \textbf{b.} La loi d'Ohm $U = R\,I$ ne s'applique qu'aux dipôles dits \emph{ohmiques} (conducteurs ohmiques), dont la résistance $R$ est constante, indépendante de la tension et du courant. Elle ne s'applique pas aux dipôles non linéaires (diodes, lampes, transistors) ni aux générateurs.
\item \textbf{c.} Dans l'équation d'état des gaz parfaits $PV = nRT$, la température thermodynamique $T$ doit impérativement être exprimée en \textbf{kelvins (K)}. L'utilisation des degrés Celsius ou Fahrenheit fausserait le calcul car le zéro absolu ($0\,\mathrm{K}$) est l'origine de l'échelle thermodynamique.
\item \textbf{a.} Toute loi physique possède un \cle{domaine de validité} (ensemble des conditions — température, pression, champ, fréquence, etc. — où ses prévisions concordent avec l'expérience). Elle ne possède pas de « valeur infinie » (réponse b absurde) et repose toujours sur des hypothèses simplificatrices (réponse c fausse).
\end{enumerate}
\end{corrige}

\begin{corrige}[Exercice 2 — Vrai ou faux justifié]
\begin{enumerate}
\item \textbf{Faux.} La loi des gaz parfaits ($PV = nRT$) est une loi limite valable pour les gaz réels à \textbf{faible pression} et \textbf{haute température} (loin de la condensation). À haute pression, le volume propre des molécules et les forces intermoléculaires ne sont plus négligeables : la loi échoue (voir équation de Van der Waals).
\item \textbf{Vrai.} C'est le cœur de la démarche scientifique (méthode hypothético-déductive) : un modèle émet des prévisions ; si l'expérience (la mesure) les contredit, le modèle est invalide pour cette situation et doit être \textbf{modifié} (affinement des hypothèses) ou \textbf{abandonné} au profit d'un modèle plus pertinent.
\item \textbf{Vrai.} Pour un conducteur ohmique, $U = R\,I$ avec $R$ constant. La caractéristique $U = f(I)$ est donc une \textbf{droite passant par l'origine} du repère $(I, U)$, de pente $R$.
\item \textbf{Faux.} Les \cle{hypothèses simplificatrices} définissent précisément le domaine de validité du modèle. Si elles ne sont pas respectées (ex. : frottements non négligeables, volume propre des molécules non nul), les prévisions du modèle s'écartent de la réalité. Leur choix conditionne donc directement la qualité et la portée des prévisions.
\end{enumerate}
\end{corrige}

\begin{corrige}[Exercice 3 — Définitions et vocabulaire]
\begin{enumerate}
\item \textbf{Modèle :} Représentation simplifiée et théorique d'un système réel, construite à partir d'hypothèses, qui permet de décrire, d'expliquer et de prévoir le comportement de ce système dans un domaine de validité donné.
\item \textbf{Hypothèse simplificatrice :} Supposition volontaire, faite lors de la construction du modèle, qui néglige certains aspects de la réalité jugés secondaires (ex. : fil sans masse, frottements nuls, molécules ponctuelles) pour rendre le problème mathématiquement traitable.
\item \textbf{Domaine de validité :} Ensemble des conditions expérimentales (valeurs des grandeurs physiques : température, pression, intensité, fréquence, etc.) pour lesquelles les prévisions du modèle concordent avec les mesures à l'intérieur des incertitudes acceptables.
\item Deux hypothèses du modèle du gaz parfait concernant les molécules :
      \begin{itemize}
      \item Les molécules sont des \textbf{particules ponctuelles} : leur volume propre est négligeable devant le volume du contenant ($V_{\text{mol}} \ll V$).
      \item Il n'existe \textbf{aucune force d'interaction} entre les molécules (ni attraction, ni répulsion) en dehors des chocs instantanés et parfaitement élastiques.
      \end{itemize}
\item \textbf{Loi des gaz parfaits :} $PV = nRT$
      \begin{center}
      \begin{tabular}{|l|l|l|}\hline
      \textbf{Grandeur} & \textbf{Symbole} & \textbf{Unité SI} \\ \hline
      Pression & $P$ & pascal (\si{\pascal}) \\
      Volume & $V$ & mètre cube (\si{\cubic\meter}) \\
      Quantité de matière & $n$ & mole (\si{\mole}) \\
      Température thermodynamique & $T$ & kelvin (\si{\kelvin}) \\
      Constante des gaz parfaits & $R$ & joule par kelvin par mole (\si{\joule\per\kelvin\per\mole}) \\ \hline
      \end{tabular}
      \end{center}
      \textit{Analyse dimensionnelle :} $[P][V] = \mathrm{Pa}\cdot\mathrm{m}^3 = \mathrm{N}\cdot\mathrm{m} = \mathrm{J}$. $[n][R][T] = \mathrm{mol}\cdot\mathrm{J}\cdot\mathrm{K}^{-1}\cdot\mathrm{mol}^{-1}\cdot\mathrm{K} = \mathrm{J}$. L'équation est homogène.
\end{enumerate}
\end{corrige}

\begin{corrige}[Exercice 4 — Calculs directs]
\begin{enumerate}
\item \textbf{Loi d'Ohm :} $U = R\,I$
      \[
      U = \SI{47}{\ohm} \times \SI{0,20}{A} = \SI{9,4}{V}.
      \]
      La tension aux bornes du conducteur est de \textbf{\SI{9,4}{V}}.
\item \textbf{Loi des gaz parfaits :} $P = \dfrac{nRT}{V}$
      \begin{itemize}
      \item Conversion du volume : $V = \SI{2,0}{L} = \SI{2,0e-3}{\cubic\meter}$.
      \item $T = \SI{300}{K}$ (déjà en SI).
      \item $n = \num{0,10}\,\mathrm{mol}$.
      \end{itemize}
      \[
      P = \frac{\num{0,10} \times \SI{8,314}{J.K^{-1}.mol^{-1}} \times \SI{300}{K}}{\SI{2,0e-3}{m^3}}
      = \frac{\SI{249,42}{J}}{\SI{2,0e-3}{m^3}}
      = \SI{124710}{Pa} \approx \SI{1,2e5}{Pa}.
      \]
      Conversion en bars : $\SI{1}{bar} = \SI{1e5}{Pa}$.
      \[
      P = \frac{\SI{1,2471e5}{Pa}}{\SI{1e5}{Pa/bar}} = \SI{1,25}{bar} \approx \textbf{\SI{1,2}{bar}}.
      \]
      \textbf{Résultat :} $P \approx \SI{1,2e5}{Pa}$ soit \textbf{\SI{1,2}{bar}} (arrondi à 2 chiffres significatifs, cohérent avec les données $n=0,10$ et $V=2,0$).
\item \textbf{Conversion °C $\to$ K :} $T(\mathrm{K}) = t(^\circ\mathrm{C}) + 273,15$
      \begin{align*}
      t_1 = \SI{27}{\degreeCelsius} &\Rightarrow T_1 = 27 + 273,15 = \SI{300,15}{K} \approx \SI{300}{K}. \\
      t_2 = \SI{-23}{\degreeCelsius} &\Rightarrow T_2 = -23 + 273,15 = \SI{250,15}{K} \approx \SI{250}{K}.
      \end{align*}
\end{enumerate}
\end{corrige}

\begin{corrige}[Exercice 5 — Caractéristique d'un dipôle inconnu]
\begin{enumerate}
\item \textbf{Trace du graphe} (à réaliser sur papier millimétré) :
      \begin{itemize}
      \item Abscisses $I$ (mA) : 1 cm pour 10 mA. Origine en bas à gauche.
      \item Ordonnées $U$ (V) : 1 cm pour 0,2 V.
      \item Points à reporter : $(0;0)$, $(20;0,44)$, $(40;0,90)$, $(60;1,32)$, $(80;1,78)$, $(100;2,21)$.
      \item Les points s'alignent quasi parfaitement sur une droite passant par l'origine.
      \end{itemize}
      \begin{popfigure}
      \begin{tikzpicture}[scale=0.8]
      \begin{axis}[
      width=\linewidth, height=8cm,
      xlabel={$I$ (mA)}, ylabel={$U$ (V)},
      xmin=0, xmax=110, ymin=0, ymax=2.5,
      grid=major, grid style={popline},
      tick label style={font=\small},
      label style={font=\small},
      ]
      \addplot[popcurve, only marks, mark=*, mark size=2pt, popblue] coordinates {
      (0,0) (20,0.44) (40,0.90) (60,1.32) (80,1.78) (100,2.21)
      };
      \addplot[poporange, thick, domain=0:110] {0.0221*x};
      \end{axis}
      \end{tikzpicture}
      \legende{Caractéristique $U=f(I)$ du dipôle inconnu. Les points expérimentaux sont alignés sur une droite passant par l'origine.}
      \end{popfigure}
\item \textbf{Preuve du caractère ohmique} :
      \begin{itemize}
      \item La courbe $U=f(I)$ est une \textbf{droite}.
      \item Cette droite \textbf{passe par l'origine} $(0;0)$.
      \item Le rapport $\dfrac{U}{I}$ est constant (voir calculs ci-dessous).
      \end{itemize}
      Ces deux critères (linéarité + passage par l'origine) prouvent que le dipôle est un \textbf{conducteur ohmique} de résistance constante.
\item \textbf{Détermination graphique de $R$} :
      La résistance $R$ est la pente de la droite $U=f(I)$.
      On choisit deux points éloignés sur la droite de tendance (ex. origine et point à $I=100\,\mathrm{mA}$) :
      \[
      R = \frac{\Delta U}{\Delta I} = \frac{\SI{2,21}{V} - \SI{0}{V}}{\SI{100}{mA} - \SI{0}{mA}} = \frac{\SI{2,21}{V}}{\SI{0,100}{A}} = \SI{22,1}{\ohm}.
      \]
      (En prenant le point $I=60\,\mathrm{mA}$ : $R = 1,32 / 0,060 = \SI{22,0}{\ohm}$).
      \textbf{Valeur retenue :} $R \approx \SI{22,1}{\ohm}$.
\item \textbf{Prévision pour $U = \SI{12}{V}$} :
      \[
      I = \frac{U}{R} = \frac{\SI{12}{V}}{\SI{22,1}{\ohm}} \approx \SI{0,543}{A} = \SI{543}{mA}.
      \]
      \textbf{Commentaire :} Cette valeur (\SI{543}{mA}) est obtenue par \textbf{extrapolation} bien au-delà du domaine des mesures (max \SI{100}{mA}). Le dipôle chaufferait fortement à ce courant, sa résistance augmenterait (effet Joule), et la loi d'Ohm ne serait plus vérifiée. La prévision n'est donc \textbf{pas fiable} hors domaine de validité (courant faible, température constante).
\end{enumerate}
\end{corrige}

\begin{corrige}[Exercice 6 — Bouteille de gaz domestique]
\begin{enumerate}
\item \textbf{Données en SI :}
      $V = \SI{12,5}{L} = \SI{12,5e-3}{\cubic\meter}$,
      $n = \num{5,0}\,\mathrm{mol}$,
      $T_1 = \SI{27}{\degreeCelsius} = \SI{300,15}{K} \approx \SI{300}{K}$,
      $R = \SI{8,314}{J.K^{-1}.mol^{-1}}$.
      \textbf{Pression $P_1$ :}
      \[
      P_1 = \frac{nRT_1}{V} = \frac{5,0 \times 8,314 \times 300,15}{12,5 \times 10^{-3}}
      = \frac{12478,7}{0,0125} = \SI{998296}{Pa} \approx \SI{9,98e5}{Pa}.
      \]
      En bars : $P_1 = \dfrac{\SI{9,98e5}{Pa}}{\SI{1e5}{Pa/bar}} = \SI{9,98}{bar} \approx \textbf{\SI{10,0}{bar}}$.
\item \textbf{Nouvelle pression $P_2$ à $t_2 = \SI{47}{\degreeCelsius}$ ($T_2 = \SI{320,15}{K} \approx \SI{320}{K}$), volume constant :}
      Loi de Gay-Lussac (cas particulier gaz parfait $V$, $n$ constants) : $\dfrac{P_1}{T_1} = \dfrac{P_2}{T_2}$.
      \[
      P_2 = P_1 \frac{T_2}{T_1} = \SI{9,98e5}{Pa} \times \frac{320,15}{300,15}
      \approx \SI{9,98e5}{Pa} \times 1,0666 = \SI{1,064e6}{Pa} \approx \textbf{\SI{10,6}{bar}}.
      \]
\item \textbf{Danger :} L'exposition au soleil fait monter la pression de \SI{10,0}{bar} à \SI{10,6}{bar} (+6\%). Si la bouteille est conçue pour une pression maximale proche de cette valeur (souvent \SI{15}{bar} à \SI{20}{bar} pour le butane, mais la soupape de sécurité peut s'ouvrir vers \SI{15}{bar}- \SI{25}{bar}), le risque est réel : \textbf{ouverture de la soupape (fuite de gaz inflammable) ou rupture de la bouteille (explosion)}.
      \textbf{Deux hypothèses du modèle utilisées :}
      \begin{enumerate}
      \item Les molécules de butane sont ponctuelles (volume propre négligeable).
      \item Il n'y a pas de forces d'attraction entre les molécules de butane.
      \end{enumerate}
      (Or le butane se liquéfie facilement sous pression : ces hypothèses sont limites).
\end{enumerate}
\end{corrige}

\begin{corrige}[Exercice 7 — Vérification de la loi de Boyle-Mariotte]
\begin{enumerate}
\item \textbf{Tableau complété :}
      \begin{center}
      \begin{tabular}{|c|c|c|c|c|c|}\hline
      $V$ (mL) & 60 & 50 & 40 & 30 & 20 \\ \hline
      $P$ (hPa) & 500 & 600 & 750 & 1000 & 1500 \\ \hline
      $1/V$ (mL$^{-1}$) & 0,0167 & 0,0200 & 0,0250 & 0,0333 & 0,0500 \\ \hline
      $P \times V$ (hPa$\cdot$mL) & 30000 & 30000 & 30000 & 30000 & 30000 \\ \hline
      \end{tabular}
      \end{center}
      Le produit $P \times V$ est \textbf{constant} (30000 hPa$\cdot$mL), ce qui confirme la relation.
\item \textbf{Courbe $P = f(1/V)$} :
      \begin{popfigure}
      \begin{tikzpicture}[scale=0.8]
      \begin{axis}[
      width=\linewidth, height=7cm,
      xlabel={$1/V$ (mL$^{-1}$)}, ylabel={$P$ (hPa)},
      xmin=0, xmax=0.055, ymin=0, ymax=1600,
      grid=major, grid style={popline},
      tick label style={font=\small},
      label style={font=\small},
      ]
      \addplot[popcurve, only marks, mark=*, mark size=2pt, popblue] coordinates {
      (0.01667,500) (0.02,600) (0.025,750) (0.03333,1000) (0.05,1500)
      };
      \addplot[poporange, thick, domain=0:0.055] {30000*x};
      \end{axis}
      \end{tikzpicture}
      \legende{Représentation graphique de $P = f(1/V)$. Les points sont alignés sur une droite passant par l'origine.}
      \end{popfigure}
      La courbe est une \textbf{droite passant par l'origine}. La relation mathématique simple est : $P = k \times \dfrac{1}{V}$ ou \textbf{$P V = \text{constante}$}.
\item \textbf{Loi de Boyle-Mariotte :} À \textbf{température constante} et pour une \textbf{quantité de matière constante} ($n$ constant), la pression $P$ d'un gaz parfait est \textbf{inversement proportionnelle} à son volume $V$ : $P V = \text{constante}$ (ou $P_1 V_1 = P_2 V_2$).
      \textbf{Conditions d'application :} Gaz parfait (ou gaz réel à basse pression / haute température), température constante (isotherme), système fermé ($n$ constant).
\item \textbf{Volume pour $P = \SI{2000}{hPa}$} :
      $P V = \text{cste} = 30000\,\mathrm{hPa\cdot mL}$.
      \[
      V = \frac{30000}{2000} = \SI{15}{mL}.
      \]
\item \textbf{Validité à très haute pression :} \textbf{Non.} À très haute pression, les molécules sont très rapprochées :
      \begin{itemize}
      \item Leur \textbf{volume propre} n'est plus négligeable devant le volume du gaz (le volume exclu $b$ dans Van der Waals devient significatif) $\Rightarrow$ $PV > \text{cste}$.
      \item Les \textbf{forces d'attraction intermoléculaires} deviennent significatives $\Rightarrow$ la pression réelle est plus faible que la pression idéale ($PV < \text{cste}$ à pression modérée).
      \end{itemize}
      La loi de Boyle-Mariotte (modèle gaz parfait) atteint sa limite de validité.
\end{enumerate}
\end{corrige}

\begin{corrige}[Exercice 8 — Lampe à filament : la loi d'Ohm en défaut]
\begin{enumerate}
\item \textbf{Calcul du rapport $U/I$ (résistance dynamique) :}
      \begin{center}
      \begin{tabular}{|c|c|c|c|c|c|}\hline
      $I$ (A) & 0,10 & 0,20 & 0,30 & 0,40 & 0,50 \\ \hline
      $U$ (V) & 0,50 & 1,60 & 3,30 & 5,60 & 8,50 \\ \hline
      $U/I$ ($\Omega$) & 5,0 & 8,0 & 11,0 & 14,0 & 17,0 \\ \hline
      \end{tabular}
      \end{center}
      Le rapport $U/I$ \textbf{n'est pas constant} : il augmente avec le courant (de 5 à 17 $\Omega$).
\item \textbf{La loi d'Ohm s'applique-t-elle ?} \textbf{Non.} La caractéristique $U=f(I)$ n'est pas une droite (le rapport $U/I$ varie). Le dipôle n'est \textbf{pas ohmmique}.
\item \textbf{Explication physique :} Le filament de la lampe est en tungstène (métal). Lorsque le courant $I$ le traverse, l'effet Joule ($P = R I^2$) l'échauffe fortement (jusqu'à 2500 K). Or, pour les métaux, \textbf{la résistance augmente avec la température}. L'hypothèse simplificatrice du modèle du conducteur ohmique qui n'est pas respectée est : \textbf{« La température du conducteur reste constante »} (ou « la résistivité du matériau est constante »). Ici, $T$ varie, donc $R$ varie.
\end{enumerate}
\end{corrige}

\begin{corrige}[Exercice 9 — Étude expérimentale d'un dipôle ohmique (Probatoire)]
\textbf{Barème indicatif :} Q1 (1 pt), Q2 (2 pts), Q3 (2 pts), Q4 (2 pts), Q5 (1 pt), Q6 (2 pts). Total 10 pts.
\begin{enumerate}
\item \textbf{Schéma du circuit :}
      \begin{popfigure}
      \begin{center}
      \begin{circuitikz}[european, scale=1.2]
      \draw (0,0) to[battery1, l_={$G$} (générateur réglable), invert] (0,3)
      to[ammeter, l={$A$}] (2,3)
      to[R, l={$D$} (dipôle), *-*] (4,3)
      to[short] (4,0)
      to[short] (0,0);
      \draw (2,3) to[voltmeter, l={$V$}] (2,0);
      \draw (4,3) to[short] (4,0);
      \end{circuitikz}
      \end{center}
      \legende{Montage : Ampèremètre $A$ en série avec $D$, Voltmètre $V$ en dérivation (parallèle) aux bornes de $D$.}
      \end{popfigure}
\item \textbf{Caractéristique $U = f(I)$} :
      \begin{popfigure}
      \begin{tikzpicture}[scale=0.8]
      \begin{axis}[
      width=\linewidth, height=7cm,
      xlabel={$I$ (mA)}, ylabel={$U$ (V)},
      xmin=0, xmax=100, ymin=0, ymax=10,
      grid=major, grid style={popline},
      tick label style={font=\small},
      label style={font=\small},
      ]
      \addplot[popcurve, only marks, mark=*, mark size=2pt, popblue] coordinates {
      (0,0) (15,1.5) (30,3.0) (45,4.5) (60,6.0) (75,7.5) (90,9.0)
      };
      \addplot[poporange, thick, domain=0:100] {0.1*x};
      \end{axis}
      \end{tikzpicture}
      \legende{Caractéristique $U=f(I)$ du dipôle $D$. Droite passant par l'origine.}
      \end{popfigure}
      Échelles respectées : 1 cm $\equiv$ 10 mA ; 1 cm $\equiv$ 1 V. Points alignés sur une droite passant par l'origine.
\item \textbf{Preuve du caractère ohmique et loi :}
      \begin{itemize}
      \item Les points expérimentaux sont alignés sur une \textbf{droite}.
      \item Cette droite \textbf{passe par l'origine} $(0;0)$.
      \item Le rapport $U/I$ est constant : $\frac{1,5}{15} = \frac{3,0}{30} = \dots = \frac{9,0}{90} = \SI{0,1}{V/mA} = \SI{100}{V/A}$.
      \end{itemize}
      Le dipôle $D$ est un \textbf{conducteur ohmique}. La loi mathématique est : $\boxed{U = R\,I}$ avec $R = \SI{100}{\ohm}$.
\item \textbf{Détermination graphique de $R$ :}
      \textbf{Méthode :} $R$ est la pente de la droite $U=f(I)$.
      On choisit deux points sur la droite de tendance (ex. $A(0;0)$ et $B(90\,\mathrm{mA}; 9,0\,\mathrm{V})$).
      \[
      R = \frac{U_B - U_A}{I_B - I_A} = \frac{\SI{9,0}{V} - \SI{0}{V}}{\SI{90}{mA} - \SI{0}{mA}} = \frac{9,0}{0,090} = \SI{100}{\ohm}.
      \]
      \textbf{Résistance :} $\boxed{R = \SI{100}{\ohm}}$.
\item \textbf{Prévision pour $U = \SI{12}{V}$} :
      \[
      I = \frac{U}{R} = \frac{\SI{12}{V}}{\SI{100}{\ohm}} = \SI{0,12}{A} = \boxed{\SI{120}{mA}}.
      \]
\item \textbf{Limite du modèle (chauffage) :}
      En réalité, le dipôle chauffe (effet Joule $P = R I^2$). Or, pour un conducteur métallique, \textbf{la résistance augmente avec la température}. Le modèle « conducteur ohmique à $R$ constant » suppose \textbf{température constante} (faible courant). À \SI{120}{mA}, la puissance dissipée est $P = U I = 12 \times 0,12 = \SI{1,44}{W}$, non négligeable pour un petit dipôle. La résistance réelle $R_{\text{réel}} > \SI{100}{\ohm}$, donc l'intensité réelle mesurée sera \textbf{inférieure à \SI{120}{mA}}.
      \textbf{Domaine de validité du modèle :} Courants faibles (échauffement négligeable), température ambiante constante.
\end{enumerate}
\textbf{Points de vigilance Probatoire :}
\begin{itemize}
\item Schéma : Ampèremètre \emph{en série} (symbole A dans un cercle), Voltmètre \emph{en parallèle} (symbole V dans un cercle). Polarités respectées (+ du générateur vers + des appareils).
\item Graphique : Échelles linéaires, régulières, unités en abscisse/ordonnée. Points repérés par des croix ou ronds. Droite de tendance tracée à la règle (ne pas relier les points).
\item Preuve ohmique : Citer explicitement « droite » ET « passage par l'origine » (ou proportionnalité $U/I = \text{cste}$).
\item Q6 : Lier chauffage $\to$ augmentation $R$ $\to$ écart prévision/mesure. Citer l'hypothèse « température constante » non respectée.
\end{itemize}
\end{corrige}

\begin{corrige}[Exercice 10 — Gaz parfait et sécurité domestique (Probatoire)]
\textbf{Barème indicatif :} Q1 (1,5 pt), Q2 (1 pt), Q3 (2 pts), Q4 (2 pts), Q5 (1,5 pt), Q6 (2 pts). Total 10 pts.
\begin{enumerate}
\item \textbf{Loi des gaz parfaits :} $P V = n R T$
      \begin{itemize}
      \item $P$ : Pression du gaz, en \textbf{pascal (\si{\pascal})}.
      \item $V$ : Volume occupé par le gaz, en \textbf{mètre cube (\si{\cubic\meter})}.
      \item $n$ : Quantité de matière de gaz, en \textbf{mole (\si{\mole})}.
      \item $T$ : Température thermodynamique, en \textbf{kelvin (\si{\kelvin})}.
      \item $R$ : Constante des gaz parfaits, $R \approx \SI{8,314}{J.K^{-1}.mol^{-1}}$.
      \end{itemize}
\item \textbf{Deux hypothèses simplificatrices :}
      \begin{enumerate}
      \item Les molécules sont des \textbf{particules ponctuelles} (volume propre négligeable).
      \item Il n'existe \textbf{aucune interaction} (force d'attraction ou de répulsion) entre les molécules hors chocs.
      \end{enumerate}
      (Autres hypothèses acceptées : chocs élastiques, durée des chocs négligeable, nombre de molécules très grand).
\item \textbf{Calcul de $P_1$ :}
      $V = \SI{12,5}{L} = \SI{12,5e-3}{m^3}$, $n = \num{5,0}\,\mathrm{mol}$, $T_1 = \SI{27}{\degreeCelsius} = \SI{300,15}{K}$.
      \[
      P_1 = \frac{n R T_1}{V} = \frac{5,0 \times 8,314 \times 300,15}{12,5 \times 10^{-3}} = \SI{9,98e5}{Pa}.
      \]
      En bars : $P_1 = \dfrac{\SI{9,98e5}{Pa}}{\SI{1e5}{Pa/bar}} = \boxed{\SI{9,98}{bar}}$.
\item \textbf{Calcul de $P_2$ à $t_2 = \SI{47}{\degreeCelsius}$ ($T_2 = \SI{320,15}{K}$) :}
      Volume $V$ et quantité $n$ constants $\Rightarrow$ Loi de Gay-Lussac : $\dfrac{P_1}{T_1} = \dfrac{P_2}{T_2}$.
      \[
      P_2 = P_1 \frac{T_2}{T_1} = \SI{9,98e5}{Pa} \times \frac{320,15}{300,15} = \SI{1,064e6}{Pa} \approx \boxed{\SI{1,06e6}{Pa}} \quad (\text{soit } \SI{10,6}{bar}).
      \]
\item \textbf{Danger ?} Pression maximale supportée $P_{\max} = \SI{1,5e6}{Pa} = \SI{15}{bar}$.
      $P_2 = \SI{1,06e6}{Pa} < P_{\max}$.
      \textbf{Conclusion :} Dans cette situation précise ($t_2 = \SI{47}{\degreeCelsius}$), la pression calculée (\SI{10,6}{bar}) reste \textbf{inférieure} à la pression maximale de rupture (\SI{15}{bar}). Il n'y a \textbf{pas de danger immédiat de rupture} de la bouteille.
      \textit{Nuance :} La soupape de sécurité est souvent tarée plus bas (ex. \SI{15}{bar} à \SI{25}{bar}), mais la marge de sécurité se réduit. Une température plus élevée (incendie) ferait dépasser $P_{\max}$.
\item \textbf{Limites du modèle à forte pression (liquéfaction) :}
      À forte pression (et/ou basse température), le butane réel s'éloigne du modèle parfait :
      \begin{itemize}
      \item \textbf{Volume propre non négligeable :} Les molécules de butane (grosses, polaires) occupent une fraction significative du volume $V$. L'espace libre pour le mouvement est $V - nb$ ($b$ = volume propre molaire). La pression réelle devient \textbf{supérieure} à la pression idéale ($P_{\text{réel}} > P_{\text{idéal}}$).
      \item \textbf{Forces d'attraction intermoléculaires :} Les molécules de butane s'attirent (forces de Van der Waals). Cela tend à les rapprocher, diminuant la pression exercée sur les parois : $P_{\text{réel}} < P_{\text{idéal}}$.
      \item \textbf{Changement de phase :} Si $P$ et $T$ franchissent la courbe de coexistence liquide-vapeur, une fraction du gaz se \textbf{liquéfie}. La quantité de matière en phase gazeuse $n_{\text{gaz}}$ diminue, la loi $PV=nRT$ (avec $n$ constant) ne décrit plus l'état du système. C'est la fin du domaine de validité du modèle « gaz parfait » pour ce système.
      \end{itemize}
\end{enumerate}
\textbf{Points de vigilance Probatoire :}
\begin{itemize}
\item Q1 : Unités SI obligatoires (pas de bar, pas de L, pas de °C dans la formule).
\item Q4 : Montrer l'égalité des rapports $P/T$ (Gay-Lussac) avant de calculer.
\item Q5 : Comparer explicitement $P_2$ et $P_{\max}$ avec les mêmes unités. Conclure par une phrase.
\item Q6 : Argumentation structurée : citer les 2 hypothèses principales $\to$ expliquer leur violation $\to$ conséquence sur $P$ ou $n$ $\to$ conclusion sur la limite du modèle.
\end{itemize}
\end{corrige}

\begin{corrige}[Exercice 11 — Détermination expérimentale de la constante $R$ et discussion (Probatoire)]
\textbf{Barème indicatif :} Q1 (1 pt), Q2 (2 pts), Q3 (2 pts), Q4 (5 pts). Total 10 pts.
\begin{enumerate}
\item \textbf{Conversion en unités SI :}
      \begin{align*}
      V &= \SI{1,24}{L} = \SI{1,24e-3}{\cubic\meter}. \\
      T &= t + 273,15 = 25 + 273,15 = \SI{298,15}{K}.
      \end{align*}
\item \textbf{Valeur expérimentale $R_{\text{exp}}$ :}
      Loi des gaz parfaits : $R = \dfrac{P V}{n T}$.
      \[
      R_{\text{exp}} = \frac{\SI{1,02e5}{Pa} \times \SI{1,24e-3}{m^3}}{\num{0,050}\,mol \times \SI{298,15}{K}}
      = \frac{\SI{126,48}{J}}{\SI{14,9075}{mol.K}}
      = \SI{8,484}{J.K^{-1}.mol^{-1}}.
      \]
      \textbf{Résultat :} $\boxed{R_{\text{exp}} \approx \SI{8,48}{J.K^{-1}.mol^{-1}}}$.
\item \textbf{Écart relatif et conclusion :}
      Valeur de référence $R_{\text{ref}} = \SI{8,314}{J.K^{-1}.mol^{-1}}$.
      \[
      \text{Écart relatif} = \frac{|R_{\text{exp}} - R_{\text{ref}}|}{R_{\text{ref}}}
      = \frac{|8,484 - 8,314|}{8,314}
      = \frac{0,170}{8,314} \approx 0,0204 = \boxed{2,04\,\%}.
      \]
      \textbf{Conclusion :} L'écart relatif est d'environ \textbf{2\%}. C'est un résultat \textbf{très satisfaisant} pour une mesure de TP lycéen (incertitudes sur $P, V, T, n$). La mesure est compatible avec la valeur de référence ; l'air se comporte bien comme un gaz parfait dans ces conditions (pression atmosphérique, température ambiante).
\item \textbf{Discussion argumentée : « La loi des gaz parfaits est toujours vraie » — FAUX.}
      \begin{quote}
      \textit{L'affirmation est fausse. La loi des gaz parfaits ($PV=nRT$) est un \textbf{modèle} fondé sur des \textbf{hypothèses simplificatrices} : les molécules sont ponctuelles (volume propre nul) et n'interagissent pas entre elles (pas de forces intermoléculaires). Ces hypothèses définissent son \textbf{domaine de validité} : pressions faibles, températures élevées (loin de la condensation).}
      
      \textit{Lorsque la pression augmente fortement (ou la température baisse), ces hypothèses cessent d'être vérifiées. Le volume propre des molécules ($b$ dans l'équation de Van der Waals) devient non négligeable devant le volume du gaz, ce qui tend à augmenter la pression par rapport à l'idéal. Surtout, les forces d'attraction intermoléculaires ($a$ dans Van der Waals) deviennent significatives : elles rapprochent les molécules et diminuent la pression exercée sur les parois.}
      
      \textit{Exemple concret : le butane dans une bouteille domestique. À \SI{20}{\degreeCelsius}, sa pression de vapeur saturante est d'environ \SI{2,1e5}{Pa} (\SI{2,1}{bar}). Si on comprime le butane au-delà de cette pression à cette température, il se \textbf{liquéfie} partiellement. La quantité de matière en phase gazeuse $n$ diminue, et la relation $PV=nRT$ (avec $n$ constant) ne décrit plus l'état du système. Le modèle « gaz parfait » a atteint sa limite ; il faut utiliser un modèle d'état réel (Van der Waals, Redlich-Kwong) ou des tables de propriétés thermodynamiques.}
      \end{quote}
\end{enumerate}
\textbf{Points de vigilance Probatoire :}
\begin{itemize}
\item Q2 : Utiliser strictement les unités SI ($m^3, K, Pa, mol$) dans la formule.
\item Q3 : Calcul de l'écart relatif en valeur absolue, conversion en \%. Conclusion quantitative (« compatible », « bon accord »).
\item Q4 : Structure attendue : 1) Rappel hypotheses/modèle. 2) Définition domaine validité. 3) Explication physique de la rupture (volume propre + forces). 4) Exemple concret (butane/liquéfaction). 5) Conclusion : modèle approximatif, pas « vérité absolue ». Longueur ~10 lignes.
\end{itemize}
\end{corrige}

\newpage

\coursec{Fiche bilan}

\begin{popfigure}
\centering
\begin{center}\tcbox[colback=popdark,colframe=popdark,coltext=white,arc=9pt,boxsep=4pt,left=6pt,right=6pt]{\bfseries\small Modéliser un phénomène}\end{center}
\vspace{-2pt}
\begin{tcbraster}[raster columns=3,raster equal height=rows,raster column skip=6pt,raster row skip=6pt,enhanced,blankest]
\begin{tcolorbox}[enhanced,colback=white,colframe=popblue,boxrule=0.9pt,arc=3pt,title={Notion de modèle},fonttitle=\bfseries\scriptsize,coltitle=white,colbacktitle=popblue,left=4pt,right=4pt,top=2pt,bottom=2pt,boxsep=1pt,toptitle=1.5pt,bottomtitle=1.5pt,halign=left]
\begin{itemize}[nosep,leftmargin=*,itemsep=1pt,font=\scriptsize]
\item Représentation simplifiée
\item Hypothèses
\end{itemize}
\end{tcolorbox}
\begin{tcolorbox}[enhanced,colback=white,colframe=poporange,boxrule=0.9pt,arc=3pt,title={Démarche expérimentale},fonttitle=\bfseries\scriptsize,coltitle=white,colbacktitle=poporange,left=4pt,right=4pt,top=2pt,bottom=2pt,boxsep=1pt,toptitle=1.5pt,bottomtitle=1.5pt,halign=left]
\begin{itemize}[nosep,leftmargin=*,itemsep=1pt,font=\scriptsize]
\item Mesures
\item Graphe
\item Équation
\end{itemize}
\end{tcolorbox}
\begin{tcolorbox}[enhanced,colback=white,colframe=popgreen,boxrule=0.9pt,arc=3pt,title={Loi d'Ohm},fonttitle=\bfseries\scriptsize,coltitle=white,colbacktitle=popgreen,left=4pt,right=4pt,top=2pt,bottom=2pt,boxsep=1pt,toptitle=1.5pt,bottomtitle=1.5pt,halign=left]
\begin{itemize}[nosep,leftmargin=*,itemsep=1pt,font=\scriptsize]
\item $U=R\,I$
\item Dipôle ohmique
\end{itemize}
\end{tcolorbox}
\begin{tcolorbox}[enhanced,colback=white,colframe=poppurple,boxrule=0.9pt,arc=3pt,title={Gaz parfait},fonttitle=\bfseries\scriptsize,coltitle=white,colbacktitle=poppurple,left=4pt,right=4pt,top=2pt,bottom=2pt,boxsep=1pt,toptitle=1.5pt,bottomtitle=1.5pt,halign=left]
\begin{itemize}[nosep,leftmargin=*,itemsep=1pt,font=\scriptsize]
\item $PV=nRT$
\item $R=\SI{8,314}{SI}$
\end{itemize}
\end{tcolorbox}
\begin{tcolorbox}[enhanced,colback=white,colframe=poppink,boxrule=0.9pt,arc=3pt,title={Domaine de validité},fonttitle=\bfseries\scriptsize,coltitle=white,colbacktitle=poppink,left=4pt,right=4pt,top=2pt,bottom=2pt,boxsep=1pt,toptitle=1.5pt,bottomtitle=1.5pt,halign=left]
\begin{itemize}[nosep,leftmargin=*,itemsep=1pt,font=\scriptsize]
\item Limites
\item Conditions
\end{itemize}
\end{tcolorbox}
\begin{tcolorbox}[enhanced,colback=white,colframe=popgold,boxrule=0.9pt,arc=3pt,title={Esprit scientifique},fonttitle=\bfseries\scriptsize,coltitle=white,colbacktitle=popgold,left=4pt,right=4pt,top=2pt,bottom=2pt,boxsep=1pt,toptitle=1.5pt,bottomtitle=1.5pt,halign=left]
\begin{itemize}[nosep,leftmargin=*,itemsep=1pt,font=\scriptsize]
\item Tester par la mesure
\end{itemize}
\end{tcolorbox}
\end{tcbraster}

\legende{Carte mentale de la leçon : les six idées forces autour de la modélisation.}
\end{popfigure}

\begin{bilan}
\textbf{Définitions clés}
\begin{itemize}
  \item \cle{Modèle} : représentation simplifiée d'un phénomène réel, construite à partir d'\textbf{hypothèses simplificatrices}, permettant de décrire, d'expliquer et de \textbf{prévoir} des résultats.
  \item \cle{Loi physique} : relation mathématique entre grandeurs mesurables, établie expérimentalement et valable dans un \textbf{domaine de validité} précis.
  \item \cle{Hypothèse simplificatrice} : approximation volontaire (frottements négligés, gaz parfait, température constante\ldots) qui rend le problème traitable.
  \item \cle{Domaine de validité} : ensemble des conditions (plages de valeurs, nature du système) dans lesquelles la loi donne des prévisions conformes aux mesures.
\end{itemize}

\textbf{Formules à connaître par cœur}
\begin{center}
\begin{tabularx}{\linewidth}{|l|X|X|}
\hline
\rowcolor{popblueL} \textbf{Loi} & \textbf{Expression} & \textbf{Conditions d'application} \\
\hline
Loi d'Ohm & $U = R\,I$ \quad ($U$ en V, $R$ en $\Omega$, $I$ en A) & Conducteur ohmique, à température constante \\
\hline
Gaz parfaits & $PV = nRT$ \quad ($P$ en Pa, $V$ en m$^3$, $T$ en K) & Gaz à pression faible et température modérée \\
\hline
Proportionnalité & $y = k\,x$ : droite passant par l'origine & Testée par le graphe $y = f(x)$ \\
\hline
\end{tabularx}
\end{center}

\textbf{Méthode express : établir une loi à partir de mesures}
\begin{enumerate}
  \item Identifier les grandeurs pertinentes et formuler des \textbf{hypothèses simplificatrices}.
  \item Réaliser des \textbf{mesures} en ne faisant varier qu'un seul paramètre à la fois.
  \item Tracer le \textbf{graphe} : une droite passant par l'origine traduit une proportionnalité.
  \item Déterminer le coefficient (pente) et écrire l'\textbf{équation} de la loi.
  \item \textbf{Tester} la loi par de nouvelles mesures et préciser son \textbf{domaine de validité}.
\end{enumerate}

\textbf{Pièges à éviter}
\begin{itemize}
  \item Confondre $T$ en degrés Celsius et en kelvins : $T(\mathrm{K}) = \theta(^\circ\mathrm{C}) + 273$.
  \item Oublier les conversions SI : $V$ en m$^3$ ($\SI{1}{L} = \SI{e-3}{m^3}$), $P$ en Pa.
  \item Conclure à une loi sans vérifier l'alignement des points ni l'origine du graphe.
  \item Appliquer une loi hors de son domaine de validité (ex. gaz parfait à très haute pression).
\end{itemize}
\end{bilan}

\begin{aretenir}[Checklist avant le Probatoire]
\begin{itemize}
  \item[$\square$] Je sais définir ce qu'est un modèle et citer au moins deux hypothèses simplificatrices.
  \item[$\square$] Je sais décrire les étapes de la démarche expérimentale (hypothèses $\to$ mesures $\to$ graphe $\to$ loi $\to$ test).
  \item[$\square$] Je connais la loi d'Ohm $U = R\,I$ avec ses unités et ses conditions d'application.
  \item[$\square$] Je sais exploiter la caractéristique $U = f(I)$ : pente $= R$ si la droite passe par l'origine.
  \item[$\square$] Je connais la loi des gaz parfaits $PV = nRT$ et la valeur $R = \SI{8,314}{J.K^{-1}.mol^{-1}}$.
  \item[$\square$] Je convertis sans erreur : $T$ en kelvins, $V$ en m$^3$, $P$ en pascals.
  \item[$\square$] Je sais reconnaître graphiquement une proportionnalité (droite passant par l'origine).
  \item[$\square$] Je sais déterminer une pente à partir de deux points : $k = \dfrac{\Delta y}{\Delta x}$.
  \item[$\square$] Je sais citer le domaine de validité de chaque loi étudiée.
  \item[$\square$] Je sais tester une loi en comparant une valeur calculée à une valeur mesurée.
  \item[$\square$] Je rédige mes résultats avec un nombre cohérent de chiffres significatifs et l'unité SI.
  \item[$\square$] Je sais expliquer pourquoi un modèle peut être abandonné ou amélioré.
\end{itemize}
\end{aretenir}

\findecours

\end{document}
